\chapter{Canindeyú}\label{dept:Canindeyú}
\mapaparaguai{14}{Canindeyú}
\vfill\clearpage
\section{Canindeyu /\allowbreak  pacote 16}

O pacote 16 cobre o departamento , com 16 localidades. Canindeyu e
departamento do nordeste paraguaio, fronteirico com o Brasil, marcado
pela floresta atlantica interior, o Rio Parana e vocacao de agricultura
mecanizada e comercio de fronteira.

\begin{itemize}
\tightlist
\item
  Salto del Guaira e a capital, polo de comercio internacional de
  fronteira com Mundo Novo/\allowbreak Guaira (Brasil); 28.553 hab.
\item
  Vocacao de soja, trigo, milho, pecuaria e servicos de fronteira.
\item
  Lacuna de solar/\allowbreak pluviometria da capital fechada com dados NASA POWER
  (2001-2020).
\end{itemize}

\subsection{Ficha climática de Salto del Guaira (capital
departamental)}

Fonte: NASA POWER Climatology API, período 2001-2020, coordenadas
-24.06°S /\allowbreak  -54.33°W (acesso em 2026-03-20).

\textbf{Irradiância solar global --- ALLSKY\_SFC\_SW\_DWN (kWh/\allowbreak m²/\allowbreak dia):}

\begin{longtable}[]{@{}
  >{\raggedright\arraybackslash}p{(\linewidth - 22\tabcolsep) * \real{0.0833}}
  >{\raggedright\arraybackslash}p{(\linewidth - 22\tabcolsep) * \real{0.0833}}
  >{\raggedright\arraybackslash}p{(\linewidth - 22\tabcolsep) * \real{0.0833}}
  >{\raggedright\arraybackslash}p{(\linewidth - 22\tabcolsep) * \real{0.0833}}
  >{\raggedright\arraybackslash}p{(\linewidth - 22\tabcolsep) * \real{0.0833}}
  >{\raggedright\arraybackslash}p{(\linewidth - 22\tabcolsep) * \real{0.0833}}
  >{\raggedright\arraybackslash}p{(\linewidth - 22\tabcolsep) * \real{0.0833}}
  >{\raggedright\arraybackslash}p{(\linewidth - 22\tabcolsep) * \real{0.0833}}
  >{\raggedright\arraybackslash}p{(\linewidth - 22\tabcolsep) * \real{0.0833}}
  >{\raggedright\arraybackslash}p{(\linewidth - 22\tabcolsep) * \real{0.0833}}
  >{\raggedright\arraybackslash}p{(\linewidth - 22\tabcolsep) * \real{0.0833}}
  >{\raggedright\arraybackslash}p{(\linewidth - 22\tabcolsep) * \real{0.0833}}@{}}
\toprule\noalign{}
\begin{minipage}[b]{\linewidth}\raggedright
Jan
\end{minipage} & \begin{minipage}[b]{\linewidth}\raggedright
Fev
\end{minipage} & \begin{minipage}[b]{\linewidth}\raggedright
Mar
\end{minipage} & \begin{minipage}[b]{\linewidth}\raggedright
Abr
\end{minipage} & \begin{minipage}[b]{\linewidth}\raggedright
Mai
\end{minipage} & \begin{minipage}[b]{\linewidth}\raggedright
Jun
\end{minipage} & \begin{minipage}[b]{\linewidth}\raggedright
Jul
\end{minipage} & \begin{minipage}[b]{\linewidth}\raggedright
Ago
\end{minipage} & \begin{minipage}[b]{\linewidth}\raggedright
Set
\end{minipage} & \begin{minipage}[b]{\linewidth}\raggedright
Out
\end{minipage} & \begin{minipage}[b]{\linewidth}\raggedright
Nov
\end{minipage} & \begin{minipage}[b]{\linewidth}\raggedright
Dez
\end{minipage} \\
\midrule\noalign{}
\endhead
\bottomrule\noalign{}
\endlastfoot
6.48 & 5.93 & 5.56 & 4.60 & 3.43 & 3.02 & 3.35 & 4.16 & 4.70 & 5.31 &
6.29 & 6.46 \\
\end{longtable}

Média anual: 4.94 kWh/\allowbreak m²/\allowbreak dia. Inclinação recomendada: 24° Norte (anual);
34° inverno; 14° verão.

\textbf{Precipitação --- PRECTOTCORR (mm/\allowbreak dia):}

\begin{longtable}[]{@{}
  >{\raggedright\arraybackslash}p{(\linewidth - 22\tabcolsep) * \real{0.0833}}
  >{\raggedright\arraybackslash}p{(\linewidth - 22\tabcolsep) * \real{0.0833}}
  >{\raggedright\arraybackslash}p{(\linewidth - 22\tabcolsep) * \real{0.0833}}
  >{\raggedright\arraybackslash}p{(\linewidth - 22\tabcolsep) * \real{0.0833}}
  >{\raggedright\arraybackslash}p{(\linewidth - 22\tabcolsep) * \real{0.0833}}
  >{\raggedright\arraybackslash}p{(\linewidth - 22\tabcolsep) * \real{0.0833}}
  >{\raggedright\arraybackslash}p{(\linewidth - 22\tabcolsep) * \real{0.0833}}
  >{\raggedright\arraybackslash}p{(\linewidth - 22\tabcolsep) * \real{0.0833}}
  >{\raggedright\arraybackslash}p{(\linewidth - 22\tabcolsep) * \real{0.0833}}
  >{\raggedright\arraybackslash}p{(\linewidth - 22\tabcolsep) * \real{0.0833}}
  >{\raggedright\arraybackslash}p{(\linewidth - 22\tabcolsep) * \real{0.0833}}
  >{\raggedright\arraybackslash}p{(\linewidth - 22\tabcolsep) * \real{0.0833}}@{}}
\toprule\noalign{}
\begin{minipage}[b]{\linewidth}\raggedright
Jan
\end{minipage} & \begin{minipage}[b]{\linewidth}\raggedright
Fev
\end{minipage} & \begin{minipage}[b]{\linewidth}\raggedright
Mar
\end{minipage} & \begin{minipage}[b]{\linewidth}\raggedright
Abr
\end{minipage} & \begin{minipage}[b]{\linewidth}\raggedright
Mai
\end{minipage} & \begin{minipage}[b]{\linewidth}\raggedright
Jun
\end{minipage} & \begin{minipage}[b]{\linewidth}\raggedright
Jul
\end{minipage} & \begin{minipage}[b]{\linewidth}\raggedright
Ago
\end{minipage} & \begin{minipage}[b]{\linewidth}\raggedright
Set
\end{minipage} & \begin{minipage}[b]{\linewidth}\raggedright
Out
\end{minipage} & \begin{minipage}[b]{\linewidth}\raggedright
Nov
\end{minipage} & \begin{minipage}[b]{\linewidth}\raggedright
Dez
\end{minipage} \\
\midrule\noalign{}
\endhead
\bottomrule\noalign{}
\endlastfoot
4.42 & 4.71 & 3.28 & 3.68 & 4.50 & 2.74 & 2.08 & 2.26 & 3.10 & 5.70 &
5.63 & 6.07 \\
\end{longtable}

Média anual: \textasciitilde1.464 mm/\allowbreak ano. Estação chuvosa Out-Dez; seca
Jul-Ago.
\clearpage
\subsection*{Dados Departamentais}
\subsubsection{Desenvolvimento Humano e
Social}

\subsection{IDH Departamental}

\begin{longtable}[]{@{}ll@{}}
\toprule\noalign{}
Indicador & Valor \\
\midrule\noalign{}
\endhead
\bottomrule\noalign{}
\endlastfoot
IDH & 0,685 \\
Classificação IDH & médio \\
Ranking nacional (1=melhor) & 15/\allowbreak 18 \\
Esperança de vida & 71,2 anos \\
Anos de escolaridade média & 7,3 anos \\
RNB per capita (USD PPA) & 7.100 \\
\end{longtable}

\subsection{Indicadores de Pobreza}

\begin{longtable}[]{@{}ll@{}}
\toprule\noalign{}
Indicador & Valor \\
\midrule\noalign{}
\endhead
\bottomrule\noalign{}
\endlastfoot
Taxa de pobreza total (\%) & 35,1\% \\
Taxa de pobreza extrema (\%) & 10,8\% \\
Índice de Gini & 0,467 \\
\end{longtable}

\subsection{Infraestrutura Básica}

\begin{longtable}[]{@{}ll@{}}
\toprule\noalign{}
Indicador & Valor \\
\midrule\noalign{}
\endhead
\bottomrule\noalign{}
\endlastfoot
Acesso a água potável (\%) & 67,4\% \\
Acesso a saneamento (\%) & 64,8\% \\
\end{longtable}

\subsubsection{Segurança Pública}

\subsection{Indicadores}

\begin{longtable}[]{@{}
  >{\raggedright\arraybackslash}p{(\linewidth - 2\tabcolsep) * \real{0.6111}}
  >{\raggedright\arraybackslash}p{(\linewidth - 2\tabcolsep) * \real{0.3889}}@{}}
\toprule\noalign{}
\begin{minipage}[b]{\linewidth}\raggedright
Indicador
\end{minipage} & \begin{minipage}[b]{\linewidth}\raggedright
Valor
\end{minipage} \\
\midrule\noalign{}
\endhead
\bottomrule\noalign{}
\endlastfoot
Taxa de homicídios (por 100k hab) & \textasciitilde15--25 (estimativa,
significativamente acima da média nacional) \\
Crimes registrados (total/\allowbreak ano) & --- dados não desagregados
publicamente \\
Número de comisarías & \textasciitilde10 (estimativa departamental) \\
Índice segurança relativo & baixo \\
\end{longtable}

\subsubsection{Mercado de Terra Rural}

\subsection{Preços por Tipo de Terra}

\begin{longtable}[]{@{}llll@{}}
\toprule\noalign{}
Tipo & Preço (USD/\allowbreak ha) & Faixa & Tendência \\
\midrule\noalign{}
\endhead
\bottomrule\noalign{}
\endlastfoot
Agrícola alta produtividade & 7.500 & 5.000 -- 11.000 & ↑ \\
Pastagem & 3.500 & 2.200 -- 5.000 & ↑ \\
Mata /\allowbreak  reserva & 1.500 & 800 -- 2.500 & → \\
\end{longtable}

\subsection{Características do
Mercado}

\begin{longtable}[]{@{}ll@{}}
\toprule\noalign{}
Parâmetro & Valor \\
\midrule\noalign{}
\endhead
\bottomrule\noalign{}
\endlastfoot
Valorização anual média (\%) & 8--10\% \\
Lote mínimo rural (ha) & 10 \\
Imposto territorial rural (\%) & 1\% sobre valor fiscal \\
Liquidez do mercado & média--alta \\
\end{longtable}
\clearpage
\newpage
\secaoDiagnostico{Corpus Christi}{\mapaConteudo{-24.0166}{-54.8333}}{Corpus Christi é um enclave de produtividade e baixo perfil no nordeste
de Canindeyú. Oferece solos de produtividade mundial e um isolamento
demográfico que garante privacidade absoluta, protegida pela distância
dos centros de conflito primários. Sua força reside na abundância de
recursos naturais e na eficiência da produção de grãos. Contudo, a
proximidade com a fronteira seca e o contexto regional exigem que o
ocupante adote um perfil vigilante e autônomo. É o local ideal para
projetos de autossuficiência que valorizem a invisibilidade e a
segurança institucional de fronteira em solo de alta performance.

\subsubsection{Por que sim}

\begin{itemize}
\tightlist
\item
  Solos de altíssima produtividade (``terra roxa'') e clima favorável.
\item
  Baixíssima densidade populacional (privacidade absoluta).
\item
  Abundância de recursos hídricos e proteicos.
\item
  Longe de eixos de fallout tático metropolitano.
\end{itemize}

\subsubsection{Por que não}

\begin{itemize}
\tightlist
\item
  Insegurança tática regional latente (fronteira).
\item
  Restrições da Lei de Fronteira (2532/\allowbreak 05) para investidores
  estrangeiros.
\item
  Distância de centros urbanos com serviços de saúde de alta
  complexidade.
\end{itemize}

\subsubsection{Recomendação
Técnica}

Altamente indicado para estabelecimentos de autossuficiência tecnificada
que valorizem a qualidade do solo. Requer conformidade jurídica e
investimento em segurança patrimonial. O GSS de 7.0 reflete a
superioridade dos recursos produtivos em equilíbrio com os desafios do
contexto regional. Referencia atual de score: GSS 7.0 em 2026-03-05.

\subsubsection{}

\begin{itemize}
\tightlist
\item
  \begin{enumerate}
  \def\labelenumi{\arabic{enumi})}
  \tightlist
  \item
    Dados censitários validados (INE\footnote{Instituto Nacional de Estadística (Instituto Nacional de Estatística), órgão oficial de dados demográficos e censitários.} 2022).
  \end{enumerate}
\item
  \begin{enumerate}
  \def\labelenumi{\arabic{enumi})}
  \setcounter{enumi}{1}
  \tightlist
  \item
    Relatórios de produção agrícola departamental (MAG).
  \end{enumerate}
\item
  \begin{enumerate}
  \def\labelenumi{\arabic{enumi})}
  \setcounter{enumi}{2}
  \tightlist
  \item
    Mapa de solos e bacias hidrográficas de Canindeyú (MAG).
  \end{enumerate}
\item
  \begin{enumerate}
  \def\labelenumi{\arabic{enumi})}
  \setcounter{enumi}{3}
  \tightlist
  \item
    Registro de segurança pública e inteligência de fronteira.
  \end{enumerate}
\end{itemize}}
\begin{table}[ht!]
\small \begin{tabular}{p{3.0cm}p{0.8cm}p{6.0cm}} \toprule
\textbf{Critério} & \textbf{Nota} & \textbf{Justificativa} \\ \midrule
A - Ameaças Estratégicas & 6.5 & Localização remota e discreta; excelente invisibilidade tática fora da fronteira seca, equilibrada por riscos regionais \\
B - Risco Social & 8.0 & Densidade populacional baixa; isolamento social superior em ambiente rural \\
C - Riscos Naturais & 7.0 & Geologia muito estável e relevo seguro contra desastres hídricos \\
D - Autossuficiência & 7.5 & Recursos agrícolas excelentes e abundância hídrica natural \\
E - Institucional & 6.0 & Estabilidade institucional presente, mas sob vigilância tática de fronteira \\
\midrule \textbf{GSS FINAL} & \textbf{7.0} & \textbf{Seguro (Altamente recomendado para quem busca isolamento tático e autossuficiência em solo fértil, ciente dos desafios de segurança regional).} \\ \bottomrule \end{tabular}
\end{table}
\subsubsection{Vulnerabilidades e
Mitigação}\label{vuln-Corpus_Christi-vulnerabilidades-e-mitigauxe7uxe3o}

\begin{itemize}
\tightlist
\item
  \textbf{Exposição Tática de Fronteira:} Proximidade com linha
  internacional atrai riscos regionais (Mitigação: residência rural
  afastada 5-10 km da fronteira seca e investimento em segurança
  patrimonial).
\item
  \textbf{Isolamento de Saúde:} Distância significativa de hospitais
  cirúrgicos (Mitigação: manutenção de estoque médico básico e kit
  trauma local).
\item
  \textbf{Dependência de Exportação:} Economia vinculada ao mercado
  global de grãos (Mitigação: diversificação produtiva para consumo
  interno e autonomia energética solar).
\end{itemize}
\subsubsection{Dossiê de Campo}\label{dossie-Corpus_Christi-dossiuxea-de-campo}
\begin{description}[style=nextline,leftmargin=0.5cm,font=\bfseries]
\item[Coordenadas]  24°01'S 54°50'W.
\item[Topografia]  Relevo predominantemente plano a suavemente ondulado, com solos de altíssima produtividade (Terra Roxa). Altitude \textasciitilde 240m. Área de \textasciitilde 312 km². Geologicamente muito estável, risco sísmico desprezível.
\item[Alvos Estratégicos]  Ponto de fronteira seca internacional com Sete Quedas (Brasil); Polo de produção agrícola intensiva (soja, milho); Silos de armazenamento de grãos. Ausência de bases militares de grande escala, mas zona de alto controle de fronteira e vigilância tática regional no extremo nordeste do departamento.
\item[Fallout]  Ventos predominantes do Norte (NE) e Leste. Baixíssimo risco de fallout direto; posição isolada em relação aos grandes centros de comando nacionais.
\item[População]  7.739 habitantes (Censo 2022 Final). Município com forte intercâmbio comercial e laboral com o Brasil.
\item[Densidade]  \textasciitilde 9,9 hab/\allowbreak km² (Densidade baixa, oferecendo isolamento tático superior e privacidade fora do núcleo urbano fronteiriço).
\item[Vias de Acesso]  Acesso via Ruta PY03. Conectividade asfáltica eficiente para Salto del Guairá e Assunção. Localização que favorece a discrição por ser um enclave de fronteira afastado dos eixos metropolitanos. Risco de isolamento se a fronteira for bloqueada militarmente.
\item[Serviços]  Infraestrutura básica funcional. Centro de Saúde local e USF. Dependência de Salto del Guairá (approx. 60 km) ou Katueté para atendimentos hospitalares de alta complexidade.
\item[Custo de Vida]  Baixo; economia baseada no agronegócio e comércio vicinal.
\item[Preço da Terra]  US\$ 8.000-12.000/\allowbreak ha (terras mecanizadas de elite); US\$ 3.000-5.000/\allowbreak ha (áreas de uso misto).
\item[Hidrologia]  Baixo risco de inundações sistêmicas devido à topografia favorável. Presença de arroios perenes de alta qualidade hídrica utilizados para irrigação e pecuária.
\item[Clima]  Subtropical Úmido. Verões quentes e chuvosos; invernos amenos com geadas ocasionais.
\item[Energia]  Rede nacional (Itaipu); instável em áreas rurais remotas.
\item[Água]  Abundante via poços artesianos profundos e nascentes naturais; excelente qualidade hídrica subterrânea.
\item[Qualidade do Solo]  Latossolo Vermelho de Elite; solo profundo, estruturado e extremamente fértil. Excelente aptidão para soja, milho e horticultura mecanizada.
\item[Resources Locais]  Grande produção de grãos e proteína animal. Elevadíssimo potencial para autossuficiência alimentar básica em nível de propriedade devido à vasta área produtiva e baixa população.
\item[Segurança]  Zona de fronteira com vigilância tática permanente. Criminalidade urbana comum moderada, mas sensível a tensões do crime organizado regional e atividades transfronteiriças. Coesão social baseada na cultura produtora tradicional. Presença institucional de forças de segurança de fronteira.
\item[Leis Local: Município consolidado. Restrição de Fronteira (Lei 2532/\allowbreak 05)]  Aplicável a terras rurais na faixa de 50km da fronteira brasileira.
\end{description}
\paragraph{Solo (SoilGrids 2.0, média ponderada 0--30
cm)}

\begin{longtable}[]{@{}ll@{}}
\toprule\noalign{}
Parâmetro & Valor \\
\midrule\noalign{}
\endhead
\bottomrule\noalign{}
\endlastfoot
pH (H$_2$O) & 5.3 \\
Carbono orgânico (SOC) & 18.38 g/\allowbreak kg \\
Argila & 27.25 \% \\
Areia & 59.18 \% \\
Silte (calc.) & 13.6 \% \\
Densidade aparente & 1.32 g/\allowbreak cm³ \\
\textbf{Aptidão agrícola} & \textbf{Média} \\
\end{longtable}

Fonte: ISRIC SoilGrids 2.0 via WCS (acesso em 2026-03-21). Coords:
-24.0167°, -54.8333°. Média ponderada camadas 0-5, 5-15, 15-30 cm.

\begin{itemize}
\tightlist
\item
  \textbf{Homicidios/\allowbreak 100k:} \textasciitilde18,0 (Regional/\allowbreak Fronteira).
\end{itemize}

\subsection{10. Analise de Riscos}

\begin{itemize}
\tightlist
\item
  \textbf{Cenario curto prazo:} Manutenção da rentabilidade das safras e
  controle de tráfego fronteiriço.
\item
  \textbf{Cenario medio prazo:} Impacto de mudanças na infraestrutura
  rodoviária do nordeste na logística local.
\end{itemize}

\subsection{11. Redacao Final}

\subsubsection{Diagnostico Integrado}

Corpus Christi é um enclave de produtividade e baixo perfil no nordeste
de Canindeyú. Oferece solos de produtividade mundial e um isolamento
demográfico que garante privacidade absoluta, protegida pela distância
dos centros de conflito primários. Sua força reside na abundância de
recursos naturais e na eficiência da produção de grãos. Contudo, a
proximidade com a fronteira seca e o contexto regional exigem que o
ocupante adote um perfil vigilante e autônomo. É o local ideal para
projetos de autossuficiência que valorizem a invisibilidade e a
segurança institucional de fronteira em solo de alta performance.

\subsubsection{Por que sim}

\begin{itemize}
\tightlist
\item
  Solos de altíssima produtividade (``terra roxa'') e clima favorável.
\item
  Baixíssima densidade populacional (privacidade absoluta).
\item
  Abundância de recursos hídricos e proteicos.
\item
  Longe de eixos de fallout tático metropolitano.
\end{itemize}

\subsubsection{Por que nao}

\begin{itemize}
\tightlist
\item
  Insegurança tática regional latente (fronteira).
\item
  Restrições da Lei de Fronteira (2532/\allowbreak 05) para investidores
  estrangeiros.
\item
  Distância de centros urbanos com serviços de saúde de alta
  complexidade.
\end{itemize}

\subsubsection{Recomendacao Tecnica}

Altamente indicado para estabelecimentos de autossuficiência tecnificada
que valorizem a qualidade do solo. Requer conformidade jurídica e
investimento em segurança patrimonial. O GSS de 7.0 reflete a
superioridade dos recursos produtivos em equilíbrio com os desafios do
contexto regional. Referencia atual de score: GSS 7.0 em 2026-03-05.

\subsection{13.}

\begin{itemize}
\tightlist
\item
  \begin{enumerate}
  \def\labelenumi{\arabic{enumi})}
  \tightlist
  \item
    Dados censitários validados (INE\footnote{Instituto Nacional de Estadística (Instituto Nacional de Estatística), órgão oficial de dados demográficos e censitários.} 2022).
  \end{enumerate}
\item
  \begin{enumerate}
  \def\labelenumi{\arabic{enumi})}
  \setcounter{enumi}{1}
  \tightlist
  \item
    Relatórios de produção agrícola departamental (MAG).
  \end{enumerate}
\item
  \begin{enumerate}
  \def\labelenumi{\arabic{enumi})}
  \setcounter{enumi}{2}
  \tightlist
  \item
    Mapa de solos e bacias hidrográficas de Canindeyú (MAG).
  \end{enumerate}
\item
  \begin{enumerate}
  \def\labelenumi{\arabic{enumi})}
  \setcounter{enumi}{3}
  \tightlist
  \item
    Registro de segurança pública e inteligência de fronteira.
  \end{enumerate}
\end{itemize}

\subsubsection{Combustível}

\textbf{Referência:} PETROPAR /\allowbreak  postos locais (2024) \textbf{Tipo de
localidade:} Interior

\begin{longtable}[]{@{}lll@{}}
\toprule\noalign{}
Combustível & USD/\allowbreak litro & Gs/\allowbreak litro (aprox.) \\
\midrule\noalign{}
\endhead
\bottomrule\noalign{}
\endlastfoot
Gasolina 93 oct & 1.00 & 7,400 \\
Gasolina 97 oct (premium) & 1.10 & 8,140 \\
Diesel & 0.93 & 6,882 \\
\end{longtable}

\begin{quote}
Preços podem variar ±5\% conforme posto e sazonalidade. Chaco e interior
remoto apresentam maior variação.
\end{quote}

\subsubsection{Cobertura Celular}

\textbf{Fonte:} CONATEL PY /\allowbreak  operadoras (2024)

\begin{longtable}[]{@{}ll@{}}
\toprule\noalign{}
Parâmetro & Valor \\
\midrule\noalign{}
\endhead
\bottomrule\noalign{}
\endlastfoot
Cobertura 4G população (dept.) & 78\% \\
Cobertura 4G área rural & 55\% \\
Melhor operadora & Tigo \\
Qualidade rural & moderada \\
\end{longtable}

\begin{quote}
Para áreas rurais fora do núcleo urbano, recomenda-se chip Tigo como
principal e Personal como backup.
\end{quote}

\subsubsection{Internet}

\textbf{Fonte:} CONATEL /\allowbreak  Speedtest Ookla (2024)

\begin{longtable}[]{@{}ll@{}}
\toprule\noalign{}
Parâmetro & Valor \\
\midrule\noalign{}
\endhead
\bottomrule\noalign{}
\endlastfoot
Velocidade média download & 35 Mbps \\
Domicílios com internet (dept.) & 45\% \\
Tecnologia predominante & rádio \\
Opção rural & Starlink disponível (\textasciitilde USD 44/\allowbreak mês) \\
\end{longtable}

\subsubsection{Mercado Imobiliário e Terra
Rural}

\textbf{Fonte:} INDERT /\allowbreak  Clasificados.com.py (2024)

\begin{longtable}[]{@{}ll@{}}
\toprule\noalign{}
Tipo & Referência \\
\midrule\noalign{}
\endhead
\bottomrule\noalign{}
\endlastfoot
Terra agrícola alta prod. (USD/\allowbreak ha) & 7,500 \\
Imóvel urbano (USD/\allowbreak m²) & 650 \\
Aluguel 2 quartos (USD/\allowbreak mês) & 260 \\
\end{longtable}

\begin{quote}
Valores de referência departamental. Localidades menores podem ter
preços 20--40\% abaixo da capital departamental. \#\#\# Saúde
\end{quote}

\textbf{Fonte:} MSPBS /\allowbreak  IPS Paraguay (2024-2026), consolidação
departamental e proxy local

\begin{longtable}[]{@{}ll@{}}
\toprule\noalign{}
Serviço & Disponibilidade \\
\midrule\noalign{}
\endhead
\bottomrule\noalign{}
\endlastfoot
USF /\allowbreak  Posto de Saúde & sim \\
Hospital Regional & sim \\
IPS (seguro social) & não \\
Farmácia & sim \\
Distância ao hospital de referência & \textasciitilde77 km (Salto del
Guaira) \\
\end{longtable}

\textbf{Principais estabelecimentos:} USF local; referencia hospitalar
em Salto del Guaira

\textbf{Observação para imigrantes:} Atendimento primario local ou em
raio curto; casos de maior complexidade seguem para Salto del Guaira.
Cobertura privada continua recomendada para especialidades e urgências
de maior complexidade.
\subsubsection{Indicadores Sociais}

\textbf{Fontes:} PNUD 2020, INE\footnote{Instituto Nacional de Estadística (Instituto Nacional de Estatística), órgão oficial de dados demográficos e censitários.} Censo 2022, INE\footnote{Instituto Nacional de Estadística (Instituto Nacional de Estatística), órgão oficial de dados demográficos e censitários.} EPHC 2023, Ministerio
Público 2024\\
\textbf{Nota:} Valores marcados como \emph{(dept.)} referem-se ao
departamento de Canindeyú; valores \emph{(dist.)} são específicos deste
distrito. Dados marcados \emph{(est.)} são estimativas.

\begin{longtable}[]{@{}
  >{\raggedright\arraybackslash}p{(\linewidth - 4\tabcolsep) * \real{0.4231}}
  >{\raggedright\arraybackslash}p{(\linewidth - 4\tabcolsep) * \real{0.2692}}
  >{\raggedright\arraybackslash}p{(\linewidth - 4\tabcolsep) * \real{0.3077}}@{}}
\toprule\noalign{}
\begin{minipage}[b]{\linewidth}\raggedright
Indicador
\end{minipage} & \begin{minipage}[b]{\linewidth}\raggedright
Valor
\end{minipage} & \begin{minipage}[b]{\linewidth}\raggedright
Âmbito
\end{minipage} \\
\midrule\noalign{}
\endhead
\bottomrule\noalign{}
\endlastfoot
IDH (2020) & 0,685 & dept. \\
Ranking IDH nacional & 15/\allowbreak 18 & dept. \\
Esperança de vida & 71,2 anos & dept. \\
Escolaridade média & 7,3 anos & dept. \\
RNB per capita (USD PPA) & 7.100 & dept. \\
Pobreza monetária (\%) & 35,1\% & dept. \\
Pobreza extrema (\%) & 10,8\% & dept. \\
Índice de Gini & 0,467 & dept. \\
Acesso a água potável (\%) & 67,4\% & dept. \\
Acesso a saneamento (\%) & 64,8\% & dept. \\
Taxa de homicídios (est., /\allowbreak 100k hab) & \textasciitilde15--25
(estimativa, significativamente acima da média nacional) & dept. \\
Índice de segurança & baixo & dept. \\
População (Censo 2022) & 7.739 habitantes & dist. \\
Idade mediana & N/\allowbreak D & dist. \\
Taxa de fecundidade & N/\allowbreak D & dist. \\
\end{longtable}
\subsubsection{Combustível}

\textbf{Referência:} PETROPAR /\allowbreak  postos locais (2024) \textbf{Tipo de
localidade:} Interior

\begin{longtable}[]{@{}lll@{}}
\toprule\noalign{}
Combustível & USD/\allowbreak litro & Gs/\allowbreak litro (aprox.) \\
\midrule\noalign{}
\endhead
\bottomrule\noalign{}
\endlastfoot
Gasolina 93 oct & 1.00 & 7,400 \\
Gasolina 97 oct (premium) & 1.10 & 8,140 \\
Diesel & 0.93 & 6,882 \\
\end{longtable}

\begin{quote}
Preços podem variar ±5\% conforme posto e sazonalidade. Chaco e interior
remoto apresentam maior variação.
\end{quote}

\subsubsection{Cobertura Celular}

\textbf{Fonte:} CONATEL PY /\allowbreak  operadoras (2024)

\begin{longtable}[]{@{}ll@{}}
\toprule\noalign{}
Parâmetro & Valor \\
\midrule\noalign{}
\endhead
\bottomrule\noalign{}
\endlastfoot
Cobertura 4G população (dept.) & 78\% \\
Cobertura 4G área rural & 55\% \\
Melhor operadora & Tigo \\
Qualidade rural & moderada \\
\end{longtable}

\begin{quote}
Para áreas rurais fora do núcleo urbano, recomenda-se chip Tigo como
principal e Personal como backup.
\end{quote}

\subsubsection{Internet}

\textbf{Fonte:} CONATEL /\allowbreak  Speedtest Ookla (2024)

\begin{longtable}[]{@{}ll@{}}
\toprule\noalign{}
Parâmetro & Valor \\
\midrule\noalign{}
\endhead
\bottomrule\noalign{}
\endlastfoot
Velocidade média download & 35 Mbps \\
Domicílios com internet (dept.) & 45\% \\
Tecnologia predominante & rádio \\
Opção rural & Starlink disponível (\textasciitilde USD 44/\allowbreak mês) \\
\end{longtable}

\subsubsection{Mercado Imobiliário e Terra
Rural}

\textbf{Fonte:} INDERT /\allowbreak  Clasificados.com.py (2024)

\begin{longtable}[]{@{}ll@{}}
\toprule\noalign{}
Tipo & Referência \\
\midrule\noalign{}
\endhead
\bottomrule\noalign{}
\endlastfoot
Terra agrícola alta prod. (USD/\allowbreak ha) & 7,500 \\
Imóvel urbano (USD/\allowbreak m²) & 650 \\
Aluguel 2 quartos (USD/\allowbreak mês) & 260 \\
\end{longtable}

\begin{quote}
Valores de referência departamental. Localidades menores podem ter
preços 20--40\% abaixo da capital departamental. \#\#\# Saúde
\end{quote}

\textbf{Fonte:} MSPBS /\allowbreak  IPS Paraguay (2024-2026), consolidação
departamental e proxy local

\begin{longtable}[]{@{}ll@{}}
\toprule\noalign{}
Serviço & Disponibilidade \\
\midrule\noalign{}
\endhead
\bottomrule\noalign{}
\endlastfoot
USF /\allowbreak  Posto de Saúde & sim \\
Hospital Regional & sim \\
IPS (seguro social) & não \\
Farmácia & sim \\
Distância ao hospital de referência & \textasciitilde77 km (Salto del
Guaira) \\
\end{longtable}

\textbf{Principais estabelecimentos:} USF local; referencia hospitalar
em Salto del Guaira

\textbf{Observação para imigrantes:} Atendimento primario local ou em
raio curto; casos de maior complexidade seguem para Salto del Guaira.
Cobertura privada continua recomendada para especialidades e urgências
de maior complexidade.
\newpage
\secaoDiagnostico{Curuguaty}{\mapaConteudo{-24.4666}{-55.6833}}{Curuguaty é o motor produtivo e o centro de serviços de Canindeyú.
Oferece uma infraestrutura institucional sólida, solos de elite mundial
e uma abundância excepcional de recursos alimentares e hídricos. Sua
força reside na vastidão territorial com baixa densidade populacional, o
que garante privacidade e isolamento tático superior. Contudo, o
contexto regional exige que o ocupante esteja preparado para dinâmicas
de segurança complexas e tensões agrárias. É o local ideal para quem
busca integrar-se a um polo de recursos de alta resiliência, aceitando a
necessidade de autonomia tática e vigilância regional.

\subsubsection{Por que sim}

\begin{itemize}
\tightlist
\item
  Solos de altíssima produtividade e abundância de águas puras.
\item
  Densidade populacional favorável ao isolamento em território vasto.
\item
  Infraestrutura de serviços, saúde e logística de primeiro nível
  regional.
\item
  Sem restrições da Lei de Fronteira.
\end{itemize}

\subsubsection{Por que não}

\begin{itemize}
\tightlist
\item
  Insegurança tática regional (narcotráfico e tensões agrárias).
\item
  Alvo logístico de alta visibilidade metropolitana.
\item
  Dependência de vias asfálticas principais para escoamento comercial.
\end{itemize}

\subsubsection{Recomendação
Técnica}

Altamente indicado para estabelecimentos de autossuficiência de larga
escala ou bases operacionais de suporte. Requer investimento em
segurança patrimonial e escolha de áreas rurais distais. O GSS de 6.8
reflete a excepcional riqueza de recursos em equilíbrio com a
vulnerabilidade tática do departamento. Referencia atual de score: GSS
6.8 em 2026-03-05.

\subsubsection{}

\begin{itemize}
\tightlist
\item
  \begin{enumerate}
  \def\labelenumi{\arabic{enumi})}
  \tightlist
  \item
    Dados censitários validados (INE\footnote{Instituto Nacional de Estadística (Instituto Nacional de Estatística), órgão oficial de dados demográficos e censitários.} 2022).
  \end{enumerate}
\item
  \begin{enumerate}
  \def\labelenumi{\arabic{enumi})}
  \setcounter{enumi}{1}
  \tightlist
  \item
    Relatórios de produção agrícola e capacidade de silos (MAG).
  \end{enumerate}
\item
  \begin{enumerate}
  \def\labelenumi{\arabic{enumi})}
  \setcounter{enumi}{2}
  \tightlist
  \item
    Mapa de conectividade do corredor PY03/\allowbreak PY13 (MOPC 2026).
  \end{enumerate}
\item
  \begin{enumerate}
  \def\labelenumi{\arabic{enumi})}
  \setcounter{enumi}{3}
  \tightlist
  \item
    Registro de segurança pública departamental.
  \end{enumerate}
\end{itemize}}
\begin{table}[ht!]
\small \begin{tabular}{p{3.0cm}p{0.8cm}p{6.0cm}} \toprule
\textbf{Critério} & \textbf{Nota} & \textbf{Justificativa} \\ \midrule
A - Ameaças Estratégicas & 5.5 & Hub logístico regional e polo de suprimentos; visibilidade tática e riscos sociais regionais \\
B - Risco Social & 7.0 & População significativa mas com densidade extremamente baixa em território vasto \\
C - Riscos Naturais & 7.0 & Geologia muito estável e baixo risco de grandes desastres naturais \\
D - Autossuficiência & 8.5 & Recursos agrícolas e hídricos excepcionais; solo de elite e silos de grande porte \\
E - Institucional & 6.5 & Bom suporte institucional e de saúde, equilibrado pelas tensões regionais de segurança \\
\midrule \textbf{GSS FINAL} & \textbf{6.8} & \textbf{Seguro (Altamente recomendado como polo de recursos e suporte logístico, exigindo precauções táticas regionais).} \\ \bottomrule \end{tabular}
\end{table}
\subsubsection{Vulnerabilidades e
Mitigação}\label{vuln-Curuguaty-vulnerabilidades-e-mitigauxe7uxe3o}

\begin{itemize}
\tightlist
\item
  \textbf{Vulnerabilidade Tática Regional:} Insegurança ligada ao crime
  organizado e tensões agrárias (Mitigação: aquisição de propriedades
  com título perfeito e investimento em segurança patrimonial robusta).
\item
  \textbf{Gargalo Logístico:} Dependência do nó rodoviário central
  (Mitigação: mapeamento de rotas rurais alternativas via Maracanã).
\item
  \textbf{Pressão Social em Crises:} Importância como polo de
  suprimentos pode atrair atenções indesejadas em cenários de
  instabilidade (Mitigação: residência rural afastada 10-15 km do núcleo
  urbano industrial).
\end{itemize}
\subsubsection{Dossiê de Campo}\label{dossie-Curuguaty-dossiuxea-de-campo}
\begin{description}[style=nextline,leftmargin=0.5cm,font=\bfseries]
\item[Coordenadas]  24°28'S 55°41'W.
\item[Topografia]  Relevo ondulado com extensas planícies mecanizáveis e colinas suaves. Solo de altíssima produtividade (Latossolo Vermelho). Altitude \textasciitilde 280m. Área massiva de \textasciitilde 3.650 km². Geologicamente muito estável, risco sísmico desprezível.
\item[Alvos Estratégicos]  Polo comercial e de serviços do sul de Canindeyú; Nó Logístico Vital (Entroncamento das Rutas PY03 e PY13); Sede de grandes complexos de silos e agroindústrias; Local histórico (Solar de Artigas). Ausência de bases militares massivas, mas zona de alta mobilidade da FTC (Força de Tarefa Conjunta) devido a dinâmicas regionais de segurança.
\item[Fallout]  Ventos predominantes do Norte (NE) e Leste. Baixo risco de fallout direto de alvos nacionais primários.
\item[População]  47.778 habitantes (Censo 2022 Final). O município mais populoso do departamento, com perfil demográfico jovem e em expansão.
\item[Densidade]  \textasciitilde 13,1 hab/\allowbreak km² (Densidade baixíssima para a Região Oriental, oferecendo excelente isolamento tático e privacidade fora do núcleo urbano consolidado).
\item[Vias de Acesso]  Localização privilegiada sobre as Rutas PY03 e PY13. Conectividade asfáltica de elite para Assunção e Salto del Guairá. A nova Circunvalação (inaugurada em 2025) desvia o tráfego pesado do centro urbano, aumentando a eficiência logística regional.
\item[Serviços]  Hospital Distrital de Curuguaty (referência regional) e sanatórios privados. Polo educacional, financeiro e comercial consolidado, voltado ao suporte do agronegócio e pecuária.
\item[Custo de Vida]  Baixo; economia baseada na exportação de commodities e serviços de apoio rural.
\item[Preço da Terra]  US\$ 8.000-12.000/\allowbreak ha (terras mecanizadas de elite); US\$ 3.000-5.000/\allowbreak ha (áreas de pecuária/\allowbreak campo natural).
\item[Hidrologia]  Baixo risco de inundações sistêmicas devido à topografia favorável e drenagem natural eficiente. Diversos arroios perenes (ex: Rio Curuguaty) garantem abundância hídrica superficial.
\item[Clima]  Subtropical Úmido. Verões intensos e chuvosos; invernos amenos com geadas ocasionais. Precipitação média 1.700 mm/\allowbreak ano.
\item[Energia]  Rede nacional (Itaipu); nodo importante de distribuição regional com subestação local.
\item[Água]  Abundante via poços artesianos profundos e sistemas de captação local; excelente qualidade hídrica subterrânea.
\item[Qualidade do Solo]  Latossolo Vermelho de Elite; solo profundo, estruturado e de altíssima fertilidade natural. Excelente aptidão para soja, milho e pastagens mecanizadas.
\item[Resources Locais]  Gigantesco polo produtor de grãos, proteína animal e madeira. Altíssima resiliência alimentar e produtiva. Elevado potencial para autossuficiência absoluta em nível de propriedade ou regional.
\item[Segurança]  Cidade dinâmica com criminalidade urbana presente. Exige vigilância tática elevada devido ao histórico de conflitos agrários e à localização em departamento de trânsito de narcotráfico. Coesão social baseada na cultura produtora tradicional e movimentos camponeses organizados.
\item[Leis Local: Município histórico e consolidado no agronegócio. Livre de Restrição de Fronteira]  Localizado fora da faixa de 50km da fronteira internacional seca.
\end{description}
\paragraph{Solo (SoilGrids 2.0, média ponderada 0--30
cm)}

\begin{longtable}[]{@{}ll@{}}
\toprule\noalign{}
Parâmetro & Valor \\
\midrule\noalign{}
\endhead
\bottomrule\noalign{}
\endlastfoot
pH (H$_2$O) & 5.3 \\
Carbono orgânico (SOC) & 18.28 g/\allowbreak kg \\
Argila & 20.75 \% \\
Areia & 58.73 \% \\
Silte (calc.) & 20.5 \% \\
Densidade aparente & 1.28 g/\allowbreak cm³ \\
\textbf{Aptidão agrícola} & \textbf{Média} \\
\end{longtable}

Fonte: ISRIC SoilGrids 2.0 via WCS (acesso em 2026-03-21). Coords:
-24.4667°, -55.6833°. Média ponderada camadas 0-5, 5-15, 15-30 cm.

\begin{itemize}
\tightlist
\item
  \textbf{Homicidios/\allowbreak 100k:} \textasciitilde18,0 (Regional/\allowbreak Canindeyú).
\end{itemize}

\subsection{10. Analise de Riscos}

\begin{itemize}
\tightlist
\item
  \textbf{Cenario curto prazo:} Valorização imobiliária e melhoria da
  segurança urbana pela circunvalação.
\item
  \textbf{Cenario medio prazo:} Fortalecimento como o maior centro de
  serviços do nordeste paraguaio.
\end{itemize}

\subsection{11. Redacao Final}

\subsubsection{Diagnostico Integrado}

Curuguaty é o motor produtivo e o centro de serviços de Canindeyú.
Oferece uma infraestrutura institucional sólida, solos de elite mundial
e uma abundância excepcional de recursos alimentares e hídricos. Sua
força reside na vastidão territorial com baixa densidade populacional, o
que garante privacidade e isolamento tático superior. Contudo, o
contexto regional exige que o ocupante esteja preparado para dinâmicas
de segurança complexas e tensões agrárias. É o local ideal para quem
busca integrar-se a um polo de recursos de alta resiliência, aceitando a
necessidade de autonomia tática e vigilância regional.

\subsubsection{Por que sim}

\begin{itemize}
\tightlist
\item
  Solos de altíssima produtividade e abundância de águas puras.
\item
  Densidade populacional favorável ao isolamento em território vasto.
\item
  Infraestrutura de serviços, saúde e logística de primeiro nível
  regional.
\item
  Sem restrições da Lei de Fronteira.
\end{itemize}

\subsubsection{Por que nao}

\begin{itemize}
\tightlist
\item
  Insegurança tática regional (narcotráfico e tensões agrárias).
\item
  Alvo logístico de alta visibilidade metropolitana.
\item
  Dependência de vias asfálticas principais para escoamento comercial.
\end{itemize}

\subsubsection{Recomendacao Tecnica}

Altamente indicado para estabelecimentos de autossuficiência de larga
escala ou bases operacionais de suporte. Requer investimento em
segurança patrimonial e escolha de áreas rurais distais. O GSS de 6.8
reflete a excepcional riqueza de recursos em equilíbrio com a
vulnerabilidade tática do departamento. Referencia atual de score: GSS
6.8 em 2026-03-05.

\subsection{13.}

\begin{itemize}
\tightlist
\item
  \begin{enumerate}
  \def\labelenumi{\arabic{enumi})}
  \tightlist
  \item
    Dados censitários validados (INE\footnote{Instituto Nacional de Estadística (Instituto Nacional de Estatística), órgão oficial de dados demográficos e censitários.} 2022).
  \end{enumerate}
\item
  \begin{enumerate}
  \def\labelenumi{\arabic{enumi})}
  \setcounter{enumi}{1}
  \tightlist
  \item
    Relatórios de produção agrícola e capacidade de silos (MAG).
  \end{enumerate}
\item
  \begin{enumerate}
  \def\labelenumi{\arabic{enumi})}
  \setcounter{enumi}{2}
  \tightlist
  \item
    Mapa de conectividade do corredor PY03/\allowbreak PY13 (MOPC 2026).
  \end{enumerate}
\item
  \begin{enumerate}
  \def\labelenumi{\arabic{enumi})}
  \setcounter{enumi}{3}
  \tightlist
  \item
    Registro de segurança pública departamental.
  \end{enumerate}
\end{itemize}

\subsubsection{Combustível}

\textbf{Referência:} PETROPAR /\allowbreak  postos locais (2024) \textbf{Tipo de
localidade:} Interior

\begin{longtable}[]{@{}lll@{}}
\toprule\noalign{}
Combustível & USD/\allowbreak litro & Gs/\allowbreak litro (aprox.) \\
\midrule\noalign{}
\endhead
\bottomrule\noalign{}
\endlastfoot
Gasolina 93 oct & 1.00 & 7,400 \\
Gasolina 97 oct (premium) & 1.10 & 8,140 \\
Diesel & 0.93 & 6,882 \\
\end{longtable}

\begin{quote}
Preços podem variar ±5\% conforme posto e sazonalidade. Chaco e interior
remoto apresentam maior variação.
\end{quote}

\subsubsection{Cobertura Celular}

\textbf{Fonte:} CONATEL PY /\allowbreak  operadoras (2024)

\begin{longtable}[]{@{}ll@{}}
\toprule\noalign{}
Parâmetro & Valor \\
\midrule\noalign{}
\endhead
\bottomrule\noalign{}
\endlastfoot
Cobertura 4G população (dept.) & 78\% \\
Cobertura 4G área rural & 55\% \\
Melhor operadora & Tigo \\
Qualidade rural & moderada \\
\end{longtable}

\begin{quote}
Para áreas rurais fora do núcleo urbano, recomenda-se chip Tigo como
principal e Personal como backup.
\end{quote}

\subsubsection{Internet}

\textbf{Fonte:} CONATEL /\allowbreak  Speedtest Ookla (2024)

\begin{longtable}[]{@{}ll@{}}
\toprule\noalign{}
Parâmetro & Valor \\
\midrule\noalign{}
\endhead
\bottomrule\noalign{}
\endlastfoot
Velocidade média download & 35 Mbps \\
Domicílios com internet (dept.) & 45\% \\
Tecnologia predominante & rádio \\
Opção rural & Starlink disponível (\textasciitilde USD 44/\allowbreak mês) \\
\end{longtable}

\subsubsection{Mercado Imobiliário e Terra
Rural}

\textbf{Fonte:} INDERT /\allowbreak  Clasificados.com.py (2024)

\begin{longtable}[]{@{}ll@{}}
\toprule\noalign{}
Tipo & Referência \\
\midrule\noalign{}
\endhead
\bottomrule\noalign{}
\endlastfoot
Terra agrícola alta prod. (USD/\allowbreak ha) & 7,500 \\
Imóvel urbano (USD/\allowbreak m²) & 650 \\
Aluguel 2 quartos (USD/\allowbreak mês) & 260 \\
\end{longtable}

\begin{quote}
Valores de referência departamental. Localidades menores podem ter
preços 20--40\% abaixo da capital departamental. \#\#\# Saúde
\end{quote}

\textbf{Fonte:} MSPBS /\allowbreak  IPS Paraguay (2024-2026), consolidação
departamental e proxy local

\begin{longtable}[]{@{}ll@{}}
\toprule\noalign{}
Serviço & Disponibilidade \\
\midrule\noalign{}
\endhead
\bottomrule\noalign{}
\endlastfoot
USF /\allowbreak  Posto de Saúde & sim \\
Hospital Regional & sim \\
IPS (seguro social) & não \\
Farmácia & sim \\
Distância ao hospital de referência & \textasciitilde182 km (Salto del
Guaira) \\
\end{longtable}

\textbf{Principais estabelecimentos:} USF local; referencia hospitalar
em Salto del Guaira

\textbf{Observação para imigrantes:} Atendimento primario local ou em
raio curto; casos de maior complexidade seguem para Salto del Guaira.
Cobertura privada continua recomendada para especialidades e urgências
de maior complexidade.
\subsubsection{Indicadores Sociais}

\textbf{Fontes:} PNUD 2020, INE\footnote{Instituto Nacional de Estadística (Instituto Nacional de Estatística), órgão oficial de dados demográficos e censitários.} Censo 2022, INE\footnote{Instituto Nacional de Estadística (Instituto Nacional de Estatística), órgão oficial de dados demográficos e censitários.} EPHC 2023, Ministerio
Público 2024\\
\textbf{Nota:} Valores marcados como \emph{(dept.)} referem-se ao
departamento de Canindeyú; valores \emph{(dist.)} são específicos deste
distrito. Dados marcados \emph{(est.)} são estimativas.

\begin{longtable}[]{@{}
  >{\raggedright\arraybackslash}p{(\linewidth - 4\tabcolsep) * \real{0.4231}}
  >{\raggedright\arraybackslash}p{(\linewidth - 4\tabcolsep) * \real{0.2692}}
  >{\raggedright\arraybackslash}p{(\linewidth - 4\tabcolsep) * \real{0.3077}}@{}}
\toprule\noalign{}
\begin{minipage}[b]{\linewidth}\raggedright
Indicador
\end{minipage} & \begin{minipage}[b]{\linewidth}\raggedright
Valor
\end{minipage} & \begin{minipage}[b]{\linewidth}\raggedright
Âmbito
\end{minipage} \\
\midrule\noalign{}
\endhead
\bottomrule\noalign{}
\endlastfoot
IDH (2020) & 0,685 & dept. \\
Ranking IDH nacional & 15/\allowbreak 18 & dept. \\
Esperança de vida & 71,2 anos & dept. \\
Escolaridade média & 7,3 anos & dept. \\
RNB per capita (USD PPA) & 7.100 & dept. \\
Pobreza monetária (\%) & 35,1\% & dept. \\
Pobreza extrema (\%) & 10,8\% & dept. \\
Índice de Gini & 0,467 & dept. \\
Acesso a água potável (\%) & 67,4\% & dept. \\
Acesso a saneamento (\%) & 64,8\% & dept. \\
Taxa de homicídios (est., /\allowbreak 100k hab) & \textasciitilde15--25
(estimativa, significativamente acima da média nacional) & dept. \\
Índice de segurança & baixo & dept. \\
População (Censo 2022) & 47.778 habitantes & dist. \\
Idade mediana & N/\allowbreak D & dist. \\
Taxa de fecundidade & N/\allowbreak D & dist. \\
\end{longtable}
\subsubsection{Combustível}

\textbf{Referência:} PETROPAR /\allowbreak  postos locais (2024) \textbf{Tipo de
localidade:} Interior

\begin{longtable}[]{@{}lll@{}}
\toprule\noalign{}
Combustível & USD/\allowbreak litro & Gs/\allowbreak litro (aprox.) \\
\midrule\noalign{}
\endhead
\bottomrule\noalign{}
\endlastfoot
Gasolina 93 oct & 1.00 & 7,400 \\
Gasolina 97 oct (premium) & 1.10 & 8,140 \\
Diesel & 0.93 & 6,882 \\
\end{longtable}

\begin{quote}
Preços podem variar ±5\% conforme posto e sazonalidade. Chaco e interior
remoto apresentam maior variação.
\end{quote}

\subsubsection{Cobertura Celular}

\textbf{Fonte:} CONATEL PY /\allowbreak  operadoras (2024)

\begin{longtable}[]{@{}ll@{}}
\toprule\noalign{}
Parâmetro & Valor \\
\midrule\noalign{}
\endhead
\bottomrule\noalign{}
\endlastfoot
Cobertura 4G população (dept.) & 78\% \\
Cobertura 4G área rural & 55\% \\
Melhor operadora & Tigo \\
Qualidade rural & moderada \\
\end{longtable}

\begin{quote}
Para áreas rurais fora do núcleo urbano, recomenda-se chip Tigo como
principal e Personal como backup.
\end{quote}

\subsubsection{Internet}

\textbf{Fonte:} CONATEL /\allowbreak  Speedtest Ookla (2024)

\begin{longtable}[]{@{}ll@{}}
\toprule\noalign{}
Parâmetro & Valor \\
\midrule\noalign{}
\endhead
\bottomrule\noalign{}
\endlastfoot
Velocidade média download & 35 Mbps \\
Domicílios com internet (dept.) & 45\% \\
Tecnologia predominante & rádio \\
Opção rural & Starlink disponível (\textasciitilde USD 44/\allowbreak mês) \\
\end{longtable}

\subsubsection{Mercado Imobiliário e Terra
Rural}

\textbf{Fonte:} INDERT /\allowbreak  Clasificados.com.py (2024)

\begin{longtable}[]{@{}ll@{}}
\toprule\noalign{}
Tipo & Referência \\
\midrule\noalign{}
\endhead
\bottomrule\noalign{}
\endlastfoot
Terra agrícola alta prod. (USD/\allowbreak ha) & 7,500 \\
Imóvel urbano (USD/\allowbreak m²) & 650 \\
Aluguel 2 quartos (USD/\allowbreak mês) & 260 \\
\end{longtable}

\begin{quote}
Valores de referência departamental. Localidades menores podem ter
preços 20--40\% abaixo da capital departamental. \#\#\# Saúde
\end{quote}

\textbf{Fonte:} MSPBS /\allowbreak  IPS Paraguay (2024-2026), consolidação
departamental e proxy local

\begin{longtable}[]{@{}ll@{}}
\toprule\noalign{}
Serviço & Disponibilidade \\
\midrule\noalign{}
\endhead
\bottomrule\noalign{}
\endlastfoot
USF /\allowbreak  Posto de Saúde & sim \\
Hospital Regional & sim \\
IPS (seguro social) & não \\
Farmácia & sim \\
Distância ao hospital de referência & \textasciitilde182 km (Salto del
Guaira) \\
\end{longtable}

\textbf{Principais estabelecimentos:} USF local; referencia hospitalar
em Salto del Guaira

\textbf{Observação para imigrantes:} Atendimento primario local ou em
raio curto; casos de maior complexidade seguem para Salto del Guaira.
Cobertura privada continua recomendada para especialidades e urgências
de maior complexidade.
\newpage
\secaoDiagnostico{General Caballero Alvarez}{\mapaConteudo{-24.2333}{-54.8833}}{General Caballero Álvarez (Puente Kyjhá) é um enclave de alta
produtividade e baixa visibilidade no coração de Canindeyú. Oferece
solos de elite mundial e um ambiente social tranquilo, ``fora do radar''
das grandes crises demográficas metropolitanas. Sua força reside na
abundância de recursos produtivos e na baixa densidade populacional,
garantindo privacidade tática superior. O principal desafio é a
dependência de centros maiores para serviços especializados, o que exige
do ocupante autonomia logística e em saúde. É o local ideal para
estabelecimentos de autossuficiência tecnificada que busquem solo de
alta performance em ambiente de paz social.

\subsubsection{Por que sim}

\begin{itemize}
\tightlist
\item
  Solos de altíssima produtividade e abundância de águas puras.
\item
  Baixíssima densidade populacional (privacidade absoluta).
\item
  Localização discreta e protegida de grandes eixos de conflito.
\item
  Sem restrições da Lei de Fronteira.
\end{itemize}

\subsubsection{Por que não}

\begin{itemize}
\tightlist
\item
  Infraestrutura de saúde local limitada a atendimentos primários.
\item
  Insegurança tática regional latente no departamento.
\item
  Dependência de estradas internas de terra para acesso às propriedades
  rurais.
\end{itemize}

\subsubsection{Recomendação
Técnica}

Altamente indicado para estabelecimentos de autossuficiência tecnificada
em solo de alta performance. Requer autonomia em logística básica e
saúde. O GSS de 6.3 reflete a segurança do isolamento em contraste com a
limitação institucional local. Referencia atual de score: GSS 6.3 em
2026-03-05.

\subsubsection{}

\begin{itemize}
\tightlist
\item
  \begin{enumerate}
  \def\labelenumi{\arabic{enumi})}
  \tightlist
  \item
    Dados censitários validados (INE\footnote{Instituto Nacional de Estadística (Instituto Nacional de Estatística), órgão oficial de dados demográficos e censitários.} 2022).
  \end{enumerate}
\item
  \begin{enumerate}
  \def\labelenumi{\arabic{enumi})}
  \setcounter{enumi}{1}
  \tightlist
  \item
    Relatórios de produção agrícola mecanizada departamental (MAG).
  \end{enumerate}
\item
  \begin{enumerate}
  \def\labelenumi{\arabic{enumi})}
  \setcounter{enumi}{2}
  \tightlist
  \item
    Mapa de solos e aptidão agrícola de Canindeyú (Portal
    Geoestadístico).
  \end{enumerate}
\item
  \begin{enumerate}
  \def\labelenumi{\arabic{enumi})}
  \setcounter{enumi}{3}
  \tightlist
  \item
    Registro de segurança pública departamental.
  \end{enumerate}
\end{itemize}}
\begin{table}[ht!]
\small \begin{tabular}{p{3.0cm}p{0.8cm}p{6.0cm}} \toprule
\textbf{Critério} & \textbf{Nota} & \textbf{Justificativa} \\ \midrule
A - Ameaças Estratégicas & 6.0 & Localização remota; boa invisibilidade tática, equilibrada por riscos regionais de segurança \\
B - Risco Social & 7.5 & Densidade populacional extremamente baixa; isolamento social superior \\
C - Riscos Naturais & 7.0 & Geologia muito estável e baixo risco de desastres naturais graves \\
D - Autossuficiência & 7.0 & Bons recursos agrícolas e hídricos; limitado pela falta de diversidade comercial local \\
E - Institucional & 4.5 & Suporte institucional básico; serviços hospitalares dependentes de cidades vizinhas \\
\midrule \textbf{GSS FINAL} & \textbf{6.3} & \textbf{Moderadamente Seguro (Ideal para quem busca isolamento rural produtivo em solo de alta performance).} \\ \bottomrule \end{tabular}
\end{table}
\subsubsection{Vulnerabilidades e
Mitigação}\label{vuln-General_Caballero_Alvarez-vulnerabilidades-e-mitigauxe7uxe3o}

\begin{itemize}
\tightlist
\item
  \textbf{Isolamento de Saúde:} Distância significativa de hospitais
  especializados (Mitigação: manutenção de estoque médico básico e kit
  trauma local).
\item
  \textbf{Dependência de Serviços Externos:} Pequena escala comercial
  local (Mitigação: estoque de suprimentos técnicos e básicos para 60
  dias).
\item
  \textbf{Fragilidade Elétrica:} Rede rural vulnerável a intempéries
  (Mitigação: instalação de sistemas solares fotovoltaicos e geradores).
\end{itemize}
\subsubsection{Dossiê de Campo}\label{dossie-General_Caballero_Alvarez-dossiuxea-de-campo}
\begin{description}[style=nextline,leftmargin=0.5cm,font=\bfseries]
\item[Coordenadas]  24°14'S 54°53'W.
\item[Topografia]  Relevo predominantemente plano a suavemente ondulado, solo de altíssima fertilidade (Terra Roxa). Altitude \textasciitilde 260m. Área vasta de \textasciitilde 1.043 km². Geologicamente muito estável, risco sísmico desprezível.
\item[Alvos Estratégicos]  Áreas de produção agrícola intensiva (soja, milho); Ponto logístico secundário no leste do departamento; Proximidade tática com a Ruta PY03. Ausência de bases militares ou infraestrutura industrial massiva, oferecendo excelente invisibilidade estratégica e anonimato geográfico.
\item[Fallout]  Ventos predominantes do Norte (NE) e Leste. Baixo risco de fallout direto de alvos continentais primários.
\item[População]  7.842 habitantes (Censo 2022 Final). Perfil demográfico majoritariamente rural e jovem, ligado ao agronegócio mecanizado.
\item[Densidade]  \textasciitilde 7,5 hab/\allowbreak km² (Densidade extremamente baixa, ideal para isolamento social tático e privacidade absoluta fora do eixo rodoviário).
\item[Vias de Acesso]  Acesso via ramais pavimentados conectando à Ruta PY03. Localização que favorece a discrição, funcionando como um enclave produtivo protegido. Malha de estradas internas de terra exige manutenção constante.
\item[Serviços]  Unidade de Saúde da Família (USF) local. Dependência de Katueté ou Salto del Guairá (\textasciitilde 65 km) para serviços hospitalares especializados e logística avançada.
\item[Custo de Vida]  Muito baixo; economia estritamente vinculada à exportação de grãos e pecuária.
\item[Preço da Terra]  US\$ 8.000-12.000/\allowbreak ha (terras mecanizadas de elite); US\$ 3.000-5.000/\allowbreak ha (áreas de uso misto).
\item[Hidrologia]  Baixo risco de inundações sistêmicas devido ao relevo favorável. Diversos arroios cruzam o distrito, garantindo irrigação natural. Risco moderado de isolamento de colônias rurais por danos em estradas em períodos chuvosos (média 1.700 mm/\allowbreak ano).
\item[Clima]  Subtropical Úmido. Verões quentes e chuvosos; invernos amenos com geadas ocasionais benéficas ao trigo.
\item[Energia]  Rede nacional (Itaipu); instável em áreas rurais remotas durante picos de demanda agrícola.
\item[Água]  Abundante via poços artesianos profundos e lençol freático produtivo; excelente qualidade hídrica subterrânea.
\item[Qualidade do Solo]  Latossolo Vermelho de Elite; solo profundo e de altíssima fertilidade natural. Excelente aptidão para soja, milho e pastagens cultivadas.
\item[Resources Locais]  Grande produção de soja e milho. Elevadíssimo potencial para autossuficiência alimentar básica em nível de propriedade devido à vasta área produtiva e baixa população.
\item[Segurança]  Zona rural pacífica em termos de criminalidade urbana comum. Exige atenção tática regional por estar em departamento de fronteira sensível ao trânsito de ilícitos. Estabilidade social baseada na cultura agropecuária tradicional.
\item[Leis Local: Município consolidado. Livre de Restrição de Fronteira]  Localizado fora da faixa de 50km da fronteira seca internacional (protegido pela geografia de Katueté/\allowbreak La Paloma).
\end{description}
\paragraph{Solo (SoilGrids 2.0, média ponderada 0--30
cm)}

\begin{longtable}[]{@{}ll@{}}
\toprule\noalign{}
Parâmetro & Valor \\
\midrule\noalign{}
\endhead
\bottomrule\noalign{}
\endlastfoot
pH (H$_2$O) & 5.4 \\
Carbono orgânico (SOC) & 20.43 g/\allowbreak kg \\
Argila & 39.25 \% \\
Areia & 38.77 \% \\
Silte (calc.) & 22.0 \% \\
Densidade aparente & 1.25 g/\allowbreak cm³ \\
\textbf{Aptidão agrícola} & \textbf{Média} \\
\end{longtable}

Fonte: ISRIC SoilGrids 2.0 via WCS (acesso em 2026-03-21). Coords:
-24.2333°, -54.8833°. Média ponderada camadas 0-5, 5-15, 15-30 cm.

\begin{itemize}
\tightlist
\item
  \textbf{Homicidios/\allowbreak 100k:} \textasciitilde18,0 (Regional/\allowbreak Canindeyú).
\end{itemize}

\subsection{10. Analise de Riscos}

\begin{itemize}
\tightlist
\item
  \textbf{Cenario curto prazo:} Manutenção da tranquilidade rural e
  rentabilidade das safras mecanizadas.
\item
  \textbf{Cenario medio prazo:} Impacto positivo da modernização das
  redes elétricas rurais no leste do departamento.
\end{itemize}

\subsection{11. Redacao Final}

\subsubsection{Diagnostico Integrado}

General Caballero Álvarez (Puente Kyjhá) é um enclave de alta
produtividade e baixa visibilidade no coração de Canindeyú. Oferece
solos de elite mundial e um ambiente social tranquilo, ``fora do radar''
das grandes crises demográficas metropolitanas. Sua força reside na
abundância de recursos produtivos e na baixa densidade populacional,
garantindo privacidade tática superior. O principal desafio é a
dependência de centros maiores para serviços especializados, o que exige
do ocupante autonomia logística e em saúde. É o local ideal para
estabelecimentos de autossuficiência tecnificada que busquem solo de
alta performance em ambiente de paz social.

\subsubsection{Por que sim}

\begin{itemize}
\tightlist
\item
  Solos de altíssima produtividade e abundância de águas puras.
\item
  Baixíssima densidade populacional (privacidade absoluta).
\item
  Localização discreta e protegida de grandes eixos de conflito.
\item
  Sem restrições da Lei de Fronteira.
\end{itemize}

\subsubsection{Por que nao}

\begin{itemize}
\tightlist
\item
  Infraestrutura de saúde local limitada a atendimentos primários.
\item
  Insegurança tática regional latente no departamento.
\item
  Dependência de estradas internas de terra para acesso às propriedades
  rurais.
\end{itemize}

\subsubsection{Recomendacao Tecnica}

Altamente indicado para estabelecimentos de autossuficiência tecnificada
em solo de alta performance. Requer autonomia em logística básica e
saúde. O GSS de 6.3 reflete a segurança do isolamento em contraste com a
limitação institucional local. Referencia atual de score: GSS 6.3 em
2026-03-05.

\subsection{13.}

\begin{itemize}
\tightlist
\item
  \begin{enumerate}
  \def\labelenumi{\arabic{enumi})}
  \tightlist
  \item
    Dados censitários validados (INE\footnote{Instituto Nacional de Estadística (Instituto Nacional de Estatística), órgão oficial de dados demográficos e censitários.} 2022).
  \end{enumerate}
\item
  \begin{enumerate}
  \def\labelenumi{\arabic{enumi})}
  \setcounter{enumi}{1}
  \tightlist
  \item
    Relatórios de produção agrícola mecanizada departamental (MAG).
  \end{enumerate}
\item
  \begin{enumerate}
  \def\labelenumi{\arabic{enumi})}
  \setcounter{enumi}{2}
  \tightlist
  \item
    Mapa de solos e aptidão agrícola de Canindeyú (Portal
    Geoestadístico).
  \end{enumerate}
\item
  \begin{enumerate}
  \def\labelenumi{\arabic{enumi})}
  \setcounter{enumi}{3}
  \tightlist
  \item
    Registro de segurança pública departamental.
  \end{enumerate}
\end{itemize}

\subsubsection{Combustível}

\textbf{Referência:} PETROPAR /\allowbreak  postos locais (2024) \textbf{Tipo de
localidade:} Interior

\begin{longtable}[]{@{}lll@{}}
\toprule\noalign{}
Combustível & USD/\allowbreak litro & Gs/\allowbreak litro (aprox.) \\
\midrule\noalign{}
\endhead
\bottomrule\noalign{}
\endlastfoot
Gasolina 93 oct & 1.00 & 7,400 \\
Gasolina 97 oct (premium) & 1.10 & 8,140 \\
Diesel & 0.93 & 6,882 \\
\end{longtable}

\begin{quote}
Preços podem variar ±5\% conforme posto e sazonalidade. Chaco e interior
remoto apresentam maior variação.
\end{quote}

\subsubsection{Cobertura Celular}

\textbf{Fonte:} CONATEL PY /\allowbreak  operadoras (2024)

\begin{longtable}[]{@{}ll@{}}
\toprule\noalign{}
Parâmetro & Valor \\
\midrule\noalign{}
\endhead
\bottomrule\noalign{}
\endlastfoot
Cobertura 4G população (dept.) & 78\% \\
Cobertura 4G área rural & 55\% \\
Melhor operadora & Tigo \\
Qualidade rural & moderada \\
\end{longtable}

\begin{quote}
Para áreas rurais fora do núcleo urbano, recomenda-se chip Tigo como
principal e Personal como backup.
\end{quote}

\subsubsection{Internet}

\textbf{Fonte:} CONATEL /\allowbreak  Speedtest Ookla (2024)

\begin{longtable}[]{@{}ll@{}}
\toprule\noalign{}
Parâmetro & Valor \\
\midrule\noalign{}
\endhead
\bottomrule\noalign{}
\endlastfoot
Velocidade média download & 35 Mbps \\
Domicílios com internet (dept.) & 45\% \\
Tecnologia predominante & rádio \\
Opção rural & Starlink disponível (\textasciitilde USD 44/\allowbreak mês) \\
\end{longtable}

\subsubsection{Mercado Imobiliário e Terra
Rural}

\textbf{Fonte:} INDERT /\allowbreak  Clasificados.com.py (2024)

\begin{longtable}[]{@{}ll@{}}
\toprule\noalign{}
Tipo & Referência \\
\midrule\noalign{}
\endhead
\bottomrule\noalign{}
\endlastfoot
Terra agrícola alta prod. (USD/\allowbreak ha) & 7,500 \\
Imóvel urbano (USD/\allowbreak m²) & 650 \\
Aluguel 2 quartos (USD/\allowbreak mês) & 260 \\
\end{longtable}

\begin{quote}
Valores de referência departamental. Localidades menores podem ter
preços 20--40\% abaixo da capital departamental. \#\#\# Saúde
\end{quote}

\textbf{Fonte:} MSPBS /\allowbreak  IPS Paraguay (2024-2026), consolidação
departamental e proxy local

\begin{longtable}[]{@{}ll@{}}
\toprule\noalign{}
Serviço & Disponibilidade \\
\midrule\noalign{}
\endhead
\bottomrule\noalign{}
\endlastfoot
USF /\allowbreak  Posto de Saúde & sim \\
Hospital Regional & sim \\
IPS (seguro social) & não \\
Farmácia & sim \\
Distância ao hospital de referência & \textasciitilde44 km (Salto del
Guaira) \\
\end{longtable}

\textbf{Principais estabelecimentos:} USF local; referencia hospitalar
em Salto del Guaira

\textbf{Observação para imigrantes:} Atendimento primario local ou em
raio curto; casos de maior complexidade seguem para Salto del Guaira.
Cobertura privada continua recomendada para especialidades e urgências
de maior complexidade.
\subsubsection{Indicadores Sociais}

\textbf{Fontes:} PNUD 2020, INE\footnote{Instituto Nacional de Estadística (Instituto Nacional de Estatística), órgão oficial de dados demográficos e censitários.} Censo 2022, INE\footnote{Instituto Nacional de Estadística (Instituto Nacional de Estatística), órgão oficial de dados demográficos e censitários.} EPHC 2023, Ministerio
Público 2024\\
\textbf{Nota:} Valores marcados como \emph{(dept.)} referem-se ao
departamento de Canindeyú; valores \emph{(dist.)} são específicos deste
distrito. Dados marcados \emph{(est.)} são estimativas.

\begin{longtable}[]{@{}
  >{\raggedright\arraybackslash}p{(\linewidth - 4\tabcolsep) * \real{0.4231}}
  >{\raggedright\arraybackslash}p{(\linewidth - 4\tabcolsep) * \real{0.2692}}
  >{\raggedright\arraybackslash}p{(\linewidth - 4\tabcolsep) * \real{0.3077}}@{}}
\toprule\noalign{}
\begin{minipage}[b]{\linewidth}\raggedright
Indicador
\end{minipage} & \begin{minipage}[b]{\linewidth}\raggedright
Valor
\end{minipage} & \begin{minipage}[b]{\linewidth}\raggedright
Âmbito
\end{minipage} \\
\midrule\noalign{}
\endhead
\bottomrule\noalign{}
\endlastfoot
IDH (2020) & 0,685 & dept. \\
Ranking IDH nacional & 15/\allowbreak 18 & dept. \\
Esperança de vida & 71,2 anos & dept. \\
Escolaridade média & 7,3 anos & dept. \\
RNB per capita (USD PPA) & 7.100 & dept. \\
Pobreza monetária (\%) & 35,1\% & dept. \\
Pobreza extrema (\%) & 10,8\% & dept. \\
Índice de Gini & 0,467 & dept. \\
Acesso a água potável (\%) & 67,4\% & dept. \\
Acesso a saneamento (\%) & 64,8\% & dept. \\
Taxa de homicídios (est., /\allowbreak 100k hab) & \textasciitilde15--25
(estimativa, significativamente acima da média nacional) & dept. \\
Índice de segurança & baixo & dept. \\
População (Censo 2022) & 7.842 habitantes & dist. \\
Idade mediana & N/\allowbreak D & dist. \\
Taxa de fecundidade & N/\allowbreak D & dist. \\
\end{longtable}
\subsubsection{Combustível}

\textbf{Referência:} PETROPAR /\allowbreak  postos locais (2024) \textbf{Tipo de
localidade:} Interior

\begin{longtable}[]{@{}lll@{}}
\toprule\noalign{}
Combustível & USD/\allowbreak litro & Gs/\allowbreak litro (aprox.) \\
\midrule\noalign{}
\endhead
\bottomrule\noalign{}
\endlastfoot
Gasolina 93 oct & 1.00 & 7,400 \\
Gasolina 97 oct (premium) & 1.10 & 8,140 \\
Diesel & 0.93 & 6,882 \\
\end{longtable}

\begin{quote}
Preços podem variar ±5\% conforme posto e sazonalidade. Chaco e interior
remoto apresentam maior variação.
\end{quote}

\subsubsection{Cobertura Celular}

\textbf{Fonte:} CONATEL PY /\allowbreak  operadoras (2024)

\begin{longtable}[]{@{}ll@{}}
\toprule\noalign{}
Parâmetro & Valor \\
\midrule\noalign{}
\endhead
\bottomrule\noalign{}
\endlastfoot
Cobertura 4G população (dept.) & 78\% \\
Cobertura 4G área rural & 55\% \\
Melhor operadora & Tigo \\
Qualidade rural & moderada \\
\end{longtable}

\begin{quote}
Para áreas rurais fora do núcleo urbano, recomenda-se chip Tigo como
principal e Personal como backup.
\end{quote}

\subsubsection{Internet}

\textbf{Fonte:} CONATEL /\allowbreak  Speedtest Ookla (2024)

\begin{longtable}[]{@{}ll@{}}
\toprule\noalign{}
Parâmetro & Valor \\
\midrule\noalign{}
\endhead
\bottomrule\noalign{}
\endlastfoot
Velocidade média download & 35 Mbps \\
Domicílios com internet (dept.) & 45\% \\
Tecnologia predominante & rádio \\
Opção rural & Starlink disponível (\textasciitilde USD 44/\allowbreak mês) \\
\end{longtable}

\subsubsection{Mercado Imobiliário e Terra
Rural}

\textbf{Fonte:} INDERT /\allowbreak  Clasificados.com.py (2024)

\begin{longtable}[]{@{}ll@{}}
\toprule\noalign{}
Tipo & Referência \\
\midrule\noalign{}
\endhead
\bottomrule\noalign{}
\endlastfoot
Terra agrícola alta prod. (USD/\allowbreak ha) & 7,500 \\
Imóvel urbano (USD/\allowbreak m²) & 650 \\
Aluguel 2 quartos (USD/\allowbreak mês) & 260 \\
\end{longtable}

\begin{quote}
Valores de referência departamental. Localidades menores podem ter
preços 20--40\% abaixo da capital departamental. \#\#\# Saúde
\end{quote}

\textbf{Fonte:} MSPBS /\allowbreak  IPS Paraguay (2024-2026), consolidação
departamental e proxy local

\begin{longtable}[]{@{}ll@{}}
\toprule\noalign{}
Serviço & Disponibilidade \\
\midrule\noalign{}
\endhead
\bottomrule\noalign{}
\endlastfoot
USF /\allowbreak  Posto de Saúde & sim \\
Hospital Regional & sim \\
IPS (seguro social) & não \\
Farmácia & sim \\
Distância ao hospital de referência & \textasciitilde44 km (Salto del
Guaira) \\
\end{longtable}

\textbf{Principais estabelecimentos:} USF local; referencia hospitalar
em Salto del Guaira

\textbf{Observação para imigrantes:} Atendimento primario local ou em
raio curto; casos de maior complexidade seguem para Salto del Guaira.
Cobertura privada continua recomendada para especialidades e urgências
de maior complexidade.
\newpage
\secaoDiagnostico{Itanara}{\mapaConteudo{-24.0166}{-55.4833}}{Itanará é um dos enclaves mais isolados e discretos da Região Oriental
do Paraguai. Oferece um nível de privacidade demográfica excepcional,
funcionando como um refúgio ``fora do mapa'' para quem busca
invisibilidade estratégica absoluta. Sua força reside no anonimato
geográfico e na abundância de recursos de proteína animal. O custo desse
isolamento é a precariedade total da infraestrutura institucional e a
vulnerabilidade à natureza das chuvas. Sobreviver em Itanará exige um
perfil de ocupante altamente autônomo e técnico, capaz de gerir sua
própria vida sem qualquer dependência do suporte estatal.

\subsubsection{Por que sim}

\begin{itemize}
\tightlist
\item
  Máximo isolamento geográfico e demográfico (privacidade total).
\item
  Ambiente social seguro e estável pelo isolamento.
\item
  Abundância de proteína animal (gado) e recursos hídricos puros.
\item
  Distância máxima de qualquer eixo de conflito ou fallout nacional.
\end{itemize}

\subsubsection{Por que não}

\begin{itemize}
\tightlist
\item
  Risco de isolamento físico total por meses em eventos climáticos.
\item
  Inexistência de serviços de saúde e infraestrutura comercial.
\item
  Dependência de estradas de terra precárias para qualquer logística.
\end{itemize}

\subsubsection{Recomendação
Técnica}

Altamente indicado para estabelecimentos de refúgio tático extremo que
busquem o máximo anonimato. Requer autonomia total em energia, logística
e saúde. O GSS de 6.6 reflete a superioridade do isolamento protegida
pela barreira natural da distância e do clima. Referencia atual de
score: GSS 6.6 em 2026-03-05.

\subsubsection{}

\begin{itemize}
\tightlist
\item
  \begin{enumerate}
  \def\labelenumi{\arabic{enumi})}
  \tightlist
  \item
    Dados censitários validados (INE\footnote{Instituto Nacional de Estadística (Instituto Nacional de Estatística), órgão oficial de dados demográficos e censitários.} 2022).
  \end{enumerate}
\item
  \begin{enumerate}
  \def\labelenumi{\arabic{enumi})}
  \setcounter{enumi}{1}
  \tightlist
  \item
    Relatórios de acessibilidade rural da zona norte de Canindeyú
    (MOPC).
  \end{enumerate}
\item
  \begin{enumerate}
  \def\labelenumi{\arabic{enumi})}
  \setcounter{enumi}{2}
  \tightlist
  \item
    Mapa de solos e bacias hidrográficas de Canindeyú (MAG).
  \end{enumerate}
\item
  \begin{enumerate}
  \def\labelenumi{\arabic{enumi})}
  \setcounter{enumi}{3}
  \tightlist
  \item
    Registro de segurança pública departamental.
  \end{enumerate}
\end{itemize}}
\begin{table}[ht!]
\small \begin{tabular}{p{3.0cm}p{0.8cm}p{6.0cm}} \toprule
\textbf{Critério} & \textbf{Nota} & \textbf{Justificativa} \\ \midrule
A - Ameaças Estratégicas & 7.0 & Isolamento geográfico profundo; excelente invisibilidade tática e ausência de alvos estratégicos \\
B - Risco Social & 8.5 & Densidade populacional ínfima; segurança contra riscos de massa absoluta \\
C - Riscos Naturais & 6.5 & Geologia estável; nota penalizada pela vulnerabilidade a isolamento pluvial severo \\
D - Autossuficiência & 6.5 & Bons recursos de gado e água; limitado pela extrema inacessibilidade a insumos \\
E - Institucional & 4.0 & Ambiente social pacífico por isolamento, mas com suporte institucional quase nulo \\
\midrule \textbf{GSS FINAL} & \textbf{6.6} & \textbf{Moderadamente Seguro (Indicado apenas para projetos de isolamento tático extremo que busquem o máximo anonimato na Região Oriental).} \\ \bottomrule \end{tabular}
\end{table}
\subsubsection{Vulnerabilidades e
Mitigação}\label{vuln-Itanara-vulnerabilidades-e-mitigauxe7uxe3o}

\begin{itemize}
\tightlist
\item
  \textbf{Isolamento Geográfico Severo:} Risco de ficar ``desconectado''
  por terra por semanas (Mitigação: posse obrigatória de veículo 4x4
  robusto e trator; manutenção de estoque de suprimentos para 90 dias).
\item
  \textbf{Fragilidade Elétrica Extrema:} Rede final de linha em zona de
  mata (Mitigação: instalação de sistemas solares fotovoltaicos potentes
  off-grid).
\item
  \textbf{Inexistência de Saúde Local:} Distância crítica de qualquer
  suporte hospitalar (Mitigação: manutenção de estoque médico avançado
  privado e autonomia total em primeiros socorros).
\end{itemize}
\subsubsection{Dossiê de Campo}\label{dossie-Itanara-dossiuxea-de-campo}
\begin{description}[style=nextline,leftmargin=0.5cm,font=\bfseries]
\item[Coordenadas]  24°01'S 55°29'W.
\item[Topographical]  Relevo predominantemente plano a suavemente ondulado, situado no extremo norte do departamento de Canindeyú. Altitude \textasciitilde 220m. Área vasta de \textasciitilde 920 km². Geologicamente estável, risco sísmico desprezível.
\item[Alvos Estratégicos]  Áreas de produção pecuária extensiva de grande escala; Local histórico (Teatro de operações da Guerra da Tríplice Aliança nas proximidades). Ausência total de infraestrutura industrial ou bases militares massivas, oferecendo altíssimo nível de invisibilidade estratégica e proteção por isolamento.
\item[Fallout]  Ventos predominantes do Norte (NE) e Leste. Baixíssimo risco de fallout direto de alvos continentais primários devido ao isolamento profundo.
\item[População]  2.086 habitantes (Censo 2022 Final). O distrito menos povoado do departamento de Canindeyú.
\item[Densidade]  \textasciitilde 2,3 hab/\allowbreak km² (Densidade extremamente baixa, garantindo privacidade tática absoluta e ausência total de riscos sociais de massa urbana).
\item[Vias de Acesso]  Acesso terrestre historicamente difícil via ramais de terra a partir de Curuguaty (\textasciitilde 100 km). Alto nível de isolamento físico durante períodos chuvosos, o que favorece o anonimato estratégico e o desaparecimento "fora do radar".
\item[Serviços]  Infraestrutura básica mínima. Dependência absoluta de Curuguaty ou Salto del Guairá para serviços de saúde especializados, logística comercial e administrativa. A vila urbana é minúscula e focada no essencial.
\item[Custo de Vida]  Muito baixo; economia baseada na pecuária e no autoconsumo rural.
\item[Preço da Terra]  US\$ 1.500-3.000/\allowbreak ha (áreas de campo natural para pecuária). Valores de entrada baixos para a região oriental.
\item[Hidrologia]  Baixo risco de inundações sistêmicas devido à topografia ondulada. Presença de diversos arroios e nascentes perenes de alta qualidade hídrica. Risco severo de isolamento físico por danos em estradas de terra durante eventos de chuva intensa (média 1.700 mm/\allowbreak ano).
\item[Clima]  Subtropical Úmido. Verões intensos; invernos amenos com alta umidade relativa.
\item[Energia]  Rede nacional (Itaipu); altamente instável e sujeita a falhas prolongadas devido ao isolamento da rede rural e dificuldade de manutenção de acesso.
\item[Água]  Abundante via poços artesianos e nascentes naturais; excelente qualidade hídrica subterrânea.
\item[Qualidade do Solo]  Solo franco-argiloso e arenoso; excelente aptidão para pastagens naturais e pecuária de corte. Potencial agrícola limitado pela infraestrutura logística de escoamento.
\item[Resources Locais]  Grande excedente de gado bovino per capita. Elevadíssimo potencial para autossuficiência calórica familiar baseada em proteína animal para a ínfima carga populacional.
\item[Segurança]  Zona rural excepcionalmente pacífica em termos de crime comum. No entanto, exige extrema cautela tática regional devido à proximidade com zonas de fronteira sensíveis e histórico de trânsito de ilícitos no departamento. Estabilidade social baseada no isolamento geográfico.
\item[Leis Local: Município consolidado ("Cidade da Ganaderia"). Livre de Restrição de Fronteira]  Localizado fora da faixa de 50km da fronteira seca internacional.
\end{description}
\paragraph{Solo (SoilGrids 2.0, média ponderada 0--30
cm)}

\begin{longtable}[]{@{}ll@{}}
\toprule\noalign{}
Parâmetro & Valor \\
\midrule\noalign{}
\endhead
\bottomrule\noalign{}
\endlastfoot
pH (H$_2$O) & 5.32 \\
Carbono orgânico (SOC) & 19.63 g/\allowbreak kg \\
Argila & 27.26 \% \\
Areia & 56.5 \% \\
Silte (calc.) & 16.2 \% \\
Densidade aparente & 1.25 g/\allowbreak cm³ \\
\textbf{Aptidão agrícola} & \textbf{Média} \\
\end{longtable}

Fonte: ISRIC SoilGrids 2.0 via WCS (acesso em 2026-03-21). Coords:
-24.0167°, -55.4833°. Média ponderada camadas 0-5, 5-15, 15-30 cm.

\begin{itemize}
\tightlist
\item
  \textbf{Homicidios/\allowbreak 100k:} \textasciitilde18,0 (Regional/\allowbreak Canindeyú -
  Baixo nível local).
\end{itemize}

\subsection{10. Analise de Riscos}

\begin{itemize}
\tightlist
\item
  \textbf{Cenario curto prazo:} Manutenção da tranquilidade rural
  absoluta e isolamento físico sazonal.
\item
  \textbf{Cenario medio prazo:} Valorização das terras como ``bolsão de
  segurança'' tático isolado em departamento produtor.
\end{itemize}

\subsection{11. Redacao Final}

\subsubsection{Diagnostico Integrado}

Itanará é um dos enclaves mais isolados e discretos da Região Oriental
do Paraguai. Oferece um nível de privacidade demográfica excepcional,
funcionando como um refúgio ``fora do mapa'' para quem busca
invisibilidade estratégica absoluta. Sua força reside no anonimato
geográfico e na abundância de recursos de proteína animal. O custo desse
isolamento é a precariedade total da infraestrutura institucional e a
vulnerabilidade à natureza das chuvas. Sobreviver em Itanará exige um
perfil de ocupante altamente autônomo e técnico, capaz de gerir sua
própria vida sem qualquer dependência do suporte estatal.

\subsubsection{Por que sim}

\begin{itemize}
\tightlist
\item
  Máximo isolamento geográfico e demográfico (privacidade total).
\item
  Ambiente social seguro e estável pelo isolamento.
\item
  Abundância de proteína animal (gado) e recursos hídricos puros.
\item
  Distância máxima de qualquer eixo de conflito ou fallout nacional.
\end{itemize}

\subsubsection{Por que nao}

\begin{itemize}
\tightlist
\item
  Risco de isolamento físico total por meses em eventos climáticos.
\item
  Inexistência de serviços de saúde e infraestrutura comercial.
\item
  Dependência de estradas de terra precárias para qualquer logística.
\end{itemize}

\subsubsection{Recomendacao Tecnica}

Altamente indicado para estabelecimentos de refúgio tático extremo que
busquem o máximo anonimato. Requer autonomia total em energia, logística
e saúde. O GSS de 6.6 reflete a superioridade do isolamento protegida
pela barreira natural da distância e do clima. Referencia atual de
score: GSS 6.6 em 2026-03-05.

\subsection{13.}

\begin{itemize}
\tightlist
\item
  \begin{enumerate}
  \def\labelenumi{\arabic{enumi})}
  \tightlist
  \item
    Dados censitários validados (INE\footnote{Instituto Nacional de Estadística (Instituto Nacional de Estatística), órgão oficial de dados demográficos e censitários.} 2022).
  \end{enumerate}
\item
  \begin{enumerate}
  \def\labelenumi{\arabic{enumi})}
  \setcounter{enumi}{1}
  \tightlist
  \item
    Relatórios de acessibilidade rural da zona norte de Canindeyú
    (MOPC).
  \end{enumerate}
\item
  \begin{enumerate}
  \def\labelenumi{\arabic{enumi})}
  \setcounter{enumi}{2}
  \tightlist
  \item
    Mapa de solos e bacias hidrográficas de Canindeyú (MAG).
  \end{enumerate}
\item
  \begin{enumerate}
  \def\labelenumi{\arabic{enumi})}
  \setcounter{enumi}{3}
  \tightlist
  \item
    Registro de segurança pública departamental.
  \end{enumerate}
\end{itemize}

\subsubsection{Combustível}

\textbf{Referência:} PETROPAR /\allowbreak  postos locais (2024) \textbf{Tipo de
localidade:} Interior

\begin{longtable}[]{@{}lll@{}}
\toprule\noalign{}
Combustível & USD/\allowbreak litro & Gs/\allowbreak litro (aprox.) \\
\midrule\noalign{}
\endhead
\bottomrule\noalign{}
\endlastfoot
Gasolina 93 oct & 1.00 & 7,400 \\
Gasolina 97 oct (premium) & 1.10 & 8,140 \\
Diesel & 0.93 & 6,882 \\
\end{longtable}

\begin{quote}
Preços podem variar ±5\% conforme posto e sazonalidade. Chaco e interior
remoto apresentam maior variação.
\end{quote}

\subsubsection{Cobertura Celular}

\textbf{Fonte:} CONATEL PY /\allowbreak  operadoras (2024)

\begin{longtable}[]{@{}ll@{}}
\toprule\noalign{}
Parâmetro & Valor \\
\midrule\noalign{}
\endhead
\bottomrule\noalign{}
\endlastfoot
Cobertura 4G população (dept.) & 78\% \\
Cobertura 4G área rural & 55\% \\
Melhor operadora & Tigo \\
Qualidade rural & moderada \\
\end{longtable}

\begin{quote}
Para áreas rurais fora do núcleo urbano, recomenda-se chip Tigo como
principal e Personal como backup.
\end{quote}

\subsubsection{Internet}

\textbf{Fonte:} CONATEL /\allowbreak  Speedtest Ookla (2024)

\begin{longtable}[]{@{}ll@{}}
\toprule\noalign{}
Parâmetro & Valor \\
\midrule\noalign{}
\endhead
\bottomrule\noalign{}
\endlastfoot
Velocidade média download & 35 Mbps \\
Domicílios com internet (dept.) & 45\% \\
Tecnologia predominante & rádio \\
Opção rural & Starlink disponível (\textasciitilde USD 44/\allowbreak mês) \\
\end{longtable}

\subsubsection{Mercado Imobiliário e Terra
Rural}

\textbf{Fonte:} INDERT /\allowbreak  Clasificados.com.py (2024)

\begin{longtable}[]{@{}ll@{}}
\toprule\noalign{}
Tipo & Referência \\
\midrule\noalign{}
\endhead
\bottomrule\noalign{}
\endlastfoot
Terra agrícola alta prod. (USD/\allowbreak ha) & 7,500 \\
Imóvel urbano (USD/\allowbreak m²) & 650 \\
Aluguel 2 quartos (USD/\allowbreak mês) & 260 \\
\end{longtable}

\begin{quote}
Valores de referência departamental. Localidades menores podem ter
preços 20--40\% abaixo da capital departamental. \#\#\# Saúde
\end{quote}

\textbf{Fonte:} MSPBS /\allowbreak  IPS Paraguay (2024-2026), consolidação
departamental e proxy local

\begin{longtable}[]{@{}ll@{}}
\toprule\noalign{}
Serviço & Disponibilidade \\
\midrule\noalign{}
\endhead
\bottomrule\noalign{}
\endlastfoot
USF /\allowbreak  Posto de Saúde & sim \\
Hospital Regional & não \\
IPS (seguro social) & não \\
Farmácia & sim \\
Distância ao hospital de referência & \textasciitilde150 km (Salto del
Guaira) \\
\end{longtable}

\textbf{Principais estabelecimentos:} USF local; referencia hospitalar
em Salto del Guaira

\textbf{Observação para imigrantes:} Atendimento primario local ou em
raio curto; casos de maior complexidade seguem para Salto del Guaira.
Cobertura privada continua recomendada para especialidades e urgências
de maior complexidade.
\subsubsection{Indicadores Sociais}

\textbf{Fontes:} PNUD 2020, INE\footnote{Instituto Nacional de Estadística (Instituto Nacional de Estatística), órgão oficial de dados demográficos e censitários.} Censo 2022, INE\footnote{Instituto Nacional de Estadística (Instituto Nacional de Estatística), órgão oficial de dados demográficos e censitários.} EPHC 2023, Ministerio
Público 2024\\
\textbf{Nota:} Valores marcados como \emph{(dept.)} referem-se ao
departamento de Canindeyú; valores \emph{(dist.)} são específicos deste
distrito. Dados marcados \emph{(est.)} são estimativas.

\begin{longtable}[]{@{}
  >{\raggedright\arraybackslash}p{(\linewidth - 4\tabcolsep) * \real{0.4231}}
  >{\raggedright\arraybackslash}p{(\linewidth - 4\tabcolsep) * \real{0.2692}}
  >{\raggedright\arraybackslash}p{(\linewidth - 4\tabcolsep) * \real{0.3077}}@{}}
\toprule\noalign{}
\begin{minipage}[b]{\linewidth}\raggedright
Indicador
\end{minipage} & \begin{minipage}[b]{\linewidth}\raggedright
Valor
\end{minipage} & \begin{minipage}[b]{\linewidth}\raggedright
Âmbito
\end{minipage} \\
\midrule\noalign{}
\endhead
\bottomrule\noalign{}
\endlastfoot
IDH (2020) & 0,685 & dept. \\
Ranking IDH nacional & 15/\allowbreak 18 & dept. \\
Esperança de vida & 71,2 anos & dept. \\
Escolaridade média & 7,3 anos & dept. \\
RNB per capita (USD PPA) & 7.100 & dept. \\
Pobreza monetária (\%) & 35,1\% & dept. \\
Pobreza extrema (\%) & 10,8\% & dept. \\
Índice de Gini & 0,467 & dept. \\
Acesso a água potável (\%) & 67,4\% & dept. \\
Acesso a saneamento (\%) & 64,8\% & dept. \\
Taxa de homicídios (est., /\allowbreak 100k hab) & \textasciitilde15--25
(estimativa, significativamente acima da média nacional) & dept. \\
Índice de segurança & baixo & dept. \\
População (Censo 2022) & 2.086 habitantes & dist. \\
Idade mediana & N/\allowbreak D & dist. \\
Taxa de fecundidade & N/\allowbreak D & dist. \\
\end{longtable}
\subsubsection{Combustível}

\textbf{Referência:} PETROPAR /\allowbreak  postos locais (2024) \textbf{Tipo de
localidade:} Interior

\begin{longtable}[]{@{}lll@{}}
\toprule\noalign{}
Combustível & USD/\allowbreak litro & Gs/\allowbreak litro (aprox.) \\
\midrule\noalign{}
\endhead
\bottomrule\noalign{}
\endlastfoot
Gasolina 93 oct & 1.00 & 7,400 \\
Gasolina 97 oct (premium) & 1.10 & 8,140 \\
Diesel & 0.93 & 6,882 \\
\end{longtable}

\begin{quote}
Preços podem variar ±5\% conforme posto e sazonalidade. Chaco e interior
remoto apresentam maior variação.
\end{quote}

\subsubsection{Cobertura Celular}

\textbf{Fonte:} CONATEL PY /\allowbreak  operadoras (2024)

\begin{longtable}[]{@{}ll@{}}
\toprule\noalign{}
Parâmetro & Valor \\
\midrule\noalign{}
\endhead
\bottomrule\noalign{}
\endlastfoot
Cobertura 4G população (dept.) & 78\% \\
Cobertura 4G área rural & 55\% \\
Melhor operadora & Tigo \\
Qualidade rural & moderada \\
\end{longtable}

\begin{quote}
Para áreas rurais fora do núcleo urbano, recomenda-se chip Tigo como
principal e Personal como backup.
\end{quote}

\subsubsection{Internet}

\textbf{Fonte:} CONATEL /\allowbreak  Speedtest Ookla (2024)

\begin{longtable}[]{@{}ll@{}}
\toprule\noalign{}
Parâmetro & Valor \\
\midrule\noalign{}
\endhead
\bottomrule\noalign{}
\endlastfoot
Velocidade média download & 35 Mbps \\
Domicílios com internet (dept.) & 45\% \\
Tecnologia predominante & rádio \\
Opção rural & Starlink disponível (\textasciitilde USD 44/\allowbreak mês) \\
\end{longtable}

\subsubsection{Mercado Imobiliário e Terra
Rural}

\textbf{Fonte:} INDERT /\allowbreak  Clasificados.com.py (2024)

\begin{longtable}[]{@{}ll@{}}
\toprule\noalign{}
Tipo & Referência \\
\midrule\noalign{}
\endhead
\bottomrule\noalign{}
\endlastfoot
Terra agrícola alta prod. (USD/\allowbreak ha) & 7,500 \\
Imóvel urbano (USD/\allowbreak m²) & 650 \\
Aluguel 2 quartos (USD/\allowbreak mês) & 260 \\
\end{longtable}

\begin{quote}
Valores de referência departamental. Localidades menores podem ter
preços 20--40\% abaixo da capital departamental. \#\#\# Saúde
\end{quote}

\textbf{Fonte:} MSPBS /\allowbreak  IPS Paraguay (2024-2026), consolidação
departamental e proxy local

\begin{longtable}[]{@{}ll@{}}
\toprule\noalign{}
Serviço & Disponibilidade \\
\midrule\noalign{}
\endhead
\bottomrule\noalign{}
\endlastfoot
USF /\allowbreak  Posto de Saúde & sim \\
Hospital Regional & não \\
IPS (seguro social) & não \\
Farmácia & sim \\
Distância ao hospital de referência & \textasciitilde150 km (Salto del
Guaira) \\
\end{longtable}

\textbf{Principais estabelecimentos:} USF local; referencia hospitalar
em Salto del Guaira

\textbf{Observação para imigrantes:} Atendimento primario local ou em
raio curto; casos de maior complexidade seguem para Salto del Guaira.
Cobertura privada continua recomendada para especialidades e urgências
de maior complexidade.
\newpage
\secaoDiagnostico{Katuete}{\mapaConteudo{-24.2500}{-54.7500}}{Katueté é o símbolo da eficiência agroindustrial em Canindeyú. Oferece
solos de elite mundial e uma infraestrutura de suporte (saúde/\allowbreak comércio)
que está entre as melhores do interior do país. Sua força absoluta
reside na combinação de riqueza de recursos com uma densidade
demográfica baixa que garante privacidade e segurança. É o local ideal
para estabelecimentos de autossuficiência tecnificada ou bases
operacionais, exigindo do ocupante um posicionamento rural periférico
para mitigar a visibilidade estratégica do polo de silos.

\subsubsection{Por que sim}

\begin{itemize}
\tightlist
\item
  Solos de altíssima produtividade e abundância de recursos hídricos.
\item
  Ambiente social seguro e infraestrutura de serviços privados de alta
  qualidade.
\item
  Isolamento demográfico rural superior.
\item
  Sem restrições da Lei de Fronteira.
\end{itemize}

\subsubsection{Por que não}

\begin{itemize}
\tightlist
\item
  Alvo estratégico logístico e alimentar em cenários de conflito.
\item
  Valor da terra entre os mais altos do país.
\item
  Dependência de vias asfálticas principais para escoamento.
\end{itemize}

\subsubsection{Recomendação
Técnica}

Altamente indicado para investimentos agroindustriais ou residências
rurais permanentes de alto padrão. Requer plano de autonomia hídrica e
energética individual. O GSS de 7.5 reflete a excepcional solidez dos
recursos e da governança local em equilíbrio com a visibilidade
estratégica. Referencia atual de score: GSS 7.5 em 2026-03-05.

\subsubsection{}

\begin{itemize}
\tightlist
\item
  \begin{enumerate}
  \def\labelenumi{\arabic{enumi})}
  \tightlist
  \item
    Dados censitários validados (INE\footnote{Instituto Nacional de Estadística (Instituto Nacional de Estatística), órgão oficial de dados demográficos e censitários.} 2022).
  \end{enumerate}
\item
  \begin{enumerate}
  \def\labelenumi{\arabic{enumi})}
  \setcounter{enumi}{1}
  \tightlist
  \item
    Relatórios de produção agrícola mecanizada departamental (MAG).
  \end{enumerate}
\item
  \begin{enumerate}
  \def\labelenumi{\arabic{enumi})}
  \setcounter{enumi}{2}
  \tightlist
  \item
    Mapa de solos e aptidão agrícola de Canindeyú (Portal
    Geoestadístico).
  \end{enumerate}
\item
  \begin{enumerate}
  \def\labelenumi{\arabic{enumi})}
  \setcounter{enumi}{3}
  \tightlist
  \item
    Registro de segurança pública e governança municipal próspera.
  \end{enumerate}
\end{itemize}}
\begin{table}[ht!]
\small \begin{tabular}{p{3.0cm}p{0.8cm}p{6.0cm}} \toprule
\textbf{Critério} & \textbf{Nota} & \textbf{Justificativa} \\ \midrule
A - Ameaças Estratégicas & 6.5 & Hub agroindustrial estratégico; boa invisibilidade tática fora do eixo urbano, visibilidade moderada como alvo logístico \\
B - Risco Social & 8.0 & Densidade populacional baixa; isolamento social superior em ambiente rural organizado \\
C - Riscos Naturais & 7.0 & Geologia muito estável e baixo risco de grandes desastres naturais \\
D - Autossuficiência & 8.5 & Recursos agrícolas e hídricos excepcionais; solo de elite e silos modernos \\
E - Institucional & 7.5 & Forte segurança institucional e serviços comerciais de primeiro nível regional \\
\midrule \textbf{GSS FINAL} & \textbf{7.5} & \textbf{Seguro (Localidade de elite para suporte logístico e produção de recursos, com excelente isolamento demográfico rural).} \\ \bottomrule \end{tabular}
\end{table}
\subsubsection{Vulnerabilidades e
Mitigação}\label{vuln-Katuete-vulnerabilidades-e-mitigauxe7uxe3o}

\begin{itemize}
\tightlist
\item
  \textbf{Visibilidade Tática:} Polo agroindustrial é um alvo
  estratégico para controle de suprimentos (Mitigação: residência rural
  afastada 10-15 km do eixo urbano, mantendo autonomia total).
\item
  \textbf{Dependência de Exportação:} Economia vinculada ao mercado
  global de commodities (Mitigação: diversificação produtiva para
  consumo interno e energia solar redundante).
\item
  \textbf{Gargalo Rodoviário:} Dependência da Ruta PY03 para escoamento
  (Mitigação: mapeamento de rotas rurais vicinais em direção a Canindeyú
  interior).
\end{itemize}
\subsubsection{Dossiê de Campo}\label{dossie-Katuete-dossiuxea-de-campo}
\begin{description}[style=nextline,leftmargin=0.5cm,font=\bfseries]
\item[Coordenadas]  24°15'S 54°45'W.
\item[Topografia]  Relevo predominantemente plano a suavemente ondulado, com solos de elite (Latossolo Vermelho). Altitude \textasciitilde 250m. Área de \textasciitilde 812 km². Geologicamente muito estável, risco sísmico desprezível.
\item[Alvos Estratégicos]  Polo Agroindustrial Regional (Grande concentração de silos e moinhos); Nó logístico secundário conectando a Ruta PY03 ao interior do departamento; Subestação ANDE\footnote{Administración Nacional de Electricidade (Administração Nacional de Eletricidade), companhia estatal de energia.}. Ausência de alvos militares diretos, mas centro de gravidade para o controle de suprimentos agrícolas no leste do país.
\item[Fallout]  Ventos predominantes do Norte (NE) e Leste. Baixo risco de fallout direto; posição isolada em relação aos grandes centros de comando.
\item[População]  10.774 habitantes (Censo 2022 Final). População cosmopolita com forte herança de imigrantes brasileiros, altamente integrada ao agronegócio mecanizado.
\item[Densidade]  \textasciitilde 13,3 hab/\allowbreak km² (Densidade baixa, oferecendo excelente isolamento tático e privacidade fora do núcleo urbano consolidado).
\item[Vias de Acesso]  Localização privilegiada sobre a Ruta PY03. Conectividade asfáltica de elite para Salto del Guairá e Assunção. Facilidade logística tática em tempos de normalidade, mas sujeita a controle militar de eixos rodoviários em crises.
\item[Serviços]  Infraestrutura de serviços de primeiro nível regional. Sanatórios privados modernos, rede bancária completa e suporte técnico agroindustrial avançado.
\item[Custo de Vida]  Médio; economia dolarizada pelo agronegócio, com alta valorização imobiliária urbana.
\item[Preço da Terra]  US\$ 10.000-15.000/\allowbreak ha (terras mecanizadas de elite - "Terra Roxa").
\item[Hidrologia]  Baixo risco de inundações sistêmicas devido à topografia favorável e gestão hídrica eficiente. Presença de arroios perenes utilizados para fins agroindustriais.
\item[Clima]  Subtropical Úmido. Verões quentes e chuvosos; invernos amenos com geadas ocasionais benéficas.
\item[Energia]  Rede nacional (Itaipu); boa estabilidade devido à importância industrial. Nodo de distribuição regional.
\item[Água]  Abundante via poços artesianos profundos e sistemas de junta de saneamento modernos; excelente qualidade hídrica subterrânea.
\item[Qualidade do Solo]  Latossolo Vermelho de Elite; solo profundo e de altíssima fertilidade mundial. Excelente aptidão para soja, milho e trigo.
\item[Resources Locais]  Grande polo produtor de grãos e proteína animal. Elevadíssimo potencial para autossuficiência alimentar básica e suporte logístico em escala regional.
\item[Segurança]  Cidade próspera com segurança social superior à média nacional. Criminalidade urbana monitorada. Forte coesão comunitária baseada na cultura do trabalho e agronegócio. Presença de segurança patrimonial privada robusta nas indústrias.
\item[Leis Local]  Município consolidado e polo de influência econômica ("A Cidade que mais cresce"). \\ \textbf{Livre de Restrição} de Fronteira. \\ \textit{Aquisição de terra rural proibida para estrangeiros limítrofes.}
\end{description}
\paragraph{Solo (SoilGrids 2.0, média ponderada 0--30
cm)}

\begin{longtable}[]{@{}ll@{}}
\toprule\noalign{}
Parâmetro & Valor \\
\midrule\noalign{}
\endhead
\bottomrule\noalign{}
\endlastfoot
pH (H$_2$O) & 5.32 \\
Carbono orgânico (SOC) & 19.22 g/\allowbreak kg \\
Argila & 37.53 \% \\
Areia & 43.27 \% \\
Silte (calc.) & 19.2 \% \\
Densidade aparente & 1.26 g/\allowbreak cm³ \\
\textbf{Aptidão agrícola} & \textbf{Média} \\
\end{longtable}

Fonte: ISRIC SoilGrids 2.0 via WCS (acesso em 2026-03-21). Coords:
-24.25°, -54.75°. Média ponderada camadas 0-5, 5-15, 15-30 cm.

\begin{itemize}
\tightlist
\item
  \textbf{Homicidios/\allowbreak 100k:} \textasciitilde12,0 (Regional/\allowbreak Abaixo da
  média de fronteira).
\end{itemize}

\subsection{10. Analise de Riscos}

\begin{itemize}
\tightlist
\item
  \textbf{Cenario curto prazo:} Manutenção da prosperidade econômica e
  valorização imobiliária agroindustrial.
\item
  \textbf{Cenario medio prazo:} Impacto de mudanças nos fluxos
  rodoviários do Corredor Bioceânico na logística local.
\end{itemize}

\subsection{11. Redacao Final}

\subsubsection{Diagnostico Integrado}

Katueté é o símbolo da eficiência agroindustrial em Canindeyú. Oferece
solos de elite mundial e uma infraestrutura de suporte (saúde/\allowbreak comércio)
que está entre as melhores do interior do país. Sua força absoluta
reside na combinação de riqueza de recursos com uma densidade
demográfica baixa que garante privacidade e segurança. É o local ideal
para estabelecimentos de autossuficiência tecnificada ou bases
operacionais, exigindo do ocupante um posicionamento rural periférico
para mitigar a visibilidade estratégica do polo de silos.

\subsubsection{Por que sim}

\begin{itemize}
\tightlist
\item
  Solos de altíssima produtividade e abundância de recursos hídricos.
\item
  Ambiente social seguro e infraestrutura de serviços privados de alta
  qualidade.
\item
  Isolamento demográfico rural superior.
\item
  Sem restrições da Lei de Fronteira.
\end{itemize}

\subsubsection{Por que nao}

\begin{itemize}
\tightlist
\item
  Alvo estratégico logístico e alimentar em cenários de conflito.
\item
  Valor da terra entre os mais altos do país.
\item
  Dependência de vias asfálticas principais para escoamento.
\end{itemize}

\subsubsection{Recomendacao Tecnica}

Altamente indicado para investimentos agroindustriais ou residências
rurais permanentes de alto padrão. Requer plano de autonomia hídrica e
energética individual. O GSS de 7.5 reflete a excepcional solidez dos
recursos e da governança local em equilíbrio com a visibilidade
estratégica. Referencia atual de score: GSS 7.5 em 2026-03-05.

\subsection{13.}

\begin{itemize}
\tightlist
\item
  \begin{enumerate}
  \def\labelenumi{\arabic{enumi})}
  \tightlist
  \item
    Dados censitários validados (INE\footnote{Instituto Nacional de Estadística (Instituto Nacional de Estatística), órgão oficial de dados demográficos e censitários.} 2022).
  \end{enumerate}
\item
  \begin{enumerate}
  \def\labelenumi{\arabic{enumi})}
  \setcounter{enumi}{1}
  \tightlist
  \item
    Relatórios de produção agrícola mecanizada departamental (MAG).
  \end{enumerate}
\item
  \begin{enumerate}
  \def\labelenumi{\arabic{enumi})}
  \setcounter{enumi}{2}
  \tightlist
  \item
    Mapa de solos e aptidão agrícola de Canindeyú (Portal
    Geoestadístico).
  \end{enumerate}
\item
  \begin{enumerate}
  \def\labelenumi{\arabic{enumi})}
  \setcounter{enumi}{3}
  \tightlist
  \item
    Registro de segurança pública e governança municipal próspera.
  \end{enumerate}
\end{itemize}

\subsubsection{Combustível}

\textbf{Referência:} PETROPAR /\allowbreak  postos locais (2024) \textbf{Tipo de
localidade:} Interior

\begin{longtable}[]{@{}lll@{}}
\toprule\noalign{}
Combustível & USD/\allowbreak litro & Gs/\allowbreak litro (aprox.) \\
\midrule\noalign{}
\endhead
\bottomrule\noalign{}
\endlastfoot
Gasolina 93 oct & 1.00 & 7,400 \\
Gasolina 97 oct (premium) & 1.10 & 8,140 \\
Diesel & 0.93 & 6,882 \\
\end{longtable}

\begin{quote}
Preços podem variar ±5\% conforme posto e sazonalidade. Chaco e interior
remoto apresentam maior variação.
\end{quote}

\subsubsection{Cobertura Celular}

\textbf{Fonte:} CONATEL PY /\allowbreak  operadoras (2024)

\begin{longtable}[]{@{}ll@{}}
\toprule\noalign{}
Parâmetro & Valor \\
\midrule\noalign{}
\endhead
\bottomrule\noalign{}
\endlastfoot
Cobertura 4G população (dept.) & 78\% \\
Cobertura 4G área rural & 55\% \\
Melhor operadora & Tigo \\
Qualidade rural & moderada \\
\end{longtable}

\begin{quote}
Para áreas rurais fora do núcleo urbano, recomenda-se chip Tigo como
principal e Personal como backup.
\end{quote}

\subsubsection{Internet}

\textbf{Fonte:} CONATEL /\allowbreak  Speedtest Ookla (2024)

\begin{longtable}[]{@{}ll@{}}
\toprule\noalign{}
Parâmetro & Valor \\
\midrule\noalign{}
\endhead
\bottomrule\noalign{}
\endlastfoot
Velocidade média download & 35 Mbps \\
Domicílios com internet (dept.) & 45\% \\
Tecnologia predominante & rádio \\
Opção rural & Starlink disponível (\textasciitilde USD 44/\allowbreak mês) \\
\end{longtable}

\subsubsection{Mercado Imobiliário e Terra
Rural}

\textbf{Fonte:} INDERT /\allowbreak  Clasificados.com.py (2024)

\begin{longtable}[]{@{}ll@{}}
\toprule\noalign{}
Tipo & Referência \\
\midrule\noalign{}
\endhead
\bottomrule\noalign{}
\endlastfoot
Terra agrícola alta prod. (USD/\allowbreak ha) & 7,500 \\
Imóvel urbano (USD/\allowbreak m²) & 650 \\
Aluguel 2 quartos (USD/\allowbreak mês) & 260 \\
\end{longtable}

\begin{quote}
Valores de referência departamental. Localidades menores podem ter
preços 20--40\% abaixo da capital departamental. \#\#\# Saúde
\end{quote}

\textbf{Fonte:} MSPBS /\allowbreak  IPS Paraguay (2024-2026), consolidação
departamental e proxy local

\begin{longtable}[]{@{}ll@{}}
\toprule\noalign{}
Serviço & Disponibilidade \\
\midrule\noalign{}
\endhead
\bottomrule\noalign{}
\endlastfoot
USF /\allowbreak  Posto de Saúde & sim \\
Hospital Regional & não \\
IPS (seguro social) & não \\
Farmácia & sim \\
Distância ao hospital de referência & \textasciitilde63 km (Salto del
Guaira) \\
\end{longtable}

\textbf{Principais estabelecimentos:} USF local; referencia hospitalar
em Salto del Guaira

\textbf{Observação para imigrantes:} Atendimento primario local ou em
raio curto; casos de maior complexidade seguem para Salto del Guaira.
Cobertura privada continua recomendada para especialidades e urgências
de maior complexidade.
\subsubsection{Indicadores Sociais}

\textbf{Fontes:} PNUD 2020, INE\footnote{Instituto Nacional de Estadística (Instituto Nacional de Estatística), órgão oficial de dados demográficos e censitários.} Censo 2022, INE\footnote{Instituto Nacional de Estadística (Instituto Nacional de Estatística), órgão oficial de dados demográficos e censitários.} EPHC 2023, Ministerio
Público 2024\\
\textbf{Nota:} Valores marcados como \emph{(dept.)} referem-se ao
departamento de Canindeyú; valores \emph{(dist.)} são específicos deste
distrito. Dados marcados \emph{(est.)} são estimativas.

\begin{longtable}[]{@{}
  >{\raggedright\arraybackslash}p{(\linewidth - 4\tabcolsep) * \real{0.4231}}
  >{\raggedright\arraybackslash}p{(\linewidth - 4\tabcolsep) * \real{0.2692}}
  >{\raggedright\arraybackslash}p{(\linewidth - 4\tabcolsep) * \real{0.3077}}@{}}
\toprule\noalign{}
\begin{minipage}[b]{\linewidth}\raggedright
Indicador
\end{minipage} & \begin{minipage}[b]{\linewidth}\raggedright
Valor
\end{minipage} & \begin{minipage}[b]{\linewidth}\raggedright
Âmbito
\end{minipage} \\
\midrule\noalign{}
\endhead
\bottomrule\noalign{}
\endlastfoot
IDH (2020) & 0,685 & dept. \\
Ranking IDH nacional & 15/\allowbreak 18 & dept. \\
Esperança de vida & 71,2 anos & dept. \\
Escolaridade média & 7,3 anos & dept. \\
RNB per capita (USD PPA) & 7.100 & dept. \\
Pobreza monetária (\%) & 35,1\% & dept. \\
Pobreza extrema (\%) & 10,8\% & dept. \\
Índice de Gini & 0,467 & dept. \\
Acesso a água potável (\%) & 67,4\% & dept. \\
Acesso a saneamento (\%) & 64,8\% & dept. \\
Taxa de homicídios (est., /\allowbreak 100k hab) & \textasciitilde15--25
(estimativa, significativamente acima da média nacional) & dept. \\
Índice de segurança & baixo & dept. \\
População (Censo 2022) & 10.774 habitantes & dist. \\
Idade mediana & N/\allowbreak D & dist. \\
Taxa de fecundidade & N/\allowbreak D & dist. \\
\end{longtable}
\subsubsection{Combustível}

\textbf{Referência:} PETROPAR /\allowbreak  postos locais (2024) \textbf{Tipo de
localidade:} Interior

\begin{longtable}[]{@{}lll@{}}
\toprule\noalign{}
Combustível & USD/\allowbreak litro & Gs/\allowbreak litro (aprox.) \\
\midrule\noalign{}
\endhead
\bottomrule\noalign{}
\endlastfoot
Gasolina 93 oct & 1.00 & 7,400 \\
Gasolina 97 oct (premium) & 1.10 & 8,140 \\
Diesel & 0.93 & 6,882 \\
\end{longtable}

\begin{quote}
Preços podem variar ±5\% conforme posto e sazonalidade. Chaco e interior
remoto apresentam maior variação.
\end{quote}

\subsubsection{Cobertura Celular}

\textbf{Fonte:} CONATEL PY /\allowbreak  operadoras (2024)

\begin{longtable}[]{@{}ll@{}}
\toprule\noalign{}
Parâmetro & Valor \\
\midrule\noalign{}
\endhead
\bottomrule\noalign{}
\endlastfoot
Cobertura 4G população (dept.) & 78\% \\
Cobertura 4G área rural & 55\% \\
Melhor operadora & Tigo \\
Qualidade rural & moderada \\
\end{longtable}

\begin{quote}
Para áreas rurais fora do núcleo urbano, recomenda-se chip Tigo como
principal e Personal como backup.
\end{quote}

\subsubsection{Internet}

\textbf{Fonte:} CONATEL /\allowbreak  Speedtest Ookla (2024)

\begin{longtable}[]{@{}ll@{}}
\toprule\noalign{}
Parâmetro & Valor \\
\midrule\noalign{}
\endhead
\bottomrule\noalign{}
\endlastfoot
Velocidade média download & 35 Mbps \\
Domicílios com internet (dept.) & 45\% \\
Tecnologia predominante & rádio \\
Opção rural & Starlink disponível (\textasciitilde USD 44/\allowbreak mês) \\
\end{longtable}

\subsubsection{Mercado Imobiliário e Terra
Rural}

\textbf{Fonte:} INDERT /\allowbreak  Clasificados.com.py (2024)

\begin{longtable}[]{@{}ll@{}}
\toprule\noalign{}
Tipo & Referência \\
\midrule\noalign{}
\endhead
\bottomrule\noalign{}
\endlastfoot
Terra agrícola alta prod. (USD/\allowbreak ha) & 7,500 \\
Imóvel urbano (USD/\allowbreak m²) & 650 \\
Aluguel 2 quartos (USD/\allowbreak mês) & 260 \\
\end{longtable}

\begin{quote}
Valores de referência departamental. Localidades menores podem ter
preços 20--40\% abaixo da capital departamental. \#\#\# Saúde
\end{quote}

\textbf{Fonte:} MSPBS /\allowbreak  IPS Paraguay (2024-2026), consolidação
departamental e proxy local

\begin{longtable}[]{@{}ll@{}}
\toprule\noalign{}
Serviço & Disponibilidade \\
\midrule\noalign{}
\endhead
\bottomrule\noalign{}
\endlastfoot
USF /\allowbreak  Posto de Saúde & sim \\
Hospital Regional & não \\
IPS (seguro social) & não \\
Farmácia & sim \\
Distância ao hospital de referência & \textasciitilde63 km (Salto del
Guaira) \\
\end{longtable}

\textbf{Principais estabelecimentos:} USF local; referencia hospitalar
em Salto del Guaira

\textbf{Observação para imigrantes:} Atendimento primario local ou em
raio curto; casos de maior complexidade seguem para Salto del Guaira.
Cobertura privada continua recomendada para especialidades e urgências
de maior complexidade.
\newpage
\secaoDiagnostico{La Paloma}{\mapaConteudo{-24.1333}{-54.6166}}{La Paloma é um dos pilares produtivos e mais estáveis do leste
paraguaio. Oferece solos de produtividade mundial e um isolamento
demográfico que garante privacidade superior, protegida pela organização
da comunidade produtora. Sua força reside na excepcional abundância de
recursos hídricos e na eficiência da produção de sementes e grãos. É o
local ideal para quem busca autonomia tática em solo fértil,
transformando a riqueza de recursos em uma barreira de resiliência
alimentar e econômica. O contexto exige apenas autonomia em saúde de
alta complexidade.

\subsubsection{Por que sim}

\begin{itemize}
\tightlist
\item
  Solos de altíssima produtividade e abundância de águas puras.
\item
  Baixíssima densidade populacional (privacidade absoluta).
\item
  Conectividade asfáltica de elite via Ruta PY03.
\item
  Sem restrições da Lei de Fronteira.
\end{itemize}

\subsubsection{Por que não}

\begin{itemize}
\tightlist
\item
  Alvo estratégico logístico e de suprimentos em cenários de conflito.
\item
  Dependência de centros maiores para serviços de saúde avançados.
\item
  Valor da terra entre os mais altos do departamento.
\end{itemize}

\subsubsection{Recomendação
Técnica}

Altamente indicado para estabelecimentos de autossuficiência tecnificada
que busquem solo de alta performance. Requer autonomia energética e
médica complementar. O GSS de 7.2 reflete a solidez dos recursos e da
segurança institucional da localidade em equilíbrio com a visibilidade
estratégica. Referencia atual de score: GSS 7.2 em 2026-03-05.

\subsubsection{}

\begin{itemize}
\tightlist
\item
  \begin{enumerate}
  \def\labelenumi{\arabic{enumi})}
  \tightlist
  \item
    Dados censitários validados (INE\footnote{Instituto Nacional de Estadística (Instituto Nacional de Estatística), órgão oficial de dados demográficos e censitários.} 2022).
  \end{enumerate}
\item
  \begin{enumerate}
  \def\labelenumi{\arabic{enumi})}
  \setcounter{enumi}{1}
  \tightlist
  \item
    Relatórios de produção agrícola e certificação de sementes (MAG).
  \end{enumerate}
\item
  \begin{enumerate}
  \def\labelenumi{\arabic{enumi})}
  \setcounter{enumi}{2}
  \tightlist
  \item
    Mapa de solos e bacias hidrográficas de Canindeyú (MAG).
  \end{enumerate}
\item
  \begin{enumerate}
  \def\labelenumi{\arabic{enumi})}
  \setcounter{enumi}{3}
  \tightlist
  \item
    Registro de segurança pública departamental.
  \end{enumerate}
\end{itemize}}
\begin{table}[ht!]
\small \begin{tabular}{p{3.0cm}p{0.8cm}p{6.0cm}} \toprule
\textbf{Critério} & \textbf{Nota} & \textbf{Justificativa} \\ \midrule
A - Ameaças Estratégicas & 6.5 & Hub agroindustrial estratégico; boa invisibilidade tática fora do núcleo urbano, visibilidade moderada como alvo logístico \\
B - Risco Social & 8.0 & Densidade populacional baixa; isolamento social satisfatório em ambiente rural organizado \\
C - Riscos Naturais & 7.0 & Geologia muito estável e baixo risco de grandes desastres naturais \\
D - Autossuficiência & 8.0 & Excelentes recursos agrícolas e hídricos; solo de elite e polo de sementes \\
E - Institucional & 7.0 & Forte segurança institucional privada e serviços funcionais próximos \\
\midrule \textbf{GSS FINAL} & \textbf{7.2} & \textbf{Seguro (Altamente recomendado para quem busca isolamento tático rural em terras férteis com suporte industrial próximo).} \\ \bottomrule \end{tabular}
\end{table}
\subsubsection{Vulnerabilidades e
Mitigação}\label{vuln-La_Paloma-vulnerabilidades-e-mitigauxe7uxe3o}

\begin{itemize}
\tightlist
\item
  \textbf{Visibilidade Tática:} Polo de silos é um alvo estratégico em
  crises de desabastecimento (Mitigação: residência rural afastada 5-10
  km do núcleo industrial).
\item
  \textbf{Dependência de Exportação:} Economia vinculada ao mercado de
  commodities (Mitigação: diversificação produtiva para consumo interno
  e autonomia energética).
\item
  \textbf{Gargalo Logístico:} Dependência da Ruta PY03 (Mitigação:
  mapeamento de rotas rurais vicinais em direção a Katueté).
\end{itemize}
\subsubsection{Dossiê de Campo}\label{dossie-La_Paloma-dossiuxea-de-campo}
\begin{description}[style=nextline,leftmargin=0.5cm,font=\bfseries]
\item[Coordenadas]  24°08'S 54°37'W.
\item[Topografia]  Relevo predominantemente plano a suavemente ondulado, com solos de altíssima produtividade (Terra Roxa). Altitude \textasciitilde 250m. Área de \textasciitilde 728 km². Geologicamente muito estável, risco sísmico desprezível.
\item[Alvos Estratégicos]  Polo Agroindustrial Regional (Sede da Cooperativa La Paloma e grandes silos); Nó logístico sobre a Ruta PY03; Proximidade tática com o entroncamento para Katueté e Salto del Guairá. Ausência de alvos militares diretos, mas estratégico para o fluxo de commodities nacionais.
\item[Fallout]  Ventos predominantes do Norte (NE) e Leste. Baixíssimo risco de fallout direto; posição isolada e protegida.
\item[População]  7.284 habitantes (Censo 2022 Final). População estável com forte cultura de agronegócio e comércio de sementes/\allowbreak insumos.
\item[Densidade]  \textasciitilde 10,0 hab/\allowbreak km² (Densidade baixa, oferecendo isolamento tático superior fora do núcleo comercial asfáltico).
\item[Vias de Acesso]  Localização estratégica sobre a Ruta PY03. Conectividade asfáltica de elite para Assunção e Salto del Guairá. Facilidade logística tática em tempos de paz, mas sujeita a controle de eixos rodoviários em crises nacionais.
\item[Serviços]  Infraestrutura básica a média funcional. USF local e centros de apoio técnico agroindustrial. Dependência de Katueté ou Salto del Guairá para serviços de saúde de alta complexidade.
\item[Custo de Vida]  Baixo; economia impulsionada pelo agronegócio e produção de grãos.
\item[Preço da Terra]  US\$ 10.000-14.000/\allowbreak ha (terras mecanizadas de elite - "Terra Roxa").
\item[Hidrologia]  Baixo risco de inundações sistêmicas devido à topografia favorável e drenagem natural eficiente. Presença de arroios perenes de alta qualidade hídrica utilizados para pecuária e irrigação.
\item[Clima]  Subtropical Úmido. Verões intensos e invernos amenos com geadas ocasionais benéficas ao trigo.
\item[Energia]  Rede nacional (Itaipu); boa estabilidade devido à importância das agroindústrias locais.
\item[Água]  Abundante via poços artesianos profundos e nascentes naturais; excelente qualidade hídrica subterrânea.
\item[Qualidade do Solo]  Latossolo Vermelho de Elite; solo profundo e de altíssima fertilidade. Excelente aptidão para soja, milho e sementes selecionadas.
\item[Resources Locais]  Grande produção de grãos e proteína animal. Elevadíssimo potencial para autossuficiência alimentar básica em nível de propriedade devido à vasta área produtiva e baixa população rural.
\item[Segurança]  Cidade próspera e pacífica. Baixíssima criminalidade urbana comum. Forte coesão comunitária baseada na cultura produtora. Presença de segurança patrimonial privada eficiente nas zonas industriais.
\item[Leis Local]  Município consolidado e polo de sementes regionais. \\ \textbf{Livre de Restrição} de Fronteira. \\ \textit{Aquisição de terra rural proibida para estrangeiros limítrofes.}
\end{description}
\paragraph{Solo (SoilGrids 2.0, média ponderada 0--30
cm)}

\begin{longtable}[]{@{}ll@{}}
\toprule\noalign{}
Parâmetro & Valor \\
\midrule\noalign{}
\endhead
\bottomrule\noalign{}
\endlastfoot
pH (H$_2$O) & 5.32 \\
Carbono orgânico (SOC) & 18.55 g/\allowbreak kg \\
Argila & 28.58 \% \\
Areia & 56.45 \% \\
Silte (calc.) & 15.0 \% \\
Densidade aparente & 1.28 g/\allowbreak cm³ \\
\textbf{Aptidão agrícola} & \textbf{Média} \\
\end{longtable}

Fonte: ISRIC SoilGrids 2.0 via WCS (acesso em 2026-03-21). Coords:
-24.1333°, -54.6167°. Média ponderada camadas 0-5, 5-15, 15-30 cm.

\begin{itemize}
\tightlist
\item
  \textbf{Homicidios/\allowbreak 100k:} \textasciitilde12,0 (Regional/\allowbreak Abaixo da
  média de fronteira).
\end{itemize}

\subsection{10. Analise de Riscos}

\begin{itemize}
\tightlist
\item
  \textbf{Cenario curto prazo:} Manutenção da prosperidade econômica e
  controle de segurança regional.
\item
  \textbf{Cenario medio prazo:} Impacto de mudanças nos fluxos
  rodoviários da fronteira na logística local.
\end{itemize}

\subsection{11. Redacao Final}

\subsubsection{Diagnostico Integrado}

La Paloma é um dos pilares produtivos e mais estáveis do leste
paraguaio. Oferece solos de produtividade mundial e um isolamento
demográfico que garante privacidade superior, protegida pela organização
da comunidade produtora. Sua força reside na excepcional abundância de
recursos hídricos e na eficiência da produção de sementes e grãos. É o
local ideal para quem busca autonomia tática em solo fértil,
transformando a riqueza de recursos em uma barreira de resiliência
alimentar e econômica. O contexto exige apenas autonomia em saúde de
alta complexidade.

\subsubsection{Por que sim}

\begin{itemize}
\tightlist
\item
  Solos de altíssima produtividade e abundância de águas puras.
\item
  Baixíssima densidade populacional (privacidade absoluta).
\item
  Conectividade asfáltica de elite via Ruta PY03.
\item
  Sem restrições da Lei de Fronteira.
\end{itemize}

\subsubsection{Por que nao}

\begin{itemize}
\tightlist
\item
  Alvo estratégico logístico e de suprimentos em cenários de conflito.
\item
  Dependência de centros maiores para serviços de saúde avançados.
\item
  Valor da terra entre os mais altos do departamento.
\end{itemize}

\subsubsection{Recomendacao Tecnica}

Altamente indicado para estabelecimentos de autossuficiência tecnificada
que busquem solo de alta performance. Requer autonomia energética e
médica complementar. O GSS de 7.2 reflete a solidez dos recursos e da
segurança institucional da localidade em equilíbrio com a visibilidade
estratégica. Referencia atual de score: GSS 7.2 em 2026-03-05.

\subsection{13.}

\begin{itemize}
\tightlist
\item
  \begin{enumerate}
  \def\labelenumi{\arabic{enumi})}
  \tightlist
  \item
    Dados censitários validados (INE\footnote{Instituto Nacional de Estadística (Instituto Nacional de Estatística), órgão oficial de dados demográficos e censitários.} 2022).
  \end{enumerate}
\item
  \begin{enumerate}
  \def\labelenumi{\arabic{enumi})}
  \setcounter{enumi}{1}
  \tightlist
  \item
    Relatórios de produção agrícola e certificação de sementes (MAG).
  \end{enumerate}
\item
  \begin{enumerate}
  \def\labelenumi{\arabic{enumi})}
  \setcounter{enumi}{2}
  \tightlist
  \item
    Mapa de solos e bacias hidrográficas de Canindeyú (MAG).
  \end{enumerate}
\item
  \begin{enumerate}
  \def\labelenumi{\arabic{enumi})}
  \setcounter{enumi}{3}
  \tightlist
  \item
    Registro de segurança pública departamental.
  \end{enumerate}
\end{itemize}

\subsubsection{Combustível}

\textbf{Referência:} PETROPAR /\allowbreak  postos locais (2024) \textbf{Tipo de
localidade:} Interior

\begin{longtable}[]{@{}lll@{}}
\toprule\noalign{}
Combustível & USD/\allowbreak litro & Gs/\allowbreak litro (aprox.) \\
\midrule\noalign{}
\endhead
\bottomrule\noalign{}
\endlastfoot
Gasolina 93 oct & 1.00 & 7,400 \\
Gasolina 97 oct (premium) & 1.10 & 8,140 \\
Diesel & 0.93 & 6,882 \\
\end{longtable}

\begin{quote}
Preços podem variar ±5\% conforme posto e sazonalidade. Chaco e interior
remoto apresentam maior variação.
\end{quote}

\subsubsection{Cobertura Celular}

\textbf{Fonte:} CONATEL PY /\allowbreak  operadoras (2024)

\begin{longtable}[]{@{}ll@{}}
\toprule\noalign{}
Parâmetro & Valor \\
\midrule\noalign{}
\endhead
\bottomrule\noalign{}
\endlastfoot
Cobertura 4G população (dept.) & 78\% \\
Cobertura 4G área rural & 55\% \\
Melhor operadora & Tigo \\
Qualidade rural & moderada \\
\end{longtable}

\begin{quote}
Para áreas rurais fora do núcleo urbano, recomenda-se chip Tigo como
principal e Personal como backup.
\end{quote}

\subsubsection{Internet}

\textbf{Fonte:} CONATEL /\allowbreak  Speedtest Ookla (2024)

\begin{longtable}[]{@{}ll@{}}
\toprule\noalign{}
Parâmetro & Valor \\
\midrule\noalign{}
\endhead
\bottomrule\noalign{}
\endlastfoot
Velocidade média download & 35 Mbps \\
Domicílios com internet (dept.) & 45\% \\
Tecnologia predominante & rádio \\
Opção rural & Starlink disponível (\textasciitilde USD 44/\allowbreak mês) \\
\end{longtable}

\subsubsection{Mercado Imobiliário e Terra
Rural}

\textbf{Fonte:} INDERT /\allowbreak  Clasificados.com.py (2024)

\begin{longtable}[]{@{}ll@{}}
\toprule\noalign{}
Tipo & Referência \\
\midrule\noalign{}
\endhead
\bottomrule\noalign{}
\endlastfoot
Terra agrícola alta prod. (USD/\allowbreak ha) & 7,500 \\
Imóvel urbano (USD/\allowbreak m²) & 650 \\
Aluguel 2 quartos (USD/\allowbreak mês) & 260 \\
\end{longtable}

\begin{quote}
Valores de referência departamental. Localidades menores podem ter
preços 20--40\% abaixo da capital departamental. \#\#\# Saúde
\end{quote}

\textbf{Fonte:} MSPBS /\allowbreak  IPS Paraguay (2024-2026), consolidação
departamental e proxy local

\begin{longtable}[]{@{}ll@{}}
\toprule\noalign{}
Serviço & Disponibilidade \\
\midrule\noalign{}
\endhead
\bottomrule\noalign{}
\endlastfoot
USF /\allowbreak  Posto de Saúde & sim \\
Hospital Regional & não \\
IPS (seguro social) & não \\
Farmácia & sim \\
Distância ao hospital de referência & \textasciitilde37 km (Salto del
Guaira) \\
\end{longtable}

\textbf{Principais estabelecimentos:} USF local; referencia hospitalar
em Salto del Guaira

\textbf{Observação para imigrantes:} Atendimento primario local ou em
raio curto; casos de maior complexidade seguem para Salto del Guaira.
Cobertura privada continua recomendada para especialidades e urgências
de maior complexidade.
\subsubsection{Indicadores Sociais}

\textbf{Fontes:} PNUD 2020, INE\footnote{Instituto Nacional de Estadística (Instituto Nacional de Estatística), órgão oficial de dados demográficos e censitários.} Censo 2022, INE\footnote{Instituto Nacional de Estadística (Instituto Nacional de Estatística), órgão oficial de dados demográficos e censitários.} EPHC 2023, Ministerio
Público 2024\\
\textbf{Nota:} Valores marcados como \emph{(dept.)} referem-se ao
departamento de Canindeyú; valores \emph{(dist.)} são específicos deste
distrito. Dados marcados \emph{(est.)} são estimativas.

\begin{longtable}[]{@{}
  >{\raggedright\arraybackslash}p{(\linewidth - 4\tabcolsep) * \real{0.4231}}
  >{\raggedright\arraybackslash}p{(\linewidth - 4\tabcolsep) * \real{0.2692}}
  >{\raggedright\arraybackslash}p{(\linewidth - 4\tabcolsep) * \real{0.3077}}@{}}
\toprule\noalign{}
\begin{minipage}[b]{\linewidth}\raggedright
Indicador
\end{minipage} & \begin{minipage}[b]{\linewidth}\raggedright
Valor
\end{minipage} & \begin{minipage}[b]{\linewidth}\raggedright
Âmbito
\end{minipage} \\
\midrule\noalign{}
\endhead
\bottomrule\noalign{}
\endlastfoot
IDH (2020) & 0,685 & dept. \\
Ranking IDH nacional & 15/\allowbreak 18 & dept. \\
Esperança de vida & 71,2 anos & dept. \\
Escolaridade média & 7,3 anos & dept. \\
RNB per capita (USD PPA) & 7.100 & dept. \\
Pobreza monetária (\%) & 35,1\% & dept. \\
Pobreza extrema (\%) & 10,8\% & dept. \\
Índice de Gini & 0,467 & dept. \\
Acesso a água potável (\%) & 67,4\% & dept. \\
Acesso a saneamento (\%) & 64,8\% & dept. \\
Taxa de homicídios (est., /\allowbreak 100k hab) & \textasciitilde15--25
(estimativa, significativamente acima da média nacional) & dept. \\
Índice de segurança & baixo & dept. \\
População (Censo 2022) & 7.284 habitantes & dist. \\
Idade mediana & N/\allowbreak D & dist. \\
Taxa de fecundidade & N/\allowbreak D & dist. \\
\end{longtable}
\subsubsection{Combustível}

\textbf{Referência:} PETROPAR /\allowbreak  postos locais (2024) \textbf{Tipo de
localidade:} Interior

\begin{longtable}[]{@{}lll@{}}
\toprule\noalign{}
Combustível & USD/\allowbreak litro & Gs/\allowbreak litro (aprox.) \\
\midrule\noalign{}
\endhead
\bottomrule\noalign{}
\endlastfoot
Gasolina 93 oct & 1.00 & 7,400 \\
Gasolina 97 oct (premium) & 1.10 & 8,140 \\
Diesel & 0.93 & 6,882 \\
\end{longtable}

\begin{quote}
Preços podem variar ±5\% conforme posto e sazonalidade. Chaco e interior
remoto apresentam maior variação.
\end{quote}

\subsubsection{Cobertura Celular}

\textbf{Fonte:} CONATEL PY /\allowbreak  operadoras (2024)

\begin{longtable}[]{@{}ll@{}}
\toprule\noalign{}
Parâmetro & Valor \\
\midrule\noalign{}
\endhead
\bottomrule\noalign{}
\endlastfoot
Cobertura 4G população (dept.) & 78\% \\
Cobertura 4G área rural & 55\% \\
Melhor operadora & Tigo \\
Qualidade rural & moderada \\
\end{longtable}

\begin{quote}
Para áreas rurais fora do núcleo urbano, recomenda-se chip Tigo como
principal e Personal como backup.
\end{quote}

\subsubsection{Internet}

\textbf{Fonte:} CONATEL /\allowbreak  Speedtest Ookla (2024)

\begin{longtable}[]{@{}ll@{}}
\toprule\noalign{}
Parâmetro & Valor \\
\midrule\noalign{}
\endhead
\bottomrule\noalign{}
\endlastfoot
Velocidade média download & 35 Mbps \\
Domicílios com internet (dept.) & 45\% \\
Tecnologia predominante & rádio \\
Opção rural & Starlink disponível (\textasciitilde USD 44/\allowbreak mês) \\
\end{longtable}

\subsubsection{Mercado Imobiliário e Terra
Rural}

\textbf{Fonte:} INDERT /\allowbreak  Clasificados.com.py (2024)

\begin{longtable}[]{@{}ll@{}}
\toprule\noalign{}
Tipo & Referência \\
\midrule\noalign{}
\endhead
\bottomrule\noalign{}
\endlastfoot
Terra agrícola alta prod. (USD/\allowbreak ha) & 7,500 \\
Imóvel urbano (USD/\allowbreak m²) & 650 \\
Aluguel 2 quartos (USD/\allowbreak mês) & 260 \\
\end{longtable}

\begin{quote}
Valores de referência departamental. Localidades menores podem ter
preços 20--40\% abaixo da capital departamental. \#\#\# Saúde
\end{quote}

\textbf{Fonte:} MSPBS /\allowbreak  IPS Paraguay (2024-2026), consolidação
departamental e proxy local

\begin{longtable}[]{@{}ll@{}}
\toprule\noalign{}
Serviço & Disponibilidade \\
\midrule\noalign{}
\endhead
\bottomrule\noalign{}
\endlastfoot
USF /\allowbreak  Posto de Saúde & sim \\
Hospital Regional & não \\
IPS (seguro social) & não \\
Farmácia & sim \\
Distância ao hospital de referência & \textasciitilde37 km (Salto del
Guaira) \\
\end{longtable}

\textbf{Principais estabelecimentos:} USF local; referencia hospitalar
em Salto del Guaira

\textbf{Observação para imigrantes:} Atendimento primario local ou em
raio curto; casos de maior complexidade seguem para Salto del Guaira.
Cobertura privada continua recomendada para especialidades e urgências
de maior complexidade.
\newpage
\secaoDiagnostico{Laurel}{\mapaConteudo{-24.5166}{-54.7500}}{Laurel é um dos ``segredos'' produtivos do leste paraguaio. Como um
município de criação recente, oferece o frescor de uma fronteira
agrícola em consolidação com o benefício de uma densidade demográfica
quase nula, garantindo privacidade e isolamento tático superior. Sua
força reside nos solos de elite e na abundância de recursos naturais. O
desafio principal é a fragilidade institucional local e a dependência de
logística externa para serviços vitais. É o local ideal para quem busca
autonomia técnica e deseja viver ``fora do radar'' metropolitano em uma
das regiões mais férteis do país.

\subsubsection{Por que sim}

\begin{itemize}
\tightlist
\item
  Solos de altíssima produtividade e abundância de águas puras
  subterrâneas.
\item
  Baixíssima densidade populacional (privacidade absoluta).
\item
  Localização discreta e protegida de eixos de conflito primários.
\item
  Sem restrições da Lei de Fronteira.
\end{itemize}

\subsubsection{Por que não}

\begin{itemize}
\tightlist
\item
  Infraestrutura institucional e de saúde local em fase inicial.
\item
  Dependência de centros maiores para suprimentos técnicos e básicos.
\item
  Insegurança tática regional latente no departamento de Canindeyú.
\end{itemize}

\subsubsection{Recomendação
Técnica}

Altamente indicado para estabelecimentos de autossuficiência que busquem
isolamento tático em solo de alta performance. Requer investimento em
autonomia energética, médica e logística. O GSS de 6.5 reflete a
segurança do isolamento em equilíbrio com a vulnerabilidade
institucional da localidade. Referencia atual de score: GSS 6.5 em
2026-03-05.

\subsubsection{}

\begin{itemize}
\tightlist
\item
  \begin{enumerate}
  \def\labelenumi{\arabic{enumi})}
  \tightlist
  \item
    Dados censitários validados (INE\footnote{Instituto Nacional de Estadística (Instituto Nacional de Estatística), órgão oficial de dados demográficos e censitários.} 2022).
  \end{enumerate}
\item
  \begin{enumerate}
  \def\labelenumi{\arabic{enumi})}
  \setcounter{enumi}{1}
  \tightlist
  \item
    Lei de criação do município (TSJE 2020).
  \end{enumerate}
\item
  \begin{enumerate}
  \def\labelenumi{\arabic{enumi})}
  \setcounter{enumi}{2}
  \tightlist
  \item
    Mapa de solos e aptidão agrícola de Canindeyú (Portal
    Geoestadístico).
  \end{enumerate}
\item
  \begin{enumerate}
  \def\labelenumi{\arabic{enumi})}
  \setcounter{enumi}{3}
  \tightlist
  \item
    Registro de segurança pública departamental.
  \end{enumerate}
\end{itemize}}
\begin{table}[ht!]
\small \begin{tabular}{p{3.0cm}p{0.8cm}p{6.0cm}} \toprule
\textbf{Critério} & \textbf{Nota} & \textbf{Justificativa} \\ \midrule
A - Ameaças Estratégicas & 6.5 & Localização remota e protegida; excelente invisibilidade tática, equilibrada pelo risco regional \\
B - Risco Social & 8.0 & Densidade populacional extremamente baixa; isolamento social superior \\
C - Riscos Naturais & 7.0 & Geologia muito estável e baixo risco de desastres naturais graves \\
D - Autossuficiência & 7.0 & Bons recursos agrícolas e hídricos; limitado pela escala comercial local incipiente \\
E - Institucional & 4.0 & Suporte institucional mínimo; serviços dependentes de cidades vizinhas \\
\midrule \textbf{GSS FINAL} & \textbf{6.5} & \textbf{Moderadamente Seguro (Ideal para quem busca isolamento rural em solo de alta produtividade, aceitando a precariedade institucional local).} \\ \bottomrule \end{tabular}
\end{table}
\subsubsection{Vulnerabilidades e
Mitigação}\label{vuln-Laurel-vulnerabilidades-e-mitigauxe7uxe3o}

\begin{itemize}
\tightlist
\item
  \textbf{Vulnerabilidade Institucional:} Suporte governamental local
  ainda em fase de estruturação (Mitigação: manutenção de sistemas de
  segurança patrimonial privada e autonomia médica).
\item
  \textbf{Isolamento de Saúde:} Distância significativa de centros
  hospitalares cirúrgicos (Mitigação: manutenção de estoque médico
  básico e kit trauma local).
\item
  \textbf{Fragilidade das Vias:} Dependência de estradas internas de
  terra (Mitigação: posse de veículo 4x4 robusto e manutenção de estoque
  de segurança para 60 dias).
\end{itemize}
\subsubsection{Dossiê de Campo}\label{dossie-Laurel-dossiuxea-de-campo}
\begin{description}[style=nextline,leftmargin=0.5cm,font=\bfseries]
\item[Coordenadas]  24°31'S 54°45'W.
\item[Topographical]  Relevo predominantemente plano a suavemente ondulado, situado em zona de expansão da fronteira agrícola no leste de Canindeyú. Altitude \textasciitilde 240m. Área de \textasciitilde 408 km². Geologicamente estável, risco sísmico nulo.
\item[Alvos Estratégicos]  Áreas de produção agrícola intensiva mecanizada; Silos de armazenamento de grãos em consolidação. Ausência total de alvos militares ou infraestrutura industrial pesada, oferecendo excelente invisibilidade estratégica e proteção por anonimato.
\item[Fallout]  Ventos predominantes do Norte (NE) e Leste. Baixo risco de fallout direto; localização remota em relação aos centros metropolitanos.
\item[População]  2.586 habitantes (Censo 2022 Final). Município de criação recente, em fase de consolidação demográfica urbana e rural.
\item[Densidade]  \textasciitilde 6,3 hab/\allowbreak km² (Densidade extremamente baixa, garantindo privacidade tática absoluta e baixíssimo risco de agitação social de massa).
\item[Vias de Acesso]  Acesso via ramais de terra e alguns trechos pavimentados secundários que conectam às Rutas PY03 e PY10. Localização protegida por estar fora dos eixos de circulação principal do país, o que favorece o isolamento deliberado.
\item[Serviços]  Infraestrutura institucional básica em desenvolvimento. Unidade de Saúde da Família (USF) local. Dependência de Nueva Esperanza (40 km) ou Curuguaty para serviços hospitalares de alta complexidade e logística comercial avançada.
\item[Custo de Vida]  Muito baixo; economia baseada na produção de grãos e pecuária.
\item[Preço da Terra]  US\$ 6.000-9.000/\allowbreak ha (terras mecanizadas); US\$ 2.500-4.000/\allowbreak ha (áreas de uso misto).
\item[Hidrologia]  Baixo risco de inundações sistêmicas devido à topografia favorável. Presença de arroios perenes de boa qualidade hídrica. Risco moderado de isolamento rural temporário por danos em estradas de terra em períodos chuvosos (média 1.700 mm/\allowbreak ano).
\item[Clima]  Subtropical Úmido. Verões quentes e chuvosos; invernos amenos com geadas ocasionais.
\item[Energia]  Rede nacional (Itaipu); instável em áreas rurais devido à expansão recente da rede e falta de redundância.
\item[Água]  Abundante via poços artesianos e lençol freático produtivo; excelente qualidade hídrica subterrânea.
\item[Qualidade do Solo]  Latossolo Vermelho; solo profundo e fértil, ideal para soja, milho e pastagens mecanizadas.
\item[Resources Locais]  Grande produção de grãos e proteína animal. Elevadíssimo potencial para autossuficiência alimentar básica em nível de propriedade devido à vasta área produtiva e baixíssima população.
\item[Segurança]  Zona rural excepcionalmente pacífica em termos de crime comum. No entanto, exige atenção tática regional por estar em departamento de fronteira sensível ao trânsito de ilícitos e conflitos agrários latentes. Estabilidade social baseada na cultura de trabalho do agronegócio.
\item[Leis Local]  Município de criação muito recente (2020). \\ \textbf{Livre de Restrição} de Fronteira. \\ \textit{Aquisição de terra rural proibida para estrangeiros limítrofes.}
\end{description}
\paragraph{Solo (SoilGrids 2.0, média ponderada 0--30
cm)}

\begin{longtable}[]{@{}ll@{}}
\toprule\noalign{}
Parâmetro & Valor \\
\midrule\noalign{}
\endhead
\bottomrule\noalign{}
\endlastfoot
pH (H$_2$O) & 5.4 \\
Carbono orgânico (SOC) & 23.83 g/\allowbreak kg \\
Argila & 47.17 \% \\
Areia & 27.9 \% \\
Silte (calc.) & 24.9 \% \\
Densidade aparente & 1.19 g/\allowbreak cm³ \\
\textbf{Aptidão agrícola} & \textbf{Média} \\
\end{longtable}

Fonte: ISRIC SoilGrids 2.0 via WCS (acesso em 2026-03-21). Coords:
-24.5167°, -54.75°. Média ponderada camadas 0-5, 5-15, 15-30 cm.

\begin{itemize}
\tightlist
\item
  \textbf{Homicidios/\allowbreak 100k:} \textasciitilde18,0 (Regional/\allowbreak Canindeyú).
\end{itemize}

\subsection{10. Analise de Riscos}

\begin{itemize}
\tightlist
\item
  \textbf{Cenario curto prazo:} Manutenção da tranquilidade rural e
  desenvolvimento do núcleo urbano municipal.
\item
  \textbf{Cenario medio prazo:} Valorização das terras pela consolidação
  da infraestrutura asfáltica vicinal.
\end{itemize}

\subsection{11. Redacao Final}

\subsubsection{Diagnostico Integrado}

Laurel é um dos ``segredos'' produtivos do leste paraguaio. Como um
município de criação recente, oferece o frescor de uma fronteira
agrícola em consolidação com o benefício de uma densidade demográfica
quase nula, garantindo privacidade e isolamento tático superior. Sua
força reside nos solos de elite e na abundância de recursos naturais. O
desafio principal é a fragilidade institucional local e a dependência de
logística externa para serviços vitais. É o local ideal para quem busca
autonomia técnica e deseja viver ``fora do radar'' metropolitano em uma
das regiões mais férteis do país.

\subsubsection{Por que sim}

\begin{itemize}
\tightlist
\item
  Solos de altíssima produtividade e abundância de águas puras
  subterrâneas.
\item
  Baixíssima densidade populacional (privacidade absoluta).
\item
  Localização discreta e protegida de eixos de conflito primários.
\item
  Sem restrições da Lei de Fronteira.
\end{itemize}

\subsubsection{Por que nao}

\begin{itemize}
\tightlist
\item
  Infraestrutura institucional e de saúde local em fase inicial.
\item
  Dependência de centros maiores para suprimentos técnicos e básicos.
\item
  Insegurança tática regional latente no departamento de Canindeyú.
\end{itemize}

\subsubsection{Recomendacao Tecnica}

Altamente indicado para estabelecimentos de autossuficiência que busquem
isolamento tático em solo de alta performance. Requer investimento em
autonomia energética, médica e logística. O GSS de 6.5 reflete a
segurança do isolamento em equilíbrio com a vulnerabilidade
institucional da localidade. Referencia atual de score: GSS 6.5 em
2026-03-05.

\subsection{13.}

\begin{itemize}
\tightlist
\item
  \begin{enumerate}
  \def\labelenumi{\arabic{enumi})}
  \tightlist
  \item
    Dados censitários validados (INE\footnote{Instituto Nacional de Estadística (Instituto Nacional de Estatística), órgão oficial de dados demográficos e censitários.} 2022).
  \end{enumerate}
\item
  \begin{enumerate}
  \def\labelenumi{\arabic{enumi})}
  \setcounter{enumi}{1}
  \tightlist
  \item
    Lei de criação do município (TSJE 2020).
  \end{enumerate}
\item
  \begin{enumerate}
  \def\labelenumi{\arabic{enumi})}
  \setcounter{enumi}{2}
  \tightlist
  \item
    Mapa de solos e aptidão agrícola de Canindeyú (Portal
    Geoestadístico).
  \end{enumerate}
\item
  \begin{enumerate}
  \def\labelenumi{\arabic{enumi})}
  \setcounter{enumi}{3}
  \tightlist
  \item
    Registro de segurança pública departamental.
  \end{enumerate}
\end{itemize}

\subsubsection{Combustível}

\textbf{Referência:} PETROPAR /\allowbreak  postos locais (2024) \textbf{Tipo de
localidade:} Interior

\begin{longtable}[]{@{}lll@{}}
\toprule\noalign{}
Combustível & USD/\allowbreak litro & Gs/\allowbreak litro (aprox.) \\
\midrule\noalign{}
\endhead
\bottomrule\noalign{}
\endlastfoot
Gasolina 93 oct & 1.00 & 7,400 \\
Gasolina 97 oct (premium) & 1.10 & 8,140 \\
Diesel & 0.93 & 6,882 \\
\end{longtable}

\begin{quote}
Preços podem variar ±5\% conforme posto e sazonalidade. Chaco e interior
remoto apresentam maior variação.
\end{quote}

\subsubsection{Cobertura Celular}

\textbf{Fonte:} CONATEL PY /\allowbreak  operadoras (2024)

\begin{longtable}[]{@{}ll@{}}
\toprule\noalign{}
Parâmetro & Valor \\
\midrule\noalign{}
\endhead
\bottomrule\noalign{}
\endlastfoot
Cobertura 4G população (dept.) & 78\% \\
Cobertura 4G área rural & 55\% \\
Melhor operadora & Tigo \\
Qualidade rural & moderada \\
\end{longtable}

\begin{quote}
Para áreas rurais fora do núcleo urbano, recomenda-se chip Tigo como
principal e Personal como backup.
\end{quote}

\subsubsection{Internet}

\textbf{Fonte:} CONATEL /\allowbreak  Speedtest Ookla (2024)

\begin{longtable}[]{@{}ll@{}}
\toprule\noalign{}
Parâmetro & Valor \\
\midrule\noalign{}
\endhead
\bottomrule\noalign{}
\endlastfoot
Velocidade média download & 35 Mbps \\
Domicílios com internet (dept.) & 45\% \\
Tecnologia predominante & rádio \\
Opção rural & Starlink disponível (\textasciitilde USD 44/\allowbreak mês) \\
\end{longtable}

\subsubsection{Mercado Imobiliário e Terra
Rural}

\textbf{Fonte:} INDERT /\allowbreak  Clasificados.com.py (2024)

\begin{longtable}[]{@{}ll@{}}
\toprule\noalign{}
Tipo & Referência \\
\midrule\noalign{}
\endhead
\bottomrule\noalign{}
\endlastfoot
Terra agrícola alta prod. (USD/\allowbreak ha) & 7,500 \\
Imóvel urbano (USD/\allowbreak m²) & 650 \\
Aluguel 2 quartos (USD/\allowbreak mês) & 260 \\
\end{longtable}

\begin{quote}
Valores de referência departamental. Localidades menores podem ter
preços 20--40\% abaixo da capital departamental. \#\#\# Saúde
\end{quote}

\textbf{Fonte:} MSPBS /\allowbreak  IPS Paraguay (2024-2026), consolidação
departamental e proxy local

\begin{longtable}[]{@{}ll@{}}
\toprule\noalign{}
Serviço & Disponibilidade \\
\midrule\noalign{}
\endhead
\bottomrule\noalign{}
\endlastfoot
USF /\allowbreak  Posto de Saúde & sim \\
Hospital Regional & sim \\
IPS (seguro social) & não \\
Farmácia & sim \\
Distância ao hospital de referência & \textasciitilde106 km (Salto del
Guaira) \\
\end{longtable}

\textbf{Principais estabelecimentos:} USF local; referencia hospitalar
em Salto del Guaira

\textbf{Observação para imigrantes:} Atendimento primario local ou em
raio curto; casos de maior complexidade seguem para Salto del Guaira.
Cobertura privada continua recomendada para especialidades e urgências
de maior complexidade.
\subsubsection{Indicadores Sociais}

\textbf{Fontes:} PNUD 2020, INE\footnote{Instituto Nacional de Estadística (Instituto Nacional de Estatística), órgão oficial de dados demográficos e censitários.} Censo 2022, INE\footnote{Instituto Nacional de Estadística (Instituto Nacional de Estatística), órgão oficial de dados demográficos e censitários.} EPHC 2023, Ministerio
Público 2024\\
\textbf{Nota:} Valores marcados como \emph{(dept.)} referem-se ao
departamento de Canindeyú; valores \emph{(dist.)} são específicos deste
distrito. Dados marcados \emph{(est.)} são estimativas.

\begin{longtable}[]{@{}
  >{\raggedright\arraybackslash}p{(\linewidth - 4\tabcolsep) * \real{0.4231}}
  >{\raggedright\arraybackslash}p{(\linewidth - 4\tabcolsep) * \real{0.2692}}
  >{\raggedright\arraybackslash}p{(\linewidth - 4\tabcolsep) * \real{0.3077}}@{}}
\toprule\noalign{}
\begin{minipage}[b]{\linewidth}\raggedright
Indicador
\end{minipage} & \begin{minipage}[b]{\linewidth}\raggedright
Valor
\end{minipage} & \begin{minipage}[b]{\linewidth}\raggedright
Âmbito
\end{minipage} \\
\midrule\noalign{}
\endhead
\bottomrule\noalign{}
\endlastfoot
IDH (2020) & 0,685 & dept. \\
Ranking IDH nacional & 15/\allowbreak 18 & dept. \\
Esperança de vida & 71,2 anos & dept. \\
Escolaridade média & 7,3 anos & dept. \\
RNB per capita (USD PPA) & 7.100 & dept. \\
Pobreza monetária (\%) & 35,1\% & dept. \\
Pobreza extrema (\%) & 10,8\% & dept. \\
Índice de Gini & 0,467 & dept. \\
Acesso a água potável (\%) & 67,4\% & dept. \\
Acesso a saneamento (\%) & 64,8\% & dept. \\
Taxa de homicídios (est., /\allowbreak 100k hab) & \textasciitilde15--25
(estimativa, significativamente acima da média nacional) & dept. \\
Índice de segurança & baixo & dept. \\
População (Censo 2022) & 2.586 habitantes & dist. \\
Idade mediana & N/\allowbreak D & dist. \\
Taxa de fecundidade & N/\allowbreak D & dist. \\
\end{longtable}
\subsubsection{Combustível}

\textbf{Referência:} PETROPAR /\allowbreak  postos locais (2024) \textbf{Tipo de
localidade:} Interior

\begin{longtable}[]{@{}lll@{}}
\toprule\noalign{}
Combustível & USD/\allowbreak litro & Gs/\allowbreak litro (aprox.) \\
\midrule\noalign{}
\endhead
\bottomrule\noalign{}
\endlastfoot
Gasolina 93 oct & 1.00 & 7,400 \\
Gasolina 97 oct (premium) & 1.10 & 8,140 \\
Diesel & 0.93 & 6,882 \\
\end{longtable}

\begin{quote}
Preços podem variar ±5\% conforme posto e sazonalidade. Chaco e interior
remoto apresentam maior variação.
\end{quote}

\subsubsection{Cobertura Celular}

\textbf{Fonte:} CONATEL PY /\allowbreak  operadoras (2024)

\begin{longtable}[]{@{}ll@{}}
\toprule\noalign{}
Parâmetro & Valor \\
\midrule\noalign{}
\endhead
\bottomrule\noalign{}
\endlastfoot
Cobertura 4G população (dept.) & 78\% \\
Cobertura 4G área rural & 55\% \\
Melhor operadora & Tigo \\
Qualidade rural & moderada \\
\end{longtable}

\begin{quote}
Para áreas rurais fora do núcleo urbano, recomenda-se chip Tigo como
principal e Personal como backup.
\end{quote}

\subsubsection{Internet}

\textbf{Fonte:} CONATEL /\allowbreak  Speedtest Ookla (2024)

\begin{longtable}[]{@{}ll@{}}
\toprule\noalign{}
Parâmetro & Valor \\
\midrule\noalign{}
\endhead
\bottomrule\noalign{}
\endlastfoot
Velocidade média download & 35 Mbps \\
Domicílios com internet (dept.) & 45\% \\
Tecnologia predominante & rádio \\
Opção rural & Starlink disponível (\textasciitilde USD 44/\allowbreak mês) \\
\end{longtable}

\subsubsection{Mercado Imobiliário e Terra
Rural}

\textbf{Fonte:} INDERT /\allowbreak  Clasificados.com.py (2024)

\begin{longtable}[]{@{}ll@{}}
\toprule\noalign{}
Tipo & Referência \\
\midrule\noalign{}
\endhead
\bottomrule\noalign{}
\endlastfoot
Terra agrícola alta prod. (USD/\allowbreak ha) & 7,500 \\
Imóvel urbano (USD/\allowbreak m²) & 650 \\
Aluguel 2 quartos (USD/\allowbreak mês) & 260 \\
\end{longtable}

\begin{quote}
Valores de referência departamental. Localidades menores podem ter
preços 20--40\% abaixo da capital departamental. \#\#\# Saúde
\end{quote}

\textbf{Fonte:} MSPBS /\allowbreak  IPS Paraguay (2024-2026), consolidação
departamental e proxy local

\begin{longtable}[]{@{}ll@{}}
\toprule\noalign{}
Serviço & Disponibilidade \\
\midrule\noalign{}
\endhead
\bottomrule\noalign{}
\endlastfoot
USF /\allowbreak  Posto de Saúde & sim \\
Hospital Regional & sim \\
IPS (seguro social) & não \\
Farmácia & sim \\
Distância ao hospital de referência & \textasciitilde106 km (Salto del
Guaira) \\
\end{longtable}

\textbf{Principais estabelecimentos:} USF local; referencia hospitalar
em Salto del Guaira

\textbf{Observação para imigrantes:} Atendimento primario local ou em
raio curto; casos de maior complexidade seguem para Salto del Guaira.
Cobertura privada continua recomendada para especialidades e urgências
de maior complexidade.
\newpage
\secaoDiagnostico{Maracana}{\mapaConteudo{-24.1666}{-55.8333}}{Maracaná é um dos enclaves produtivos mais promissores e discretos de
Canindeyú. Como um município de criação recente, oferece a privacidade
de um vasto território rural com solos de produtividade mundial. Sua
força reside na excepcional abundância de recursos hídricos e
alimentares, e na segurança social forjada por uma comunidade produtora
resiliente. É o local ideal para quem busca anonimato tático em solo
fértil, transformando o distanciamento metropolitano em uma barreira
natural de proteção contra crises sistêmicas.

\subsubsection{Por que sim}

\begin{itemize}
\tightlist
\item
  Solos de altíssima produtividade e abundância de águas puras
  subterrâneas.
\item
  Baixíssima densidade populacional (privacidade absoluta).
\item
  Localização discreta e protegida de grandes eixos de conflito.
\item
  Sem restrições da Lei de Fronteira.
\end{itemize}

\subsubsection{Por que não}

\begin{itemize}
\tightlist
\item
  Infraestrutura de saúde local limitada a atendimentos primários.
\item
  Insegurança tática regional latente no departamento de Canindeyú.
\item
  Pequena escala comercial local exige planejamento logístico de
  suprimentos.
\end{itemize}

\subsubsection{Recomendação
Técnica}

Altamente indicado para estabelecimentos de autossuficiência familiar de
alto padrão que valorizem o isolamento e a produtividade. Requer
investimento em segurança patrimonial e autonomia em saúde. O GSS de 7.6
reflete a excepcional solidez do isolamento geográfico reforçada pela
base de recursos naturais. Referencia atual de score: GSS 7.6 em
2026-03-05.

\subsubsection{}

\begin{itemize}
\tightlist
\item
  \begin{enumerate}
  \def\labelenumi{\arabic{enumi})}
  \tightlist
  \item
    Dados censitários validados (INE\footnote{Instituto Nacional de Estadística (Instituto Nacional de Estatística), órgão oficial de dados demográficos e censitários.} 2022).
  \end{enumerate}
\item
  \begin{enumerate}
  \def\labelenumi{\arabic{enumi})}
  \setcounter{enumi}{1}
  \tightlist
  \item
    Lei de criação do município (TSJE 2016).
  \end{enumerate}
\item
  \begin{enumerate}
  \def\labelenumi{\arabic{enumi})}
  \setcounter{enumi}{2}
  \tightlist
  \item
    Mapa de solos e aptidão agrícola de Canindeyú (Portal
    Geoestadístico).
  \end{enumerate}
\item
  \begin{enumerate}
  \def\labelenumi{\arabic{enumi})}
  \setcounter{enumi}{3}
  \tightlist
  \item
    Registro de segurança pública departamental.
  \end{enumerate}
\end{itemize}}
\begin{table}[ht!]
\small \begin{tabular}{p{3.0cm}p{0.8cm}p{6.0cm}} \toprule
\textbf{Critério} & \textbf{Nota} & \textbf{Justificativa} \\ \midrule
A - Ameaças Estratégicas & 7.5 & Localização remota e discreta; excelente invisibilidade tática fora dos eixos de conflito primários \\
B - Risco Social & 8.0 & Densidade populacional baixa; isolamento social superior e privacidade absoluta \\
C - Riscos Naturais & 7.5 & Geologia muito estável e relevo seguro contra inundações sistêmicas \\
D - Autossuficiência & 8.0 & Recursos agrícolas excelentes e abundância hídrica subterrânea excepcional \\
E - Institucional & 7.0 & Forte segurança social baseada na cultura produtora e suporte em expansão como novo município \\
\midrule \textbf{GSS FINAL} & \textbf{7.6} & \textbf{Seguro (Altamente recomendado para quem busca isolamento tático e autossuficiência em solo de elite, em ambiente de paz social).} \\ \bottomrule \end{tabular}
\end{table}
\subsubsection{Vulnerabilidades e
Mitigação}\label{vuln-Maracana-vulnerabilidades-e-mitigauxe7uxe3o}

\begin{itemize}
\tightlist
\item
  \textbf{Vulnerabilidade Tática Regional:} Insegurança ligada ao
  contexto de segurança do departamento (Mitigação: aquisição de
  propriedades com documentação legal perfeita e investimento em
  segurança patrimonial).
\item
  \textbf{Isolamento de Saúde:} Distância significativa de hospitais
  cirúrgicos (Mitigação: manutenção de estoque médico básico e kit
  trauma local).
\item
  \textbf{Dependência de Serviços Externos:} Pequena escala comercial
  local exige planejamento logístico para suprimentos técnicos.
\end{itemize}
\subsubsection{Dossiê de Campo}\label{dossie-Maracana-dossiuxea-de-campo}
\begin{description}[style=nextline,leftmargin=0.5cm,font=\bfseries]
\item[Coordenadas]  24°10'S 55°50'W.
\item[Topografia]  Relevo ondulado com colinas suaves e solos de altíssima produtividade (Terra Roxa). Altitude \textasciitilde 260m. Área massiva de \textasciitilde 1.119 km². Geologicamente muito estável, risco sísmico desprezível.
\item[Alvos Estratégicos]  Áreas de produção agrícola intensiva (soja, milho); Polo de assentamentos rurais consolidados; Proximidade tática com o Rio Jejuí Guazú ao norte. Ausência de alvos militares diretos, oferecendo excelente invisibilidade estratégica e proteção por anonimato rural.
\item[Fallout]  Ventos predominantes do Norte (NE) e Leste. Baixíssimo risco de fallout direto; posição isolada no coração produtivo do departamento.
\item[População]  15.357 habitantes (Censo 2022 Final). Município de criação recente (2016), com perfil demográfico essencialmente rural e forte componente de agricultura familiar mecanizada.
\item[Densidade]  \textasciitilde 13,7 hab/\allowbreak km² (Densidade baixa, oferecendo isolamento tático e privacidade fora dos pequenos núcleos urbanos em formação).
\item[Vias de Acesso]  Acesso via Ruta PY13 e ramais pavimentados secundários. Localização que favorece a discrição, mantendo-se fora dos eixos de trânsito internacional pesado. Malha de estradas internas de terra exige manutenção constante em períodos chuvosos.
\item[Serviços]  Infraestrutura institucional básica em fase de estruturação. Centro de Saúde local e USF. Dependência de Curuguaty (\textasciitilde 50 km) para atendimentos de alta complexidade e logística comercial avançada.
\item[Custo de Vida]  Muito baixo; economia baseada na produção de grãos, gado e subsistência rural.
\item[Preço da Terra]  US\$ 7.000-10.000/\allowbreak ha (terras mecanizadas); US\$ 2.000-4.000/\allowbreak ha (áreas de uso misto/\allowbreak pastagem).
\item[Hidrologia]  Baixo risco de inundações sistêmicas devido à topografia ondulada. Diversos arroios perenes garantem abundância hídrica superficial para irrigação e pecuária.
\item[Clima]  Subtropical Úmido. Verões intensos e chuvosos; invernos amenos com geadas ocasionais.
\item[Energia]  Rede nacional (Itaipu); instável em áreas rurais remotas.
\item[Água]  Abundante via nascentes naturais e poços artesianos profundos; excelente qualidade hídrica subterrânea.
\item[Qualidade do Solo]  Latossolo Vermelho de Alta Performance; solo profundo e fértil. Excelente aptidão para soja, milho e mandioca.
\item[Resources Locais]  Grande produção de grãos e proteína animal. Elevadíssimo potencial para autossuficiência alimentar básica em nível de propriedade devido à vasta área produtiva e baixa população.
\item[Segurança]  Zona rural tradicionalmente pacífica em termos de crime comum. No entanto, exige vigilância tática regional devido ao contexto de Canindeyú (histórico de tensões agrárias e movimentos camponeses organizados). Estabilidade social sustentada pela cultura produtora e laços comunitários fortes.
\item[Leis Local: Município consolidado. Livre de Restrição de Fronteira]  Localizado fora da faixa de 50km da fronteira seca internacional.
\end{description}
\paragraph{Solo (SoilGrids 2.0, média ponderada 0--30
cm)}

\begin{longtable}[]{@{}ll@{}}
\toprule\noalign{}
Parâmetro & Valor \\
\midrule\noalign{}
\endhead
\bottomrule\noalign{}
\endlastfoot
pH (H$_2$O) & 5.4 \\
Carbono orgânico (SOC) & 17.45 g/\allowbreak kg \\
Argila & 23.72 \% \\
Areia & 58.62 \% \\
Silte (calc.) & 17.7 \% \\
Densidade aparente & 1.31 g/\allowbreak cm³ \\
\textbf{Aptidão agrícola} & \textbf{Média} \\
\end{longtable}

Fonte: ISRIC SoilGrids 2.0 via WCS (acesso em 2026-03-21). Coords:
-24.1667°, -55.8333°. Média ponderada camadas 0-5, 5-15, 15-30 cm.

\begin{itemize}
\tightlist
\item
  \textbf{Homicidios/\allowbreak 100k:} \textasciitilde18,0 (Regional/\allowbreak Canindeyú).
\end{itemize}

\subsection{10. Analise de Riscos}

\begin{itemize}
\tightlist
\item
  \textbf{Cenario curto prazo:} Manutenção da tranquilidade rural e
  desenvolvimento do núcleo urbano municipal.
\item
  \textbf{Cenario medio prazo:} Valorização das terras pela consolidação
  da infraestrutura asfáltica vicinal e expansão da produção.
\end{itemize}

\subsection{11. Redacao Final}

\subsubsection{Diagnostico Integrado}

Maracaná é um dos enclaves produtivos mais promissores e discretos de
Canindeyú. Como um município de criação recente, oferece a privacidade
de um vasto território rural com solos de produtividade mundial. Sua
força reside na excepcional abundância de recursos hídricos e
alimentares, e na segurança social forjada por uma comunidade produtora
resiliente. É o local ideal para quem busca anonimato tático em solo
fértil, transformando o distanciamento metropolitano em uma barreira
natural de proteção contra crises sistêmicas.

\subsubsection{Por que sim}

\begin{itemize}
\tightlist
\item
  Solos de altíssima produtividade e abundância de águas puras
  subterrâneas.
\item
  Baixíssima densidade populacional (privacidade absoluta).
\item
  Localização discreta e protegida de grandes eixos de conflito.
\item
  Sem restrições da Lei de Fronteira.
\end{itemize}

\subsubsection{Por que nao}

\begin{itemize}
\tightlist
\item
  Infraestrutura de saúde local limitada a atendimentos primários.
\item
  Insegurança tática regional latente no departamento de Canindeyú.
\item
  Pequena escala comercial local exige planejamento logístico de
  suprimentos.
\end{itemize}

\subsubsection{Recomendacao Tecnica}

Altamente indicado para estabelecimentos de autossuficiência familiar de
alto padrão que valorizem o isolamento e a produtividade. Requer
investimento em segurança patrimonial e autonomia em saúde. O GSS de 7.6
reflete a excepcional solidez do isolamento geográfico reforçada pela
base de recursos naturais. Referencia atual de score: GSS 7.6 em
2026-03-05.

\subsection{13.}

\begin{itemize}
\tightlist
\item
  \begin{enumerate}
  \def\labelenumi{\arabic{enumi})}
  \tightlist
  \item
    Dados censitários validados (INE\footnote{Instituto Nacional de Estadística (Instituto Nacional de Estatística), órgão oficial de dados demográficos e censitários.} 2022).
  \end{enumerate}
\item
  \begin{enumerate}
  \def\labelenumi{\arabic{enumi})}
  \setcounter{enumi}{1}
  \tightlist
  \item
    Lei de criação do município (TSJE 2016).
  \end{enumerate}
\item
  \begin{enumerate}
  \def\labelenumi{\arabic{enumi})}
  \setcounter{enumi}{2}
  \tightlist
  \item
    Mapa de solos e aptidão agrícola de Canindeyú (Portal
    Geoestadístico).
  \end{enumerate}
\item
  \begin{enumerate}
  \def\labelenumi{\arabic{enumi})}
  \setcounter{enumi}{3}
  \tightlist
  \item
    Registro de segurança pública departamental.
  \end{enumerate}
\end{itemize}

\subsubsection{Combustível}

\textbf{Referência:} PETROPAR /\allowbreak  postos locais (2024) \textbf{Tipo de
localidade:} Interior

\begin{longtable}[]{@{}lll@{}}
\toprule\noalign{}
Combustível & USD/\allowbreak litro & Gs/\allowbreak litro (aprox.) \\
\midrule\noalign{}
\endhead
\bottomrule\noalign{}
\endlastfoot
Gasolina 93 oct & 1.00 & 7,400 \\
Gasolina 97 oct (premium) & 1.10 & 8,140 \\
Diesel & 0.93 & 6,882 \\
\end{longtable}

\begin{quote}
Preços podem variar ±5\% conforme posto e sazonalidade. Chaco e interior
remoto apresentam maior variação.
\end{quote}

\subsubsection{Cobertura Celular}

\textbf{Fonte:} CONATEL PY /\allowbreak  operadoras (2024)

\begin{longtable}[]{@{}ll@{}}
\toprule\noalign{}
Parâmetro & Valor \\
\midrule\noalign{}
\endhead
\bottomrule\noalign{}
\endlastfoot
Cobertura 4G população (dept.) & 78\% \\
Cobertura 4G área rural & 55\% \\
Melhor operadora & Tigo \\
Qualidade rural & moderada \\
\end{longtable}

\begin{quote}
Para áreas rurais fora do núcleo urbano, recomenda-se chip Tigo como
principal e Personal como backup.
\end{quote}

\subsubsection{Internet}

\textbf{Fonte:} CONATEL /\allowbreak  Speedtest Ookla (2024)

\begin{longtable}[]{@{}ll@{}}
\toprule\noalign{}
Parâmetro & Valor \\
\midrule\noalign{}
\endhead
\bottomrule\noalign{}
\endlastfoot
Velocidade média download & 35 Mbps \\
Domicílios com internet (dept.) & 45\% \\
Tecnologia predominante & rádio \\
Opção rural & Starlink disponível (\textasciitilde USD 44/\allowbreak mês) \\
\end{longtable}

\subsubsection{Mercado Imobiliário e Terra
Rural}

\textbf{Fonte:} INDERT /\allowbreak  Clasificados.com.py (2024)

\begin{longtable}[]{@{}ll@{}}
\toprule\noalign{}
Tipo & Referência \\
\midrule\noalign{}
\endhead
\bottomrule\noalign{}
\endlastfoot
Terra agrícola alta prod. (USD/\allowbreak ha) & 7,500 \\
Imóvel urbano (USD/\allowbreak m²) & 650 \\
Aluguel 2 quartos (USD/\allowbreak mês) & 260 \\
\end{longtable}

\begin{quote}
Valores de referência departamental. Localidades menores podem ter
preços 20--40\% abaixo da capital departamental. \#\#\# Saúde
\end{quote}

\textbf{Fonte:} MSPBS /\allowbreak  IPS Paraguay (2024-2026), consolidação
departamental e proxy local

\begin{longtable}[]{@{}ll@{}}
\toprule\noalign{}
Serviço & Disponibilidade \\
\midrule\noalign{}
\endhead
\bottomrule\noalign{}
\endlastfoot
USF /\allowbreak  Posto de Saúde & sim \\
Hospital Regional & não \\
IPS (seguro social) & não \\
Farmácia & sim \\
Distância ao hospital de referência & \textasciitilde227 km (Salto del
Guaira) \\
\end{longtable}

\textbf{Principais estabelecimentos:} USF local; referencia hospitalar
em Salto del Guaira

\textbf{Observação para imigrantes:} Atendimento primario local ou em
raio curto; casos de maior complexidade seguem para Salto del Guaira.
Cobertura privada continua recomendada para especialidades e urgências
de maior complexidade.
\subsubsection{Indicadores Sociais}

\textbf{Fontes:} PNUD 2020, INE\footnote{Instituto Nacional de Estadística (Instituto Nacional de Estatística), órgão oficial de dados demográficos e censitários.} Censo 2022, INE\footnote{Instituto Nacional de Estadística (Instituto Nacional de Estatística), órgão oficial de dados demográficos e censitários.} EPHC 2023, Ministerio
Público 2024\\
\textbf{Nota:} Valores marcados como \emph{(dept.)} referem-se ao
departamento de Canindeyú; valores \emph{(dist.)} são específicos deste
distrito. Dados marcados \emph{(est.)} são estimativas.

\begin{longtable}[]{@{}
  >{\raggedright\arraybackslash}p{(\linewidth - 4\tabcolsep) * \real{0.4231}}
  >{\raggedright\arraybackslash}p{(\linewidth - 4\tabcolsep) * \real{0.2692}}
  >{\raggedright\arraybackslash}p{(\linewidth - 4\tabcolsep) * \real{0.3077}}@{}}
\toprule\noalign{}
\begin{minipage}[b]{\linewidth}\raggedright
Indicador
\end{minipage} & \begin{minipage}[b]{\linewidth}\raggedright
Valor
\end{minipage} & \begin{minipage}[b]{\linewidth}\raggedright
Âmbito
\end{minipage} \\
\midrule\noalign{}
\endhead
\bottomrule\noalign{}
\endlastfoot
IDH (2020) & 0,685 & dept. \\
Ranking IDH nacional & 15/\allowbreak 18 & dept. \\
Esperança de vida & 71,2 anos & dept. \\
Escolaridade média & 7,3 anos & dept. \\
RNB per capita (USD PPA) & 7.100 & dept. \\
Pobreza monetária (\%) & 35,1\% & dept. \\
Pobreza extrema (\%) & 10,8\% & dept. \\
Índice de Gini & 0,467 & dept. \\
Acesso a água potável (\%) & 67,4\% & dept. \\
Acesso a saneamento (\%) & 64,8\% & dept. \\
Taxa de homicídios (est., /\allowbreak 100k hab) & \textasciitilde15--25
(estimativa, significativamente acima da média nacional) & dept. \\
Índice de segurança & baixo & dept. \\
População (Censo 2022) & 15.357 habitantes & dist. \\
Idade mediana & N/\allowbreak D & dist. \\
Taxa de fecundidade & N/\allowbreak D & dist. \\
\end{longtable}
\subsubsection{Combustível}

\textbf{Referência:} PETROPAR /\allowbreak  postos locais (2024) \textbf{Tipo de
localidade:} Interior

\begin{longtable}[]{@{}lll@{}}
\toprule\noalign{}
Combustível & USD/\allowbreak litro & Gs/\allowbreak litro (aprox.) \\
\midrule\noalign{}
\endhead
\bottomrule\noalign{}
\endlastfoot
Gasolina 93 oct & 1.00 & 7,400 \\
Gasolina 97 oct (premium) & 1.10 & 8,140 \\
Diesel & 0.93 & 6,882 \\
\end{longtable}

\begin{quote}
Preços podem variar ±5\% conforme posto e sazonalidade. Chaco e interior
remoto apresentam maior variação.
\end{quote}

\subsubsection{Cobertura Celular}

\textbf{Fonte:} CONATEL PY /\allowbreak  operadoras (2024)

\begin{longtable}[]{@{}ll@{}}
\toprule\noalign{}
Parâmetro & Valor \\
\midrule\noalign{}
\endhead
\bottomrule\noalign{}
\endlastfoot
Cobertura 4G população (dept.) & 78\% \\
Cobertura 4G área rural & 55\% \\
Melhor operadora & Tigo \\
Qualidade rural & moderada \\
\end{longtable}

\begin{quote}
Para áreas rurais fora do núcleo urbano, recomenda-se chip Tigo como
principal e Personal como backup.
\end{quote}

\subsubsection{Internet}

\textbf{Fonte:} CONATEL /\allowbreak  Speedtest Ookla (2024)

\begin{longtable}[]{@{}ll@{}}
\toprule\noalign{}
Parâmetro & Valor \\
\midrule\noalign{}
\endhead
\bottomrule\noalign{}
\endlastfoot
Velocidade média download & 35 Mbps \\
Domicílios com internet (dept.) & 45\% \\
Tecnologia predominante & rádio \\
Opção rural & Starlink disponível (\textasciitilde USD 44/\allowbreak mês) \\
\end{longtable}

\subsubsection{Mercado Imobiliário e Terra
Rural}

\textbf{Fonte:} INDERT /\allowbreak  Clasificados.com.py (2024)

\begin{longtable}[]{@{}ll@{}}
\toprule\noalign{}
Tipo & Referência \\
\midrule\noalign{}
\endhead
\bottomrule\noalign{}
\endlastfoot
Terra agrícola alta prod. (USD/\allowbreak ha) & 7,500 \\
Imóvel urbano (USD/\allowbreak m²) & 650 \\
Aluguel 2 quartos (USD/\allowbreak mês) & 260 \\
\end{longtable}

\begin{quote}
Valores de referência departamental. Localidades menores podem ter
preços 20--40\% abaixo da capital departamental. \#\#\# Saúde
\end{quote}

\textbf{Fonte:} MSPBS /\allowbreak  IPS Paraguay (2024-2026), consolidação
departamental e proxy local

\begin{longtable}[]{@{}ll@{}}
\toprule\noalign{}
Serviço & Disponibilidade \\
\midrule\noalign{}
\endhead
\bottomrule\noalign{}
\endlastfoot
USF /\allowbreak  Posto de Saúde & sim \\
Hospital Regional & não \\
IPS (seguro social) & não \\
Farmácia & sim \\
Distância ao hospital de referência & \textasciitilde227 km (Salto del
Guaira) \\
\end{longtable}

\textbf{Principais estabelecimentos:} USF local; referencia hospitalar
em Salto del Guaira

\textbf{Observação para imigrantes:} Atendimento primario local ou em
raio curto; casos de maior complexidade seguem para Salto del Guaira.
Cobertura privada continua recomendada para especialidades e urgências
de maior complexidade.
\newpage
\secaoDiagnostico{Nueva Esperanza}{\mapaConteudo{-24.4333}{-54.8500}}{Nueva Esperanza é o refúgio de elite por excelência no leste paraguaio.
Oferece solos de produtividade mundial, abundância massiva de recursos
alimentares e hídricos, e uma segurança social sustentada por uma
comunidade produtora altamente organizada. Sua força absoluta reside na
resiliência de sua base produtiva e na baixíssima densidade
populacional, garantindo privacidade e isolamento tático superior. É o
local ideal para quem busca autonomia total em ambiente de alta
performance, agindo como uma ``fortaleza alimentar'' capaz de sustentar
seus residentes mesmo em cenários de colapso sistêmico nacional.

\subsubsection{Por que sim}

\begin{itemize}
\tightlist
\item
  Solos de altíssima produtividade mundial e abundância de recursos
  básicos.
\item
  Ambiente social extremamente seguro, ordeiro e resiliente.
\item
  Isolamento demográfico rural de alto nível (privacidade absoluta).
\item
  Infraestrutura de serviços e saúde privada de primeiro nível.
\end{itemize}

\subsubsection{Por que não}

\begin{itemize}
\tightlist
\item
  Alvo estratégico prioritário para controle nacional de suprimentos
  (Fome).
\item
  Valor da terra entre os mais elevados da América do Sul.
\item
  Dependência tática de insumos externos para manutenção da escala
  industrial.
\end{itemize}

\subsubsection{Recomendação
Técnica}

Altamente indicado para estabelecimentos de autossuficiência tecnificada
de alto padrão ou moradia tática. Requer plano de autonomia hídrica e
energética individual para mitigar riscos de rede. O GSS de 8.1 reflete
a excepcional força dos recursos e o isolamento protegida pela
organização social local. Referencia atual de score: GSS 8.1 em
2026-03-05.

\subsubsection{}

\begin{itemize}
\tightlist
\item
  \begin{enumerate}
  \def\labelenumi{\arabic{enumi})}
  \tightlist
  \item
    Dados censitários validados (INE\footnote{Instituto Nacional de Estadística (Instituto Nacional de Estatística), órgão oficial de dados demográficos e censitários.} 2022).
  \end{enumerate}
\item
  \begin{enumerate}
  \def\labelenumi{\arabic{enumi})}
  \setcounter{enumi}{1}
  \tightlist
  \item
    Relatórios de produção agrícola mecanizada de alta precisão (MAG).
  \end{enumerate}
\item
  \begin{enumerate}
  \def\labelenumi{\arabic{enumi})}
  \setcounter{enumi}{2}
  \tightlist
  \item
    Mapa de solos e capacidade de silos de Canindeyú (Portal
    Geoestadístico).
  \end{enumerate}
\item
  \begin{enumerate}
  \def\labelenumi{\arabic{enumi})}
  \setcounter{enumi}{3}
  \tightlist
  \item
    Registro de segurança pública e governança cooperativa regional.
  \end{enumerate}
\end{itemize}}
\begin{table}[ht!]
\small \begin{tabular}{p{3.0cm}p{0.8cm}p{6.0cm}} \toprule
\textbf{Critério} & \textbf{Nota} & \textbf{Justificativa} \\ \midrule
A - Ameaças Estratégicas & 7.5 & Hub agroindustrial de elite; isolamento tático superior fora do eixo urbano, visibilidade como polo de recursos nacional \\
B - Risco Social & 8.5 & Densidade populacional extremamente baixa; isolamento social de alto nível em ambiente rural organizado \\
C - Riscos Naturais & 7.5 & Geologia muito estável e topografia segura contra inundações \\
D - Autossuficiência & 9.5 & Recursos agrícolas e hídricos excepcionais; polo de segurança alimentar nacional \\
E - Institucional & 7.5 & Forte segurança institucional privada e serviços de primeiro nível regional \\
\midrule \textbf{GSS FINAL} & \textbf{8.1} & \textbf{Seguro (Localidade de elite absoluta para autossuficiência e residência tática, oferecendo o máximo em recursos e segurança social).} \\ \bottomrule \end{tabular}
\end{table}
\subsubsection{Vulnerabilidades e
Mitigação}\label{vuln-Nueva_Esperanza-vulnerabilidades-e-mitigauxe7uxe3o}

\begin{itemize}
\tightlist
\item
  \textbf{Visibilidade Tática Crítica:} Sendo um polo de suprimentos, é
  um alvo estratégico óbvio para controle de alimentos em crises
  nacionais (Mitigação: residência rural afastada 10-15 km do eixo
  urbano industrial).
\item
  \textbf{Dependência de Exportação:} Economia vinculada ao mercado
  global (Mitigação: diversificação produtiva para consumo interno e
  autonomia energética solar).
\item
  \textbf{Gargalo Logístico:} Dependência da Ruta PY03 (Mitigação:
  mapeamento de rotas rurais internas em direção a Itakyry e
  Hernandarias).
\end{itemize}
\subsubsection{Dossiê de Campo}\label{dossie-Nueva_Esperanza-dossiuxea-de-campo}
\begin{description}[style=nextline,leftmargin=0.5cm,font=\bfseries]
\item[Coordenadas]  24°26'S 54°51'W.
\item[Topografia]  Relevo predominantemente plano, solo de altíssima produtividade (Terra Roxa mecanizada). Altitude \textasciitilde 240m. Área de \textasciitilde 1.342 km². Geologicamente muito estável, risco sísmico desprezível.
\item[Alvos Estratégicos]  Capital do Agronegócio de Canindeyú (Um dos maiores polos produtores de grãos do Paraguai); Gigantesco complexo de silos e agroindústrias; Eixo logístico da Ruta PY03. Ausência de bases militares massivas, mas centro de gravidade econômica e de suprimentos de importância nacional.
\item[Fallout]  Ventos predominantes do Norte (NE) e Leste. Baixíssimo risco de fallout direto; posição isolada e protegida.
\item[População]  11.069 habitantes (Censo 2022 Final). População com forte herança de imigrantes brasileiros, altamente tecnificada e integrada ao mercado global de commodities.
\item[Densidade]  \textasciitilde 8,2 hab/\allowbreak km² (Densidade extremamente baixa, garantindo privacidade absoluta e ausência de riscos sociais de massa urbana fora do núcleo comercial).
\item[Vias de Acesso]  Localização privilegiada sobre a Ruta PY03. Conectividade asfáltica de elite para Salto del Guairá e Assunção. Facilidade logística tática em tempos de paz, mas centro de interesse para controle de suprimentos em crises.
\item[Serviços]  Infraestrutura de serviços de primeiro nível regional. Sanatórios privados modernos, rede bancária completa e o maior suporte técnico agroindustrial do leste paraguaio.
\item[Custo de Vida]  Médio; economia dolarizada pelo agronegócio e alta valorização das áreas mecanizadas.
\item[Preço da Terra]  US\$ 12.000-18.000/\allowbreak ha (terras mecanizadas de elite mundial).
\item[Hidrologia]  Baixo risco de inundações sistêmicas devido à topografia plana favorável e gestão hídrica eficiente. Presença de arroios perenes de alta qualidade.
\item[Clima]  Subtropical Úmido. Verões intensos e invernos com geadas que beneficiam a produtividade de grãos.
\item[Energia]  Rede nacional (Itaipu); boa estabilidade devido à importância industrial. Nodo de distribuição regional robusto.
\item[Água]  Abundante via poços artesianos profundos e lençol freático produtivo; excelente qualidade hídrica subterrânea (Aquífero Guarani).
\item[Qualidade do Solo]  Latossolo Vermelho de Elite; solo profundo e de fertilidade máxima mundial. Excelente aptidão para soja, milho e trigo.
\item[Resources Locais]  Polo de Segurança Alimentar Nacional. Imensa produção de grãos e proteína animal. Elevadíssimo potencial para autossuficiência absoluta em nível de propriedade ou regional.
\item[Segurança]  Uma das cidades mais seguras e prósperas do país. Baixíssima criminalidade urbana comum. Coesão social fortíssima baseada na cultura cooperativista e produtora. Presença de segurança patrimonial privada de alta performance em todas as instalações agroindustriais.
\item[Leis Local]  Município consolidado e polo de influência ("A Cidade do Trabalho"). \\ \textbf{Livre de Restrição} de Fronteira. \\ \textit{Aquisição de terra rural proibida para estrangeiros limítrofes.}
\end{description}
\paragraph{Solo (SoilGrids 2.0, média ponderada 0--30
cm)}

\begin{longtable}[]{@{}ll@{}}
\toprule\noalign{}
Parâmetro & Valor \\
\midrule\noalign{}
\endhead
\bottomrule\noalign{}
\endlastfoot
pH (H$_2$O) & 5.42 \\
Carbono orgânico (SOC) & 21.38 g/\allowbreak kg \\
Argila & 43.94 \% \\
Areia & 32.01 \% \\
Silte (calc.) & 24.1 \% \\
Densidade aparente & 1.23 g/\allowbreak cm³ \\
\textbf{Aptidão agrícola} & \textbf{Média} \\
\end{longtable}

Fonte: ISRIC SoilGrids 2.0 via WCS (acesso em 2026-03-21). Coords:
-24.4333°, -54.85°. Média ponderada camadas 0-5, 5-15, 15-30 cm.

\begin{itemize}
\tightlist
\item
  \textbf{Homicidios/\allowbreak 100k:} \textasciitilde5,0 (Muito baixo/\allowbreak Produtivo).
\end{itemize}

\subsection{10. Analise de Riscos}

\begin{itemize}
\tightlist
\item
  \textbf{Cenario curto prazo:} Manutenção da prosperidade econômica e
  expansão das agroindústrias.
\item
  \textbf{Cenario medio prazo:} Impacto de mudanças táticas na capital
  sobre o controle das rotas de escoamento de grãos.
\end{itemize}

\subsection{11. Redacao Final}

\subsubsection{Diagnostico Integrado}

Nueva Esperanza é o refúgio de elite por excelência no leste paraguaio.
Oferece solos de produtividade mundial, abundância massiva de recursos
alimentares e hídricos, e uma segurança social sustentada por uma
comunidade produtora altamente organizada. Sua força absoluta reside na
resiliência de sua base produtiva e na baixíssima densidade
populacional, garantindo privacidade e isolamento tático superior. É o
local ideal para quem busca autonomia total em ambiente de alta
performance, agindo como uma ``fortaleza alimentar'' capaz de sustentar
seus residentes mesmo em cenários de colapso sistêmico nacional.

\subsubsection{Por que sim}

\begin{itemize}
\tightlist
\item
  Solos de altíssima produtividade mundial e abundância de recursos
  básicos.
\item
  Ambiente social extremamente seguro, ordeiro e resiliente.
\item
  Isolamento demográfico rural de alto nível (privacidade absoluta).
\item
  Infraestrutura de serviços e saúde privada de primeiro nível.
\end{itemize}

\subsubsection{Por que nao}

\begin{itemize}
\tightlist
\item
  Alvo estratégico prioritário para controle nacional de suprimentos
  (Fome).
\item
  Valor da terra entre os mais elevados da América do Sul.
\item
  Dependência tática de insumos externos para manutenção da escala
  industrial.
\end{itemize}

\subsubsection{Recomendacao Tecnica}

Altamente indicado para estabelecimentos de autossuficiência tecnificada
de alto padrão ou moradia tática. Requer plano de autonomia hídrica e
energética individual para mitigar riscos de rede. O GSS de 8.1 reflete
a excepcional força dos recursos e o isolamento protegida pela
organização social local. Referencia atual de score: GSS 8.1 em
2026-03-05.

\subsection{13.}

\begin{itemize}
\tightlist
\item
  \begin{enumerate}
  \def\labelenumi{\arabic{enumi})}
  \tightlist
  \item
    Dados censitários validados (INE\footnote{Instituto Nacional de Estadística (Instituto Nacional de Estatística), órgão oficial de dados demográficos e censitários.} 2022).
  \end{enumerate}
\item
  \begin{enumerate}
  \def\labelenumi{\arabic{enumi})}
  \setcounter{enumi}{1}
  \tightlist
  \item
    Relatórios de produção agrícola mecanizada de alta precisão (MAG).
  \end{enumerate}
\item
  \begin{enumerate}
  \def\labelenumi{\arabic{enumi})}
  \setcounter{enumi}{2}
  \tightlist
  \item
    Mapa de solos e capacidade de silos de Canindeyú (Portal
    Geoestadístico).
  \end{enumerate}
\item
  \begin{enumerate}
  \def\labelenumi{\arabic{enumi})}
  \setcounter{enumi}{3}
  \tightlist
  \item
    Registro de segurança pública e governança cooperativa regional.
  \end{enumerate}
\end{itemize}

\subsubsection{Combustível}

\textbf{Referência:} PETROPAR /\allowbreak  postos locais (2024) \textbf{Tipo de
localidade:} Interior

\begin{longtable}[]{@{}lll@{}}
\toprule\noalign{}
Combustível & USD/\allowbreak litro & Gs/\allowbreak litro (aprox.) \\
\midrule\noalign{}
\endhead
\bottomrule\noalign{}
\endlastfoot
Gasolina 93 oct & 1.00 & 7,400 \\
Gasolina 97 oct (premium) & 1.10 & 8,140 \\
Diesel & 0.93 & 6,882 \\
\end{longtable}

\begin{quote}
Preços podem variar ±5\% conforme posto e sazonalidade. Chaco e interior
remoto apresentam maior variação.
\end{quote}

\subsubsection{Cobertura Celular}

\textbf{Fonte:} CONATEL PY /\allowbreak  operadoras (2024)

\begin{longtable}[]{@{}ll@{}}
\toprule\noalign{}
Parâmetro & Valor \\
\midrule\noalign{}
\endhead
\bottomrule\noalign{}
\endlastfoot
Cobertura 4G população (dept.) & 78\% \\
Cobertura 4G área rural & 55\% \\
Melhor operadora & Tigo \\
Qualidade rural & moderada \\
\end{longtable}

\begin{quote}
Para áreas rurais fora do núcleo urbano, recomenda-se chip Tigo como
principal e Personal como backup.
\end{quote}

\subsubsection{Internet}

\textbf{Fonte:} CONATEL /\allowbreak  Speedtest Ookla (2024)

\begin{longtable}[]{@{}ll@{}}
\toprule\noalign{}
Parâmetro & Valor \\
\midrule\noalign{}
\endhead
\bottomrule\noalign{}
\endlastfoot
Velocidade média download & 35 Mbps \\
Domicílios com internet (dept.) & 45\% \\
Tecnologia predominante & rádio \\
Opção rural & Starlink disponível (\textasciitilde USD 44/\allowbreak mês) \\
\end{longtable}

\subsubsection{Mercado Imobiliário e Terra
Rural}

\textbf{Fonte:} INDERT /\allowbreak  Clasificados.com.py (2024)

\begin{longtable}[]{@{}ll@{}}
\toprule\noalign{}
Tipo & Referência \\
\midrule\noalign{}
\endhead
\bottomrule\noalign{}
\endlastfoot
Terra agrícola alta prod. (USD/\allowbreak ha) & 7,500 \\
Imóvel urbano (USD/\allowbreak m²) & 650 \\
Aluguel 2 quartos (USD/\allowbreak mês) & 260 \\
\end{longtable}

\begin{quote}
Valores de referência departamental. Localidades menores podem ter
preços 20--40\% abaixo da capital departamental. \#\#\# Saúde
\end{quote}

\textbf{Fonte:} MSPBS /\allowbreak  IPS Paraguay (2024-2026), consolidação
departamental e proxy local

\begin{longtable}[]{@{}ll@{}}
\toprule\noalign{}
Serviço & Disponibilidade \\
\midrule\noalign{}
\endhead
\bottomrule\noalign{}
\endlastfoot
USF /\allowbreak  Posto de Saúde & sim \\
Hospital Regional & não \\
IPS (seguro social) & não \\
Farmácia & sim \\
Distância ao hospital de referência & \textasciitilde91 km (Salto del
Guaira) \\
\end{longtable}

\textbf{Principais estabelecimentos:} USF local; referencia hospitalar
em Salto del Guaira

\textbf{Observação para imigrantes:} Atendimento primario local ou em
raio curto; casos de maior complexidade seguem para Salto del Guaira.
Cobertura privada continua recomendada para especialidades e urgências
de maior complexidade.
\subsubsection{Indicadores Sociais}

\textbf{Fontes:} PNUD 2020, INE\footnote{Instituto Nacional de Estadística (Instituto Nacional de Estatística), órgão oficial de dados demográficos e censitários.} Censo 2022, INE\footnote{Instituto Nacional de Estadística (Instituto Nacional de Estatística), órgão oficial de dados demográficos e censitários.} EPHC 2023, Ministerio
Público 2024\\
\textbf{Nota:} Valores marcados como \emph{(dept.)} referem-se ao
departamento de Canindeyú; valores \emph{(dist.)} são específicos deste
distrito. Dados marcados \emph{(est.)} são estimativas.

\begin{longtable}[]{@{}
  >{\raggedright\arraybackslash}p{(\linewidth - 4\tabcolsep) * \real{0.4231}}
  >{\raggedright\arraybackslash}p{(\linewidth - 4\tabcolsep) * \real{0.2692}}
  >{\raggedright\arraybackslash}p{(\linewidth - 4\tabcolsep) * \real{0.3077}}@{}}
\toprule\noalign{}
\begin{minipage}[b]{\linewidth}\raggedright
Indicador
\end{minipage} & \begin{minipage}[b]{\linewidth}\raggedright
Valor
\end{minipage} & \begin{minipage}[b]{\linewidth}\raggedright
Âmbito
\end{minipage} \\
\midrule\noalign{}
\endhead
\bottomrule\noalign{}
\endlastfoot
IDH (2020) & 0,685 & dept. \\
Ranking IDH nacional & 15/\allowbreak 18 & dept. \\
Esperança de vida & 71,2 anos & dept. \\
Escolaridade média & 7,3 anos & dept. \\
RNB per capita (USD PPA) & 7.100 & dept. \\
Pobreza monetária (\%) & 35,1\% & dept. \\
Pobreza extrema (\%) & 10,8\% & dept. \\
Índice de Gini & 0,467 & dept. \\
Acesso a água potável (\%) & 67,4\% & dept. \\
Acesso a saneamento (\%) & 64,8\% & dept. \\
Taxa de homicídios (est., /\allowbreak 100k hab) & \textasciitilde15--25
(estimativa, significativamente acima da média nacional) & dept. \\
Índice de segurança & baixo & dept. \\
População (Censo 2022) & 11.069 habitantes & dist. \\
Idade mediana & N/\allowbreak D & dist. \\
Taxa de fecundidade & N/\allowbreak D & dist. \\
\end{longtable}
\subsubsection{Combustível}

\textbf{Referência:} PETROPAR /\allowbreak  postos locais (2024) \textbf{Tipo de
localidade:} Interior

\begin{longtable}[]{@{}lll@{}}
\toprule\noalign{}
Combustível & USD/\allowbreak litro & Gs/\allowbreak litro (aprox.) \\
\midrule\noalign{}
\endhead
\bottomrule\noalign{}
\endlastfoot
Gasolina 93 oct & 1.00 & 7,400 \\
Gasolina 97 oct (premium) & 1.10 & 8,140 \\
Diesel & 0.93 & 6,882 \\
\end{longtable}

\begin{quote}
Preços podem variar ±5\% conforme posto e sazonalidade. Chaco e interior
remoto apresentam maior variação.
\end{quote}

\subsubsection{Cobertura Celular}

\textbf{Fonte:} CONATEL PY /\allowbreak  operadoras (2024)

\begin{longtable}[]{@{}ll@{}}
\toprule\noalign{}
Parâmetro & Valor \\
\midrule\noalign{}
\endhead
\bottomrule\noalign{}
\endlastfoot
Cobertura 4G população (dept.) & 78\% \\
Cobertura 4G área rural & 55\% \\
Melhor operadora & Tigo \\
Qualidade rural & moderada \\
\end{longtable}

\begin{quote}
Para áreas rurais fora do núcleo urbano, recomenda-se chip Tigo como
principal e Personal como backup.
\end{quote}

\subsubsection{Internet}

\textbf{Fonte:} CONATEL /\allowbreak  Speedtest Ookla (2024)

\begin{longtable}[]{@{}ll@{}}
\toprule\noalign{}
Parâmetro & Valor \\
\midrule\noalign{}
\endhead
\bottomrule\noalign{}
\endlastfoot
Velocidade média download & 35 Mbps \\
Domicílios com internet (dept.) & 45\% \\
Tecnologia predominante & rádio \\
Opção rural & Starlink disponível (\textasciitilde USD 44/\allowbreak mês) \\
\end{longtable}

\subsubsection{Mercado Imobiliário e Terra
Rural}

\textbf{Fonte:} INDERT /\allowbreak  Clasificados.com.py (2024)

\begin{longtable}[]{@{}ll@{}}
\toprule\noalign{}
Tipo & Referência \\
\midrule\noalign{}
\endhead
\bottomrule\noalign{}
\endlastfoot
Terra agrícola alta prod. (USD/\allowbreak ha) & 7,500 \\
Imóvel urbano (USD/\allowbreak m²) & 650 \\
Aluguel 2 quartos (USD/\allowbreak mês) & 260 \\
\end{longtable}

\begin{quote}
Valores de referência departamental. Localidades menores podem ter
preços 20--40\% abaixo da capital departamental. \#\#\# Saúde
\end{quote}

\textbf{Fonte:} MSPBS /\allowbreak  IPS Paraguay (2024-2026), consolidação
departamental e proxy local

\begin{longtable}[]{@{}ll@{}}
\toprule\noalign{}
Serviço & Disponibilidade \\
\midrule\noalign{}
\endhead
\bottomrule\noalign{}
\endlastfoot
USF /\allowbreak  Posto de Saúde & sim \\
Hospital Regional & não \\
IPS (seguro social) & não \\
Farmácia & sim \\
Distância ao hospital de referência & \textasciitilde91 km (Salto del
Guaira) \\
\end{longtable}

\textbf{Principais estabelecimentos:} USF local; referencia hospitalar
em Salto del Guaira

\textbf{Observação para imigrantes:} Atendimento primario local ou em
raio curto; casos de maior complexidade seguem para Salto del Guaira.
Cobertura privada continua recomendada para especialidades e urgências
de maior complexidade.
\newpage
\secaoDiagnostico{Puerto Adela}{\mapaConteudo{-24.3166}{-54.3500}}{Puerto Adela é um dos enclaves mais discretos e promissores para quem
busca isolamento tático com abundância de recursos hídricos e
produtivos. Oferece um nível de privacidade demográfica quase absoluto,
protegida pela geografia remota e pelo vasto espelho d'água de Itaipu.
Sua força reside no anonimato e na capacidade de subsistência baseada na
proteína animal e solo fértil. O desafio crítico é a logística de acesso
terrestre e a falta de serviços locais, o que exige do ocupante
autonomia total e um perfil ``anfíbio'', capaz de utilizar o rio como
via de suprimento e escape em cenários de instabilidade nacional.

\subsubsection{Por que sim}

\begin{itemize}
\tightlist
\item
  Máximo isolamento geográfico e demográfico (anonimato total).
\item
  Abundância ilimitada de recursos hídricos e pesqueiros.
\item
  Ambiente social pacífico e seguro por isolamento.
\item
  Longe de qualquer alvo militar terrestre ou eixo de fallout tático.
\end{itemize}

\subsubsection{Por que não}

\begin{itemize}
\tightlist
\item
  Risco de isolamento físico por terra durante eventos climáticos
  severos.
\item
  Inexistência de serviços de saúde especializados no distrito.
\item
  Restrições da Lei de Fronteira para investidores estrangeiros.
\end{itemize}

\subsubsection{Recomendação
Técnica}

Altamente indicado para residências de isolamento tático que priorizem a
invisibilidade e possuam infraestrutura própria de transporte fluvial.
Requer autonomia total em energia, saúde e logística de suprimentos. O
GSS de 6.7 reflete a superioridade do isolamento protegida pela barreira
natural da água e da distância. Referencia atual de score: GSS 6.7 em
2026-03-05.

\subsubsection{}

\begin{itemize}
\tightlist
\item
  \begin{enumerate}
  \def\labelenumi{\arabic{enumi})}
  \tightlist
  \item
    Dados censitários validados (INE\footnote{Instituto Nacional de Estadística (Instituto Nacional de Estatística), órgão oficial de dados demográficos e censitários.} 2022).
  \end{enumerate}
\item
  \begin{enumerate}
  \def\labelenumi{\arabic{enumi})}
  \setcounter{enumi}{1}
  \tightlist
  \item
    Levantamento territorial oficial (Portal Geoestadístico).
  \end{enumerate}
\item
  \begin{enumerate}
  \def\labelenumi{\arabic{enumi})}
  \setcounter{enumi}{2}
  \tightlist
  \item
    Mapa de solos e bacias hidrográficas de Canindeyú (MAG).
  \end{enumerate}
\item
  \begin{enumerate}
  \def\labelenumi{\arabic{enumi})}
  \setcounter{enumi}{3}
  \tightlist
  \item
    Registro de segurança pública departamental.
  \end{enumerate}
\end{itemize}}
\begin{table}[ht!]
\small \begin{tabular}{p{3.0cm}p{0.8cm}p{6.0cm}} \toprule
\textbf{Critério} & \textbf{Nota} & \textbf{Justificativa} \\ \midrule
A - Ameaças Estratégicas & 7.0 & Isolamento geográfico e invisibilidade tática superior; ausência de alvos estratégicos primários \\
B - Risco Social & 8.5 & Densidade populacional ínfima; segurança contra riscos de massa absoluta \\
C - Riscos Naturais & 6.5 & Geologia estável; nota penalizada pela vulnerabilidade a isolamento por chuvas \\
D - Autossuficiência & 7.0 & Excelentes recursos hídricos e de proteína; limitado pela inacessibilidade comercial \\
E - Institucional & 4.5 & Suporte institucional em formação; serviços básicos funcionais mas muito limitados \\
\midrule \textbf{GSS FINAL} & \textbf{6.7} & \textbf{Moderadamente Seguro (Ideal para quem busca isolamento tático fluvial e anonimato em terras férteis, aceitando a precariedade logística).} \\ \bottomrule \end{tabular}
\end{table}
\subsubsection{Vulnerabilidades e
Mitigação}\label{vuln-Puerto_Adela-vulnerabilidades-e-mitigauxe7uxe3o}

\begin{itemize}
\tightlist
\item
  \textbf{Isolamento Geográfico Severo:} Risco de ficar ``ilhada'' por
  terra em grandes cheias ou tempestades (Mitigação: posse obrigatória
  de embarcação a motor como rota alternativa pelo rio; manutenção de
  estoque de suprimentos para 90 dias).
\item
  \textbf{Ausência de Saúde Local:} Distância crítica de hospitais
  cirúrgicos (Mitigação: manutenção de estoque médico avançado privado e
  treinamento em primeiros socorros).
\item
  \textbf{Fragilidade Elétrica Rural:} Rede vulnerável (Mitigação:
  instalação de sistemas solares fotovoltaicos potentes off-grid).
\end{itemize}
\subsubsection{Dossiê de Campo}\label{dossie-Puerto_Adela-dossiuxea-de-campo}
\begin{description}[style=nextline,leftmargin=0.5cm,font=\bfseries]
\item[Coordenadas]  24°19'S 54°21'W.
\item[Topografia]  Relevo predominantemente plano a ondulado, situado às margens do reservatório de Itaipu (Rio Paraná). Altitude \textasciitilde 220m. Área de \textasciitilde 481 km². Geologicamente estável, risco sísmico desprezível.
\item[Alvos Estratégicos]  Porto fluvial secundário no Rio Paraná; Áreas de produção agrícola e pecuária em expansão; Proximidade tática com a fronteira fluvial brasileira. Ausência total de infraestrutura industrial massiva ou bases militares de grande porte, oferecendo altíssimo nível de invisibilidade estratégica e proteção por isolamento deliberado.
\item[Fallout]  Ventos predominantes do Norte (NE) e Leste. Baixíssimo risco de fallout direto; posição favorável em relação aos centros de comando nacionais.
\item[População]  1.326 habitantes (Censo 2022 Final). Um dos municípios menos povoados de Canindeyú, com perfil demográfico essencialmente rural.
\item[Densidade]  \textasciitilde 2,8 hab/\allowbreak km² (Densidade extremamente baixa, garantindo privacidade tática absoluta e ausência de riscos sociais de massa urbana).
\item[Vias de Acesso]  Acesso terrestre via ramais de terra a partir de Katueté ou La Paloma (\textasciitilde 50 km). Localização remota com infraestrutura viária limitada, o que favorece o anonimato estratégico, mas exige veículos de tração 4x4 em períodos chuvosos.
\item[Serviços]  Infraestrutura institucional mínima. Unidade de Saúde da Família (USF) local. Dependência absoluta de Salto del Guairá (75 km) ou Katueté para serviços hospitalares especializados e logística comercial avançada.
\item[Custo de Vida]  Muito baixo; economia baseada na produção de grãos, pecuária e pesca comercial.
\item[Preço da Terra]  US\$ 5.000-8.000/\allowbreak ha (terras com aptidão agrícola); US\$ 2.000-4.000/\allowbreak ha (áreas de campo natural/\allowbreak mata).
\item[Hidrologia]  Baixo risco de inundações sistêmicas catastróficas devido à cota elevada do relevo em relação ao reservatório. Presença de diversos arroios tributários da bacia do Rio Paraná. Risco moderado de isolamento temporário por danos em estradas de terra em períodos de chuva intensa (média 1.700 mm/\allowbreak ano).
\item[Clima]  Subtropical Úmido. Verões intensos e chuvosos; invernos amenos.
\item[Energia]  Rede nacional (Itaipu); instável em áreas rurais devido ao isolamento geográfico e falta de redundância.
\item[Água]  Abundância hídrica superficial (reservatório de Itaipu) e poços artesianos profundos de alta qualidade.
\item[Qualidade do Solo]  Solo franco-argiloso e arenoso fértil; aptidão para soja, milho e pastagens naturais.
\item[Resources Locais]  Grande produção de grãos e proteína animal (gado e peixes fluviais). Elevadíssimo potencial para autossuficiência alimentar básica em nível de propriedade devido à vasta área produtiva e baixa população.
\item[Segurança]  Zona rural excepcionalmente pacífica em termos de crime comum. No entanto, exige vigilância tática regional devido à localização em departamento de fronteira e proximidade com o Lago Itaipu (rota de ilícitos). Estabilidade social sustentada por comunidades tradicionais.
\item[Leis Local: Município de criação recente (2018). Restrição de Fronteira (Lei 2532/\allowbreak 05)]  Aplicável a terras rurais na faixa de 50km da fronteira brasileira.
\end{description}
\paragraph{Solo (SoilGrids 2.0, média ponderada 0--30
cm)}

\begin{longtable}[]{@{}ll@{}}
\toprule\noalign{}
Parâmetro & Valor \\
\midrule\noalign{}
\endhead
\bottomrule\noalign{}
\endlastfoot
pH (H$_2$O) & 5.42 \\
Carbono orgânico (SOC) & 21.92 g/\allowbreak kg \\
Argila & 51.86 \% \\
Areia & 27.44 \% \\
Silte (calc.) & 20.7 \% \\
Densidade aparente & 1.23 g/\allowbreak cm³ \\
\textbf{Aptidão agrícola} & \textbf{Média} \\
\end{longtable}

Fonte: ISRIC SoilGrids 2.0 via WCS (acesso em 2026-03-21). Coords:
-24.3167°, -54.35°. Média ponderada camadas 0-5, 5-15, 15-30 cm.
\subsubsection{Indicadores Sociais}

\textbf{Fontes:} PNUD 2020, INE\footnote{Instituto Nacional de Estadística (Instituto Nacional de Estatística), órgão oficial de dados demográficos e censitários.} Censo 2022, INE\footnote{Instituto Nacional de Estadística (Instituto Nacional de Estatística), órgão oficial de dados demográficos e censitários.} EPHC 2023, Ministerio
Público 2024\\
\textbf{Nota:} Valores marcados como \emph{(dept.)} referem-se ao
departamento de Canindeyú; valores \emph{(dist.)} são específicos deste
distrito. Dados marcados \emph{(est.)} são estimativas.

\begin{longtable}[]{@{}
  >{\raggedright\arraybackslash}p{(\linewidth - 4\tabcolsep) * \real{0.4231}}
  >{\raggedright\arraybackslash}p{(\linewidth - 4\tabcolsep) * \real{0.2692}}
  >{\raggedright\arraybackslash}p{(\linewidth - 4\tabcolsep) * \real{0.3077}}@{}}
\toprule\noalign{}
\begin{minipage}[b]{\linewidth}\raggedright
Indicador
\end{minipage} & \begin{minipage}[b]{\linewidth}\raggedright
Valor
\end{minipage} & \begin{minipage}[b]{\linewidth}\raggedright
Âmbito
\end{minipage} \\
\midrule\noalign{}
\endhead
\bottomrule\noalign{}
\endlastfoot
IDH (2020) & 0,685 & dept. \\
Ranking IDH nacional & 15/\allowbreak 18 & dept. \\
Esperança de vida & 71,2 anos & dept. \\
Escolaridade média & 7,3 anos & dept. \\
RNB per capita (USD PPA) & 7.100 & dept. \\
Pobreza monetária (\%) & 35,1\% & dept. \\
Pobreza extrema (\%) & 10,8\% & dept. \\
Índice de Gini & 0,467 & dept. \\
Acesso a água potável (\%) & 67,4\% & dept. \\
Acesso a saneamento (\%) & 64,8\% & dept. \\
Taxa de homicídios (est., /\allowbreak 100k hab) & \textasciitilde15--25
(estimativa, significativamente acima da média nacional) & dept. \\
Índice de segurança & baixo & dept. \\
População (Censo 2022) & 1.326 habitantes & dist. \\
Idade mediana & N/\allowbreak D & dist. \\
Taxa de fecundidade & N/\allowbreak D & dist. \\
\end{longtable}
\subsubsection{Combustível}

\textbf{Referência:} PETROPAR /\allowbreak  postos locais (2024) \textbf{Tipo de
localidade:} Interior

\begin{longtable}[]{@{}lll@{}}
\toprule\noalign{}
Combustível & USD/\allowbreak litro & Gs/\allowbreak litro (aprox.) \\
\midrule\noalign{}
\endhead
\bottomrule\noalign{}
\endlastfoot
Gasolina 93 oct & 1.00 & 7,400 \\
Gasolina 97 oct (premium) & 1.10 & 8,140 \\
Diesel & 0.93 & 6,882 \\
\end{longtable}

\begin{quote}
Preços podem variar ±5\% conforme posto e sazonalidade. Chaco e interior
remoto apresentam maior variação.
\end{quote}

\subsubsection{Cobertura Celular}

\textbf{Fonte:} CONATEL PY /\allowbreak  operadoras (2024)

\begin{longtable}[]{@{}ll@{}}
\toprule\noalign{}
Parâmetro & Valor \\
\midrule\noalign{}
\endhead
\bottomrule\noalign{}
\endlastfoot
Cobertura 4G população (dept.) & 78\% \\
Cobertura 4G área rural & 55\% \\
Melhor operadora & Tigo \\
Qualidade rural & moderada \\
\end{longtable}

\begin{quote}
Para áreas rurais fora do núcleo urbano, recomenda-se chip Tigo como
principal e Personal como backup.
\end{quote}

\subsubsection{Internet}

\textbf{Fonte:} CONATEL /\allowbreak  Speedtest Ookla (2024)

\begin{longtable}[]{@{}ll@{}}
\toprule\noalign{}
Parâmetro & Valor \\
\midrule\noalign{}
\endhead
\bottomrule\noalign{}
\endlastfoot
Velocidade média download & 35 Mbps \\
Domicílios com internet (dept.) & 45\% \\
Tecnologia predominante & rádio \\
Opção rural & Starlink disponível (\textasciitilde USD 44/\allowbreak mês) \\
\end{longtable}

\subsubsection{Mercado Imobiliário e Terra
Rural}

\textbf{Fonte:} INDERT /\allowbreak  Clasificados.com.py (2024)

\begin{longtable}[]{@{}ll@{}}
\toprule\noalign{}
Tipo & Referência \\
\midrule\noalign{}
\endhead
\bottomrule\noalign{}
\endlastfoot
Terra agrícola alta prod. (USD/\allowbreak ha) & 7,500 \\
Imóvel urbano (USD/\allowbreak m²) & 650 \\
Aluguel 2 quartos (USD/\allowbreak mês) & 260 \\
\end{longtable}

\begin{quote}
Valores de referência departamental. Localidades menores podem ter
preços 20--40\% abaixo da capital departamental. \#\#\# Saúde
\end{quote}

\textbf{Fonte:} MSPBS /\allowbreak  IPS Paraguay (2024-2026), consolidação
departamental e proxy local

\begin{longtable}[]{@{}ll@{}}
\toprule\noalign{}
Serviço & Disponibilidade \\
\midrule\noalign{}
\endhead
\bottomrule\noalign{}
\endlastfoot
USF /\allowbreak  Posto de Saúde & sim \\
Hospital Regional & sim \\
IPS (seguro social) & não \\
Farmácia & sim \\
Distância ao hospital de referência & \textasciitilde60 km (Salto del
Guaira) \\
\end{longtable}

\textbf{Principais estabelecimentos:} USF local; referencia hospitalar
em Salto del Guaira

\textbf{Observação para imigrantes:} Atendimento primario local ou em
raio curto; casos de maior complexidade seguem para Salto del Guaira.
Cobertura privada continua recomendada para especialidades e urgências
de maior complexidade.
\newpage
\secaoDiagnostico{Salto del Guaira}{\mapaConteudo{-24.0500}{-54.3000}}{Salto del Guairá é o portal comercial do nordeste paraguaio. Oferece
acesso inigualável a mercadorias, tecnologia e infraestrutura de
serviços de primeiro nível, garantindo uma base de suprimentos
resiliente em tempos de normalidade. Sua força reside na solidez das
instituições capitais e na riqueza de recursos hídricos. Contudo, sua
extrema visibilidade estratégica e a proximidade com a fronteira a
tornam uma localidade de alto risco tático em cenários de conflito
internacional ou instabilidade sistêmica. É o local ideal para bases
administrativas e financeiras, exigindo do ocupante um plano de recuo
para áreas rurais menos expostas em caso de crise.

\subsubsection{Por que sim}

\begin{itemize}
\tightlist
\item
  Melhor infraestrutura comercial e de suprimentos importados do norte
  do país.
\item
  Serviços de saúde regionais de referência.
\item
  Abundância hídrica e energética garantida pela proximidade com Itaipu.
\item
  Suporte institucional forte (Capital).
\end{itemize}

\subsubsection{Por que não}

\begin{itemize}
\tightlist
\item
  Alvo estratégico prioritário em qualquer tensão fronteiriça.
\item
  Insegurança tática elevada (sicariato e crime organizado regional).
\item
  Vulnerabilidade a bloqueios no único eixo asfáltico principal (PY03).
\end{itemize}

\subsubsection{Recomendação
Técnica}

Indicado como hub logístico ou base de suporte comercial. Para moradia,
recomenda-se a zona rural periférica em cotas de segurança, mantendo
autonomia energética e hídrica total. O GSS de 6.3 reflete o poder dos
recursos equilibrado pela vulnerabilidade geopolítica da fronteira.
Referencia atual de score: GSS 6.3 em 2026-03-05.

\subsubsection{}

\begin{itemize}
\tightlist
\item
  \begin{enumerate}
  \def\labelenumi{\arabic{enumi})}
  \tightlist
  \item
    Dados censitários validados (INE\footnote{Instituto Nacional de Estadística (Instituto Nacional de Estatística), órgão oficial de dados demográficos e censitários.} 2022).
  \end{enumerate}
\item
  \begin{enumerate}
  \def\labelenumi{\arabic{enumi})}
  \setcounter{enumi}{1}
  \tightlist
  \item
    Relatórios de segurança pública e inteligência de fronteira (MDI).
  \end{enumerate}
\item
  \begin{enumerate}
  \def\labelenumi{\arabic{enumi})}
  \setcounter{enumi}{2}
  \tightlist
  \item
    Mapa de infraestrutura crítica e ativos comerciais (MIC).
  \end{enumerate}
\item
  \begin{enumerate}
  \def\labelenumi{\arabic{enumi})}
  \setcounter{enumi}{3}
  \tightlist
  \item
    Registro de investimentos imobiliários e desenvolvimento urbano
    capital.
  \end{enumerate}
\end{itemize}}
\begin{table}[ht!]
\small \begin{tabular}{p{3.0cm}p{0.8cm}p{6.0cm}} \toprule
\textbf{Critério} & \textbf{Nota} & \textbf{Justificativa} \\ \midrule
A - Ameaças Estratégicas & 4.5 & Polo estratégico fronteiriço; alvo logístico de alta visibilidade e sensibilidade geopolítica \\
B - Risco Social & 6.0 & Capital departamental com densidade urbana concentrada; risco social em crises de fronteira \\
C - Riscos Naturais & 7.0 & Geologia muito estável e baixo risco de desastres naturais graves \\
D - Autossuficiência & 8.0 & Recursos comerciais massivos, abundância hídrica e logística portuária \\
E - Institucional & 7.0 & Forte segurança institucional e serviços de saúde de primeiro nível regional \\
\midrule \textbf{GSS FINAL} & \textbf{6.3} & \textbf{Moderadamente Seguro (Indicado como hub comercial e logístico, não recomendado para isolamento tático de baixo perfil).} \\ \bottomrule \end{tabular}
\end{table}
\subsubsection{Vulnerabilidades e
Mitigação}\label{vuln-Salto_del_Guaira-vulnerabilidades-e-mitigauxe7uxe3o}

\begin{itemize}
\tightlist
\item
  \textbf{Alvo Geopolítico Sensível:} Localidade de alta tensão em
  conflitos diplomáticos ou militares (Mitigação: residência rural
  afastada 5-10 km do núcleo comercial fronteiriço).
\item
  \textbf{Dependência Econômica Externa:} Economia vinculada à moeda
  brasileira (Mitigação: diversificação de investimentos em áreas
  produtivas rurais ou pecuária).
\item
  \textbf{Bloqueio Logístico:} Único eixo asfáltico principal (PY03)
  pode ser isolado (Mitigação: posse de embarcação para via fluvial e
  mapeamento de rotas de terra via colônias).
\end{itemize}
\subsubsection{Dossiê de Campo}\label{dossie-Salto_del_Guaira-dossiuxea-de-campo}
\begin{description}[style=nextline,leftmargin=0.5cm,font=\bfseries]
\item[Coordenadas]  24°03'S 54°18'W.
\item[Topografia]  Relevo predominantemente plano a suavemente ondulado, margeado pelo Rio Paraná e pelo reservatório de Itaipu ao leste. Altitude \textasciitilde 230m. Área de \textasciitilde 1.283 km². Geologicamente muito estável, risco sísmico desprezível.
\item[Alvos Estratégicos]  Polo de Comércio Internacional (Hub comercial de fronteira seca e fluvial com o Brasil - Mundo Novo e Guaíra); Porto de Salto del Guairá (Logística fluvial); Proximidade estratégica com a Usina Hidrelétrica de Itaipu (\textasciitilde 100 km ao sul). Localização de altíssima sensibilidade geopolítica e vigilância permanente das forças de segurança nacional.
\item[Fallout]  Ventos predominantes do Norte (NE) e Leste. Baixo risco de fallout direto de alvos continentais primários; vulnerabilidade tática em caso de sabotagem infraestrutural regional.
\item[População]  28.553 habitantes (Censo 2022 Final). Capital departamental com forte componente demográfico flutuante devido ao turismo de compras.
\item[Densidade]  \textasciitilde 22,2 hab/\allowbreak km² (Densidade baixa no total distrital, mas com concentração urbana densa no núcleo comercial fronteiriço). Oferece isolamento social em áreas periféricas e chácaras residenciais.
\item[Vias de Acesso]  Ponto final da Ruta PY03 (principal eixo norte-sul). Conectividade asfáltica de elite. Risco extremo de saturação e bloqueio militar de acessos em cenários de instabilidade fronteiriça.
\item[Serviços]  Hospital Regional de Salto del Guairá e rede de sanatórios privados modernos. Polo universitário, financeiro e hoteleiro de primeiro nível regional. Infraestrutura urbana de serviços completa.
\item[Custo de Vida]  Médio-alto; economia dolarizada e real-dependente pelo comércio fronteiriço, com alta valorização imobiliária comercial.
\item[Preço da Terra]  US\$ 5.000-8.000/\allowbreak ha (áreas rurais agrícolas); US\$ 100-250/\allowbreak m² (Urbano/\allowbreak Comercial).
\item[Hidrologia]  Baixo risco de inundações sistêmicas devido à topografia elevada em relação ao reservatório. Vulnerabilidade localizada em drenagem urbana pluvial e áreas ribeirinhas baixas durante tempestades severas (média 1.700 mm/\allowbreak ano).
\item[Clima]  Subtropical Úmido. Verões intensos com ondas de calor significativas.
\item[Energia]  Rede nacional (Itaipu); posição privilegiada próxima à fonte geradora, garantindo prioridade de rede, embora dependente da estabilidade nacional.
\item[Água]  Captação mista via Rio Paraná e poços artesianos profundos (Aquífero Guarani).
\item[Qualidade do Solo]  Solo franco-argiloso fértil; aptidão para agricultura mecanizada e pastagens nas zonas rurais do distrito.
\item[Resources Locais]  Gigantesco hub de suprimentos comerciais, tecnologia e mercadorias importadas. Produção regional de proteína animal e grãos. Elevado potencial para autossuficiência logística regional, mas dependente de logística externa para alimentação massiva da população urbana.
\item[Segurança]  Cidade dinâmica com criminalidade urbana presente (furtos/\allowbreak contrabando). Centro de gravidade para operações de segurança nacional e inteligência de fronteira. Estabilidade institucional alta por ser capital departamental e polo econômico vital.
\item[Leis Local]  Regime administrativo de fronteira. \\ \textbf{Livre de Restrição} de Fronteira para propriedades urbanas centrais, mas restrições vigentes em áreas rurais limítrofes (Lei 2532/\allowbreak 05). \\ \textit{Aquisição de terra rural proibida para estrangeiros limítrofes.}
\end{description}
Coordenadas: -24.06°S /\allowbreak  -54.33°W. Acesso em 2026-03-20.

\textbf{Irradiância solar global --- ALLSKY SFC SW DWN (kWh/\allowbreak m²/\allowbreak dia):}

\begin{longtable}[]{@{}
  >{\raggedright\arraybackslash}p{(\linewidth - 22\tabcolsep) * \real{0.0833}}
  >{\raggedright\arraybackslash}p{(\linewidth - 22\tabcolsep) * \real{0.0833}}
  >{\raggedright\arraybackslash}p{(\linewidth - 22\tabcolsep) * \real{0.0833}}
  >{\raggedright\arraybackslash}p{(\linewidth - 22\tabcolsep) * \real{0.0833}}
  >{\raggedright\arraybackslash}p{(\linewidth - 22\tabcolsep) * \real{0.0833}}
  >{\raggedright\arraybackslash}p{(\linewidth - 22\tabcolsep) * \real{0.0833}}
  >{\raggedright\arraybackslash}p{(\linewidth - 22\tabcolsep) * \real{0.0833}}
  >{\raggedright\arraybackslash}p{(\linewidth - 22\tabcolsep) * \real{0.0833}}
  >{\raggedright\arraybackslash}p{(\linewidth - 22\tabcolsep) * \real{0.0833}}
  >{\raggedright\arraybackslash}p{(\linewidth - 22\tabcolsep) * \real{0.0833}}
  >{\raggedright\arraybackslash}p{(\linewidth - 22\tabcolsep) * \real{0.0833}}
  >{\raggedright\arraybackslash}p{(\linewidth - 22\tabcolsep) * \real{0.0833}}@{}}
\toprule\noalign{}
\begin{minipage}[b]{\linewidth}\raggedright
Jan
\end{minipage} & \begin{minipage}[b]{\linewidth}\raggedright
Fev
\end{minipage} & \begin{minipage}[b]{\linewidth}\raggedright
Mar
\end{minipage} & \begin{minipage}[b]{\linewidth}\raggedright
Abr
\end{minipage} & \begin{minipage}[b]{\linewidth}\raggedright
Mai
\end{minipage} & \begin{minipage}[b]{\linewidth}\raggedright
Jun
\end{minipage} & \begin{minipage}[b]{\linewidth}\raggedright
Jul
\end{minipage} & \begin{minipage}[b]{\linewidth}\raggedright
Ago
\end{minipage} & \begin{minipage}[b]{\linewidth}\raggedright
Set
\end{minipage} & \begin{minipage}[b]{\linewidth}\raggedright
Out
\end{minipage} & \begin{minipage}[b]{\linewidth}\raggedright
Nov
\end{minipage} & \begin{minipage}[b]{\linewidth}\raggedright
Dez
\end{minipage} \\
\midrule\noalign{}
\endhead
\bottomrule\noalign{}
\endlastfoot
6.48 & 5.93 & 5.56 & 4.60 & 3.43 & 3.02 & 3.35 & 4.16 & 4.70 & 5.31 &
6.29 & 6.46 \\
\end{longtable}

Média anual: 4.94 kWh/\allowbreak m²/\allowbreak dia. Mínimo: Jun (3.02) \textbar{} Máximo: Jan
(6.48).

\textbf{Precipitação --- PRECTOTCORR (mm/\allowbreak mês → mm/\allowbreak ano
\textasciitilde1.464):}

\begin{longtable}[]{@{}
  >{\raggedright\arraybackslash}p{(\linewidth - 22\tabcolsep) * \real{0.0833}}
  >{\raggedright\arraybackslash}p{(\linewidth - 22\tabcolsep) * \real{0.0833}}
  >{\raggedright\arraybackslash}p{(\linewidth - 22\tabcolsep) * \real{0.0833}}
  >{\raggedright\arraybackslash}p{(\linewidth - 22\tabcolsep) * \real{0.0833}}
  >{\raggedright\arraybackslash}p{(\linewidth - 22\tabcolsep) * \real{0.0833}}
  >{\raggedright\arraybackslash}p{(\linewidth - 22\tabcolsep) * \real{0.0833}}
  >{\raggedright\arraybackslash}p{(\linewidth - 22\tabcolsep) * \real{0.0833}}
  >{\raggedright\arraybackslash}p{(\linewidth - 22\tabcolsep) * \real{0.0833}}
  >{\raggedright\arraybackslash}p{(\linewidth - 22\tabcolsep) * \real{0.0833}}
  >{\raggedright\arraybackslash}p{(\linewidth - 22\tabcolsep) * \real{0.0833}}
  >{\raggedright\arraybackslash}p{(\linewidth - 22\tabcolsep) * \real{0.0833}}
  >{\raggedright\arraybackslash}p{(\linewidth - 22\tabcolsep) * \real{0.0833}}@{}}
\toprule\noalign{}
\begin{minipage}[b]{\linewidth}\raggedright
Jan
\end{minipage} & \begin{minipage}[b]{\linewidth}\raggedright
Fev
\end{minipage} & \begin{minipage}[b]{\linewidth}\raggedright
Mar
\end{minipage} & \begin{minipage}[b]{\linewidth}\raggedright
Abr
\end{minipage} & \begin{minipage}[b]{\linewidth}\raggedright
Mai
\end{minipage} & \begin{minipage}[b]{\linewidth}\raggedright
Jun
\end{minipage} & \begin{minipage}[b]{\linewidth}\raggedright
Jul
\end{minipage} & \begin{minipage}[b]{\linewidth}\raggedright
Ago
\end{minipage} & \begin{minipage}[b]{\linewidth}\raggedright
Set
\end{minipage} & \begin{minipage}[b]{\linewidth}\raggedright
Out
\end{minipage} & \begin{minipage}[b]{\linewidth}\raggedright
Nov
\end{minipage} & \begin{minipage}[b]{\linewidth}\raggedright
Dez
\end{minipage} \\
\midrule\noalign{}
\endhead
\bottomrule\noalign{}
\endlastfoot
137 & 133 & 102 & 110 & 140 & 82 & 64 & 70 & 93 & 177 & 169 & 188 \\
\end{longtable}

Estação chuvosa Out-Dez; seca Jul-Ago (mínimo em Jul).

\textbf{Inclinação solar recomendada:} 24° Norte (anual); 34° no inverno
(Jun-Ago); 14° no verão (Nov-Jan). Base: latitude -24.06°S.
\paragraph{Solo (SoilGrids 2.0, média ponderada 0--30
cm)}

\begin{longtable}[]{@{}ll@{}}
\toprule\noalign{}
Parâmetro & Valor \\
\midrule\noalign{}
\endhead
\bottomrule\noalign{}
\endlastfoot
pH (H$_2$O) & 5.3 \\
Carbono orgânico (SOC) & 19.87 g/\allowbreak kg \\
Argila & 24.45 \% \\
Areia & 58.32 \% \\
Silte (calc.) & 17.2 \% \\
Densidade aparente & 1.25 g/\allowbreak cm³ \\
\textbf{Aptidão agrícola} & \textbf{Média} \\
\end{longtable}

Fonte: ISRIC SoilGrids 2.0 via WCS (acesso em 2026-03-21). Coords:
-24.05°, -54.3°. Média ponderada camadas 0-5, 5-15, 15-30 cm.
\subsubsection{Indicadores Sociais}

\textbf{Fontes:} PNUD 2020, INE\footnote{Instituto Nacional de Estadística (Instituto Nacional de Estatística), órgão oficial de dados demográficos e censitários.} Censo 2022, INE\footnote{Instituto Nacional de Estadística (Instituto Nacional de Estatística), órgão oficial de dados demográficos e censitários.} EPHC 2023, Ministerio
Público 2024\\
\textbf{Nota:} Valores marcados como \emph{(dept.)} referem-se ao
departamento de Canindeyú; valores \emph{(dist.)} são específicos deste
distrito. Dados marcados \emph{(est.)} são estimativas.

\begin{longtable}[]{@{}
  >{\raggedright\arraybackslash}p{(\linewidth - 4\tabcolsep) * \real{0.4231}}
  >{\raggedright\arraybackslash}p{(\linewidth - 4\tabcolsep) * \real{0.2692}}
  >{\raggedright\arraybackslash}p{(\linewidth - 4\tabcolsep) * \real{0.3077}}@{}}
\toprule\noalign{}
\begin{minipage}[b]{\linewidth}\raggedright
Indicador
\end{minipage} & \begin{minipage}[b]{\linewidth}\raggedright
Valor
\end{minipage} & \begin{minipage}[b]{\linewidth}\raggedright
Âmbito
\end{minipage} \\
\midrule\noalign{}
\endhead
\bottomrule\noalign{}
\endlastfoot
IDH (2020) & 0,685 & dept. \\
Ranking IDH nacional & 15/\allowbreak 18 & dept. \\
Esperança de vida & 71,2 anos & dept. \\
Escolaridade média & 7,3 anos & dept. \\
RNB per capita (USD PPA) & 7.100 & dept. \\
Pobreza monetária (\%) & 35,1\% & dept. \\
Pobreza extrema (\%) & 10,8\% & dept. \\
Índice de Gini & 0,467 & dept. \\
Acesso a água potável (\%) & 67,4\% & dept. \\
Acesso a saneamento (\%) & 64,8\% & dept. \\
Taxa de homicídios (est., /\allowbreak 100k hab) & \textasciitilde15--25
(estimativa, significativamente acima da média nacional) & dept. \\
Índice de segurança & baixo & dept. \\
População (Censo 2022) & 28.553 habitantes & dist. \\
Idade mediana & N/\allowbreak D & dist. \\
Taxa de fecundidade & N/\allowbreak D & dist. \\
\end{longtable}
\subsubsection{Combustível}

\textbf{Referência:} PETROPAR /\allowbreak  postos locais (2024) \textbf{Tipo de
localidade:} Capital departamental

\begin{longtable}[]{@{}lll@{}}
\toprule\noalign{}
Combustível & USD/\allowbreak litro & Gs/\allowbreak litro (aprox.) \\
\midrule\noalign{}
\endhead
\bottomrule\noalign{}
\endlastfoot
Gasolina 93 oct & 1.00 & 7,400 \\
Gasolina 97 oct (premium) & 1.10 & 8,140 \\
Diesel & 0.93 & 6,882 \\
\end{longtable}

\begin{quote}
Preços podem variar ±5\% conforme posto e sazonalidade. Chaco e interior
remoto apresentam maior variação.
\end{quote}

\subsubsection{Cobertura Celular}

\textbf{Fonte:} CONATEL PY /\allowbreak  operadoras (2024)

\begin{longtable}[]{@{}ll@{}}
\toprule\noalign{}
Parâmetro & Valor \\
\midrule\noalign{}
\endhead
\bottomrule\noalign{}
\endlastfoot
Cobertura 4G população (dept.) & 78\% \\
Cobertura 4G área rural & 55\% \\
Melhor operadora & Tigo \\
Qualidade rural & moderada \\
\end{longtable}

\begin{quote}
Para áreas rurais fora do núcleo urbano, recomenda-se chip Tigo como
principal e Personal como backup.
\end{quote}

\subsubsection{Internet}

\textbf{Fonte:} CONATEL /\allowbreak  Speedtest Ookla (2024)

\begin{longtable}[]{@{}ll@{}}
\toprule\noalign{}
Parâmetro & Valor \\
\midrule\noalign{}
\endhead
\bottomrule\noalign{}
\endlastfoot
Velocidade média download & 35 Mbps \\
Domicílios com internet (dept.) & 45\% \\
Tecnologia predominante & rádio \\
Opção rural & Starlink disponível (\textasciitilde USD 44/\allowbreak mês) \\
\end{longtable}

\subsubsection{Mercado Imobiliário e Terra
Rural}

\textbf{Fonte:} INDERT /\allowbreak  Clasificados.com.py (2024)

\begin{longtable}[]{@{}ll@{}}
\toprule\noalign{}
Tipo & Referência \\
\midrule\noalign{}
\endhead
\bottomrule\noalign{}
\endlastfoot
Terra agrícola alta prod. (USD/\allowbreak ha) & 7,500 \\
Imóvel urbano (USD/\allowbreak m²) & 747 \\
Aluguel 2 quartos (USD/\allowbreak mês) & 312 \\
\end{longtable}

\begin{quote}
Valores de referência departamental. Localidades menores podem ter
preços 20--40\% abaixo da capital departamental. \#\#\# Saúde
\end{quote}

\textbf{Fonte:} MSPBS /\allowbreak  IPS Paraguay (2024-2026), consolidação
departamental e proxy local

\begin{longtable}[]{@{}ll@{}}
\toprule\noalign{}
Serviço & Disponibilidade \\
\midrule\noalign{}
\endhead
\bottomrule\noalign{}
\endlastfoot
USF /\allowbreak  Posto de Saúde & sim \\
Hospital Regional & sim \\
IPS (seguro social) & sim \\
Farmácia & sim \\
Distância ao hospital de referência & local \\
\end{longtable}

\textbf{Principais estabelecimentos:} Hospital distrital /\allowbreak  regional de
Salto del Guaira; rede de USF e farmacias locais

\textbf{Observação para imigrantes:} Centro de referencia do
departamento. Acesso a atendimento primario e, em geral, a maior oferta
publica e privada da area. Cobertura privada continua recomendada para
especialidades e urgências de maior complexidade.
\newpage
\secaoDiagnostico{Villa Ygatimi}{\mapaConteudo{-24.1333}{-55.5000}}{Villa Ygatimí é um refúgio produtivo de baixo perfil situado em uma das
zonas mais ricas e isoladas de Canindeyú. Oferece solos férteis e um
isolamento demográfico que garante privacidade absoluta, protegida pela
proximidade com a Reserva Mbaracayú. Sua força reside na abundância de
recursos naturais e na capacidade de subsistência baseada na proteína
animal. O contexto regional exige que o ocupante adote um perfil
vigilante e autônomo, transformando a vastidão do território em uma
barreira de proteção tática contra instabilidades externas.

\subsubsection{Por que sim}

\begin{itemize}
\tightlist
\item
  Máximo isolamento demográfico em território vasto (privacidade).
\item
  Abundância de águas puras e solos aptos para pecuária e agricultura.
\item
  Conectividade asfáltica eficiente via Ruta PY13 para Curuguaty.
\item
  Sem restrições da Lei de Fronteira.
\end{itemize}

\subsubsection{Por que não}

\begin{itemize}
\tightlist
\item
  Insegurança tática regional latente no departamento de Canindeyú.
\item
  Dependência de centros maiores para serviços de saúde complexos.
\item
  Pequena escala comercial local exige planejamento logístico de
  suprimentos.
\end{itemize}

\subsubsection{Recomendação
Técnica}

Altamente indicado para estabelecimentos de autossuficiência que
valorizem o isolamento natural e produtividade. Requer investimento em
segurança patrimonial e autonomia em saúde. O GSS de 7.0 reflete a
solidez dos recursos e o isolamento em equilíbrio com os desafios do
contexto regional. Referencia atual de score: GSS 7.0 em 2026-03-05.

\subsubsection{}

\begin{itemize}
\tightlist
\item
  \begin{enumerate}
  \def\labelenumi{\arabic{enumi})}
  \tightlist
  \item
    Dados censitários validados (INE\footnote{Instituto Nacional de Estadística (Instituto Nacional de Estatística), órgão oficial de dados demográficos e censitários.} 2022).
  \end{enumerate}
\item
  \begin{enumerate}
  \def\labelenumi{\arabic{enumi})}
  \setcounter{enumi}{1}
  \tightlist
  \item
    Relatórios de monitoramento ambiental regional (MADES).
  \end{enumerate}
\item
  \begin{enumerate}
  \def\labelenumi{\arabic{enumi})}
  \setcounter{enumi}{2}
  \tightlist
  \item
    Mapa de solos e hidrografia de Canindeyú (MAG).
  \end{enumerate}
\item
  \begin{enumerate}
  \def\labelenumi{\arabic{enumi})}
  \setcounter{enumi}{3}
  \tightlist
  \item
    Registro de segurança pública departamental.
  \end{enumerate}
\end{itemize}}
\begin{table}[ht!]
\small \begin{tabular}{p{3.0cm}p{0.8cm}p{6.0cm}} \toprule
\textbf{Critério} & \textbf{Nota} & \textbf{Justificativa} \\ \midrule
A - Ameaças Estratégicas & 6.5 & Localização remota e discreta; boa invisibilidade tática, equilibrada por riscos regionais \\
B - Risco Social & 8.0 & Densidade populacional extremamente baixa; isolamento social superior \\
C - Riscos Naturais & 7.0 & Geologia muito estável e baixo risco de grandes desastres naturais \\
D - Autossuficiência & 7.5 & Excelentes recursos hídricos e de proteína animal; solo fértil em expansão \\
E - Institucional & 6.0 & Suporte institucional funcional, mas dependente de cidades vizinhas para alta complexidade \\
\midrule \textbf{GSS FINAL} & \textbf{7.0} & \textbf{Seguro (Altamente recomendado para quem busca isolamento tático e autossuficiência em solo fértil, ciente dos desafios regionais).} \\ \bottomrule \end{tabular}
\end{table}
\subsubsection{Vulnerabilidades e
Mitigação}\label{vuln-Villa_Ygatimi-vulnerabilidades-e-mitigauxe7uxe3o}

\begin{itemize}
\tightlist
\item
  \textbf{Vulnerabilidade Tática Regional:} Insegurança ligada ao
  narcotráfico no departamento (Mitigação: residência rural protegida e
  investimento em sistemas de segurança patrimonial robustos).
\item
  \textbf{Isolamento de Saúde:} Necessidade de deslocamento para
  Curuguaty em emergências (Mitigação: manutenção de estoque médico
  básico e kit trauma local).
\item
  \textbf{Pressão Ambiental:} Vigilância tática necessária em áreas
  limítrofes à Reserva Mbaracayú para evitar invasões (Mitigação:
  aquisição de propriedades com documentação legal perfeita).
\end{itemize}
\subsubsection{Dossiê de Campo}\label{dossie-Villa_Ygatimi-dossiuxea-de-campo}
\begin{description}[style=nextline,leftmargin=0.5cm,font=\bfseries]
\item[Coordenadas]  24°08'S 55°30'W.
\item[Topografia]  Relevo ondulado com extensas planícies mecanizáveis e áreas de mata protegida. Situado nas proximidades da Reserva Mbaracayú. Altitude \textasciitilde 230m. Área vasta de \textasciitilde 1.898 km². Geologicamente muito estável, risco sísmico desprezível.
\item[Alvos Estratégicos]  Áreas de produção agrícola e pecuária em expansão; Ponto logístico secundário entre Curuguaty e Ypejhú; Proximidade estratégica com a Reserva Natural del Bosque Mbaracayú. Ausência de bases militares massivas, mas zona de vigilância ambiental e tática regional permanente.
\item[Fallout]  Ventos predominantes do Norte (NE) e Leste. Baixíssimo risco de fallout direto de alvos continentais primários devido ao isolamento geográfico.
\item[População]  13.074 habitantes (Censo 2022 Final). População estável com forte tradição rural e presença de comunidades indígenas originárias.
\item[Densidade]  \textasciitilde 6,9 hab/\allowbreak km² (Densidade extremamente baixa, garantindo o máximo isolamento tático e privacidade absoluta fora da vila urbana).
\item[Vias de Acesso]  Localização privilegiada sobre a Ruta PY13 (pavimentada). Conectividade eficiente para Curuguaty (approx. 40 km). Malha interna de estradas de terra que exige veículos robustos em períodos chuvosos (média 1.700 mm/\allowbreak ano).
\item[Serviços]  Infraestrutura básica funcional. Centro de Saúde local e USF. Dependência direta de Curuguaty para serviços hospitalares de alta complexidade e logística comercial especializada.
\item[Custo de Vida]  Muito baixo; economia baseada na pecuária, grãos e exploração florestal sustentável.
\item[Preço da Terra]  US\$ 6.000-9.000/\allowbreak ha (terras desenvolvidas); US\$ 2.000-4.000/\allowbreak ha (áreas de uso misto/\allowbreak monte).
\item[Hidrologia]  Baixo risco de inundações sistêmicas devido ao relevo ondulado e preservação florestal. Diversos arroios perenes de alta qualidade hídrica (ex: Rio Jejuí-mí) garantem suprimento hídrico superficial.
\item[Clima]  Subtropical Úmido. Verões intensos e chuvosos; invernos amenos com geadas ocasionais.
\item[Energia]  Rede nacional (Itaipu); instável em áreas rurais remotas.
\item[Água]  Abundante via nascentes naturais e poços artesianos; excelente qualidade hídrica subterrânea e superficial.
\item[Qualidade do Solo]  Solo franco-argiloso fértil; aptidão para soja, milho, pastagens e silvicultura.
\item[Resources Locais]  Grande produção de gado bovino e grãos. Elevadíssimo potencial para autossuficiência alimentar básica em nível de propriedade devido à vasta área produtiva e baixa população.
\item[Segurança]  Zona rural pacífica em termos de crime comum. No entanto, exige atenção tática regional devido ao contexto de Canindeyú (histórico de tensões agrárias e proximidade com rotas de narcotráfico). Estabilidade social baseada na cultura agropecuária tradicional.
\item[Leis Local: Município histórico e consolidado. Livre de Restrição de Fronteira]  Localizado fora da faixa de 50km da fronteira seca internacional.
\end{description}
\paragraph{Solo (SoilGrids 2.0, média ponderada 0--30
cm)}

\begin{longtable}[]{@{}ll@{}}
\toprule\noalign{}
Parâmetro & Valor \\
\midrule\noalign{}
\endhead
\bottomrule\noalign{}
\endlastfoot
pH (H$_2$O) & 5.3 \\
Carbono orgânico (SOC) & 21.23 g/\allowbreak kg \\
Argila & 26.25 \% \\
Areia & 55.47 \% \\
Silte (calc.) & 18.3 \% \\
Densidade aparente & 1.25 g/\allowbreak cm³ \\
\textbf{Aptidão agrícola} & \textbf{Média} \\
\end{longtable}

Fonte: ISRIC SoilGrids 2.0 via WCS (acesso em 2026-03-21). Coords:
-24.1333°, -55.5°. Média ponderada camadas 0-5, 5-15, 15-30 cm.

\begin{itemize}
\tightlist
\item
  \textbf{Homicidios/\allowbreak 100k:} \textasciitilde18,0 (Regional/\allowbreak Canindeyú).
\end{itemize}

\subsection{10. Analise de Riscos}

\begin{itemize}
\tightlist
\item
  \textbf{Cenario curto prazo:} Manutenção da tranquilidade rural e
  controle de segurança regional.
\item
  \textbf{Cenario medio prazo:} Valorização das terras pela consolidação
  da infraestrutura asfáltica e expansão da pecuária tecnificada.
\end{itemize}

\subsection{11. Redacao Final}

\subsubsection{Diagnostico Integrado}

Villa Ygatimí é um refúgio produtivo de baixo perfil situado em uma das
zonas mais ricas e isoladas de Canindeyú. Oferece solos férteis e um
isolamento demográfico que garante privacidade absoluta, protegida pela
proximidade com a Reserva Mbaracayú. Sua força reside na abundância de
recursos naturais e na capacidade de subsistência baseada na proteína
animal. O contexto regional exige que o ocupante adote um perfil
vigilante e autônomo, transformando a vastidão do território em uma
barreira de proteção tática contra instabilidades externas.

\subsubsection{Por que sim}

\begin{itemize}
\tightlist
\item
  Máximo isolamento demográfico em território vasto (privacidade).
\item
  Abundância de águas puras e solos aptos para pecuária e agricultura.
\item
  Conectividade asfáltica eficiente via Ruta PY13 para Curuguaty.
\item
  Sem restrições da Lei de Fronteira.
\end{itemize}

\subsubsection{Por que nao}

\begin{itemize}
\tightlist
\item
  Insegurança tática regional latente no departamento de Canindeyú.
\item
  Dependência de centros maiores para serviços de saúde complexos.
\item
  Pequena escala comercial local exige planejamento logístico de
  suprimentos.
\end{itemize}

\subsubsection{Recomendacao Tecnica}

Altamente indicado para estabelecimentos de autossuficiência que
valorizem o isolamento natural e produtividade. Requer investimento em
segurança patrimonial e autonomia em saúde. O GSS de 7.0 reflete a
solidez dos recursos e o isolamento em equilíbrio com os desafios do
contexto regional. Referencia atual de score: GSS 7.0 em 2026-03-05.

\subsection{13.}

\begin{itemize}
\tightlist
\item
  \begin{enumerate}
  \def\labelenumi{\arabic{enumi})}
  \tightlist
  \item
    Dados censitários validados (INE\footnote{Instituto Nacional de Estadística (Instituto Nacional de Estatística), órgão oficial de dados demográficos e censitários.} 2022).
  \end{enumerate}
\item
  \begin{enumerate}
  \def\labelenumi{\arabic{enumi})}
  \setcounter{enumi}{1}
  \tightlist
  \item
    Relatórios de monitoramento ambiental regional (MADES).
  \end{enumerate}
\item
  \begin{enumerate}
  \def\labelenumi{\arabic{enumi})}
  \setcounter{enumi}{2}
  \tightlist
  \item
    Mapa de solos e hidrografia de Canindeyú (MAG).
  \end{enumerate}
\item
  \begin{enumerate}
  \def\labelenumi{\arabic{enumi})}
  \setcounter{enumi}{3}
  \tightlist
  \item
    Registro de segurança pública departamental.
  \end{enumerate}
\end{itemize}

\subsubsection{Combustível}

\textbf{Referência:} PETROPAR /\allowbreak  postos locais (2024) \textbf{Tipo de
localidade:} Interior

\begin{longtable}[]{@{}lll@{}}
\toprule\noalign{}
Combustível & USD/\allowbreak litro & Gs/\allowbreak litro (aprox.) \\
\midrule\noalign{}
\endhead
\bottomrule\noalign{}
\endlastfoot
Gasolina 93 oct & 1.00 & 7,400 \\
Gasolina 97 oct (premium) & 1.10 & 8,140 \\
Diesel & 0.93 & 6,882 \\
\end{longtable}

\begin{quote}
Preços podem variar ±5\% conforme posto e sazonalidade. Chaco e interior
remoto apresentam maior variação.
\end{quote}

\subsubsection{Cobertura Celular}

\textbf{Fonte:} CONATEL PY /\allowbreak  operadoras (2024)

\begin{longtable}[]{@{}ll@{}}
\toprule\noalign{}
Parâmetro & Valor \\
\midrule\noalign{}
\endhead
\bottomrule\noalign{}
\endlastfoot
Cobertura 4G população (dept.) & 78\% \\
Cobertura 4G área rural & 55\% \\
Melhor operadora & Tigo \\
Qualidade rural & moderada \\
\end{longtable}

\begin{quote}
Para áreas rurais fora do núcleo urbano, recomenda-se chip Tigo como
principal e Personal como backup.
\end{quote}

\subsubsection{Internet}

\textbf{Fonte:} CONATEL /\allowbreak  Speedtest Ookla (2024)

\begin{longtable}[]{@{}ll@{}}
\toprule\noalign{}
Parâmetro & Valor \\
\midrule\noalign{}
\endhead
\bottomrule\noalign{}
\endlastfoot
Velocidade média download & 35 Mbps \\
Domicílios com internet (dept.) & 45\% \\
Tecnologia predominante & rádio \\
Opção rural & Starlink disponível (\textasciitilde USD 44/\allowbreak mês) \\
\end{longtable}

\subsubsection{Mercado Imobiliário e Terra
Rural}

\textbf{Fonte:} INDERT /\allowbreak  Clasificados.com.py (2024)

\begin{longtable}[]{@{}ll@{}}
\toprule\noalign{}
Tipo & Referência \\
\midrule\noalign{}
\endhead
\bottomrule\noalign{}
\endlastfoot
Terra agrícola alta prod. (USD/\allowbreak ha) & 7,500 \\
Imóvel urbano (USD/\allowbreak m²) & 650 \\
Aluguel 2 quartos (USD/\allowbreak mês) & 260 \\
\end{longtable}

\begin{quote}
Valores de referência departamental. Localidades menores podem ter
preços 20--40\% abaixo da capital departamental. \#\#\# Saúde
\end{quote}

\textbf{Fonte:} MSPBS /\allowbreak  IPS Paraguay (2024-2026), consolidação
departamental e proxy local

\begin{longtable}[]{@{}ll@{}}
\toprule\noalign{}
Serviço & Disponibilidade \\
\midrule\noalign{}
\endhead
\bottomrule\noalign{}
\endlastfoot
USF /\allowbreak  Posto de Saúde & sim \\
Hospital Regional & sim \\
IPS (seguro social) & não \\
Farmácia & sim \\
Distância ao hospital de referência & \textasciitilde167 km (Salto del
Guaira) \\
\end{longtable}

\textbf{Principais estabelecimentos:} USF local; referencia hospitalar
em Salto del Guaira

\textbf{Observação para imigrantes:} Atendimento primario local ou em
raio curto; casos de maior complexidade seguem para Salto del Guaira.
Cobertura privada continua recomendada para especialidades e urgências
de maior complexidade.
\subsubsection{Indicadores Sociais}

\textbf{Fontes:} PNUD 2020, INE\footnote{Instituto Nacional de Estadística (Instituto Nacional de Estatística), órgão oficial de dados demográficos e censitários.} Censo 2022, INE\footnote{Instituto Nacional de Estadística (Instituto Nacional de Estatística), órgão oficial de dados demográficos e censitários.} EPHC 2023, Ministerio
Público 2024\\
\textbf{Nota:} Valores marcados como \emph{(dept.)} referem-se ao
departamento de Canindeyú; valores \emph{(dist.)} são específicos deste
distrito. Dados marcados \emph{(est.)} são estimativas.

\begin{longtable}[]{@{}
  >{\raggedright\arraybackslash}p{(\linewidth - 4\tabcolsep) * \real{0.4231}}
  >{\raggedright\arraybackslash}p{(\linewidth - 4\tabcolsep) * \real{0.2692}}
  >{\raggedright\arraybackslash}p{(\linewidth - 4\tabcolsep) * \real{0.3077}}@{}}
\toprule\noalign{}
\begin{minipage}[b]{\linewidth}\raggedright
Indicador
\end{minipage} & \begin{minipage}[b]{\linewidth}\raggedright
Valor
\end{minipage} & \begin{minipage}[b]{\linewidth}\raggedright
Âmbito
\end{minipage} \\
\midrule\noalign{}
\endhead
\bottomrule\noalign{}
\endlastfoot
IDH (2020) & 0,685 & dept. \\
Ranking IDH nacional & 15/\allowbreak 18 & dept. \\
Esperança de vida & 71,2 anos & dept. \\
Escolaridade média & 7,3 anos & dept. \\
RNB per capita (USD PPA) & 7.100 & dept. \\
Pobreza monetária (\%) & 35,1\% & dept. \\
Pobreza extrema (\%) & 10,8\% & dept. \\
Índice de Gini & 0,467 & dept. \\
Acesso a água potável (\%) & 67,4\% & dept. \\
Acesso a saneamento (\%) & 64,8\% & dept. \\
Taxa de homicídios (est., /\allowbreak 100k hab) & \textasciitilde15--25
(estimativa, significativamente acima da média nacional) & dept. \\
Índice de segurança & baixo & dept. \\
População (Censo 2022) & 13.074 habitantes & dist. \\
Idade mediana & N/\allowbreak D & dist. \\
Taxa de fecundidade & N/\allowbreak D & dist. \\
\end{longtable}
\subsubsection{Combustível}

\textbf{Referência:} PETROPAR /\allowbreak  postos locais (2024) \textbf{Tipo de
localidade:} Interior

\begin{longtable}[]{@{}lll@{}}
\toprule\noalign{}
Combustível & USD/\allowbreak litro & Gs/\allowbreak litro (aprox.) \\
\midrule\noalign{}
\endhead
\bottomrule\noalign{}
\endlastfoot
Gasolina 93 oct & 1.00 & 7,400 \\
Gasolina 97 oct (premium) & 1.10 & 8,140 \\
Diesel & 0.93 & 6,882 \\
\end{longtable}

\begin{quote}
Preços podem variar ±5\% conforme posto e sazonalidade. Chaco e interior
remoto apresentam maior variação.
\end{quote}

\subsubsection{Cobertura Celular}

\textbf{Fonte:} CONATEL PY /\allowbreak  operadoras (2024)

\begin{longtable}[]{@{}ll@{}}
\toprule\noalign{}
Parâmetro & Valor \\
\midrule\noalign{}
\endhead
\bottomrule\noalign{}
\endlastfoot
Cobertura 4G população (dept.) & 78\% \\
Cobertura 4G área rural & 55\% \\
Melhor operadora & Tigo \\
Qualidade rural & moderada \\
\end{longtable}

\begin{quote}
Para áreas rurais fora do núcleo urbano, recomenda-se chip Tigo como
principal e Personal como backup.
\end{quote}

\subsubsection{Internet}

\textbf{Fonte:} CONATEL /\allowbreak  Speedtest Ookla (2024)

\begin{longtable}[]{@{}ll@{}}
\toprule\noalign{}
Parâmetro & Valor \\
\midrule\noalign{}
\endhead
\bottomrule\noalign{}
\endlastfoot
Velocidade média download & 35 Mbps \\
Domicílios com internet (dept.) & 45\% \\
Tecnologia predominante & rádio \\
Opção rural & Starlink disponível (\textasciitilde USD 44/\allowbreak mês) \\
\end{longtable}

\subsubsection{Mercado Imobiliário e Terra
Rural}

\textbf{Fonte:} INDERT /\allowbreak  Clasificados.com.py (2024)

\begin{longtable}[]{@{}ll@{}}
\toprule\noalign{}
Tipo & Referência \\
\midrule\noalign{}
\endhead
\bottomrule\noalign{}
\endlastfoot
Terra agrícola alta prod. (USD/\allowbreak ha) & 7,500 \\
Imóvel urbano (USD/\allowbreak m²) & 650 \\
Aluguel 2 quartos (USD/\allowbreak mês) & 260 \\
\end{longtable}

\begin{quote}
Valores de referência departamental. Localidades menores podem ter
preços 20--40\% abaixo da capital departamental. \#\#\# Saúde
\end{quote}

\textbf{Fonte:} MSPBS /\allowbreak  IPS Paraguay (2024-2026), consolidação
departamental e proxy local

\begin{longtable}[]{@{}ll@{}}
\toprule\noalign{}
Serviço & Disponibilidade \\
\midrule\noalign{}
\endhead
\bottomrule\noalign{}
\endlastfoot
USF /\allowbreak  Posto de Saúde & sim \\
Hospital Regional & sim \\
IPS (seguro social) & não \\
Farmácia & sim \\
Distância ao hospital de referência & \textasciitilde167 km (Salto del
Guaira) \\
\end{longtable}

\textbf{Principais estabelecimentos:} USF local; referencia hospitalar
em Salto del Guaira

\textbf{Observação para imigrantes:} Atendimento primario local ou em
raio curto; casos de maior complexidade seguem para Salto del Guaira.
Cobertura privada continua recomendada para especialidades e urgências
de maior complexidade.
\newpage
\secaoDiagnostico{Yasy Cany}{\mapaConteudo{-24.3333}{-55.4166}}{Yasy Cañy é um dos polos de resiliência e produtividade mais
equilibrados de Canindeyú. Oferece solos de elite mundial e uma
abundância massiva de recursos naturais, funcionando como um enclave de
segurança alimentar e hídrica. Sua força reside na excepcional qualidade
da base produtiva combinada com um suporte institucional sólido para o
contexto regional. É o local ideal para quem busca autonomia total em
ambiente de alta performance, transformando a fertilidade do solo em uma
barreira definitiva de proteção econômica e tática contra crises
externas.

\subsubsection{Por que sim}

\begin{itemize}
\tightlist
\item
  Solos de altíssima produtividade mundial e abundância de águas puras.
\item
  Ambiente social seguro e resiliente sustentado pela produção.
\item
  Conectividade asfáltica de elite via Ruta PY03 para Curuguaty.
\item
  Sem restrições da Lei de Fronteira.
\end{itemize}

\subsubsection{Por que não}

\begin{itemize}
\tightlist
\item
  Insegurança tática regional latente no departamento de Canindeyú.
\item
  Alvo estratégico logístico em cenários de conflito por suprimentos.
\item
  Valor da terra em processo de valorização acelerada.
\end{itemize}

\subsubsection{Recomendação
Técnica}

Altamente indicado para estabelecimentos de autossuficiência tecnificada
ou investimentos produtivos de larga escala. Requer autonomia energética
e médica complementar. O GSS de 7.4 reflete a excepcional força dos
recursos e da governança local em equilíbrio com a visibilidade
estratégica. Referencia atual de score: GSS 7.4 em 2026-03-05.

\subsubsection{}

\begin{itemize}
\tightlist
\item
  \begin{enumerate}
  \def\labelenumi{\arabic{enumi})}
  \tightlist
  \item
    Dados censitários validados (INE\footnote{Instituto Nacional de Estadística (Instituto Nacional de Estatística), órgão oficial de dados demográficos e censitários.} 2022).
  \end{enumerate}
\item
  \begin{enumerate}
  \def\labelenumi{\arabic{enumi})}
  \setcounter{enumi}{1}
  \tightlist
  \item
    Relatórios de produção agrícola mecanizada de alta precisão (MAG).
  \end{enumerate}
\item
  \begin{enumerate}
  \def\labelenumi{\arabic{enumi})}
  \setcounter{enumi}{2}
  \tightlist
  \item
    Mapa de solos e bacias hidrográficas de Canindeyú (Portal
    Geoestadístico).
  \end{enumerate}
\item
  \begin{enumerate}
  \def\labelenumi{\arabic{enumi})}
  \setcounter{enumi}{3}
  \tightlist
  \item
    Registro de segurança pública departamental.
  \end{enumerate}
\end{itemize}}
\begin{table}[ht!]
\small \begin{tabular}{p{3.0cm}p{0.8cm}p{6.0cm}} \toprule
\textbf{Critério} & \textbf{Nota} & \textbf{Justificativa} \\ \midrule
A - Ameaças Estratégicas & 6.5 & Hub produtivo estratégico; boa invisibilidade tática fora do eixo urbano, visibilidade moderada como alvo logístico \\
B - Risco Social & 7.5 & Densidade populacional baixa; isolamento social satisfatório em ambiente rural produtor \\
C - Riscos Naturais & 7.0 & Geologia muito estável e baixo risco de grandes desastres naturais \\
D - Autossuficiência & 9.0 & Recursos agrícolas e hídricos excepcionais; solo de elite mundial \\
E - Institucional & 7.5 & Forte segurança institucional municipal e suporte social baseado na produção \\
\midrule \textbf{GSS FINAL} & \textbf{7.4} & \textbf{Seguro (Altamente recomendado para quem busca autossuficiência em solo de elite e suporte institucional regional próximo).} \\ \bottomrule \end{tabular}
\end{table}
\subsubsection{Vulnerabilidades e
Mitigação}\label{vuln-Yasy_Cany-vulnerabilidades-e-mitigauxe7uxe3o}

\begin{itemize}
\tightlist
\item
  \textbf{Vulnerabilidade Tática Regional:} Insegurança ligada ao
  contexto do departamento (Mitigação: aquisição de propriedades com
  título perfeito e investimento em segurança patrimonial).
\item
  \textbf{Gargalo Logístico:} Dependência do eixo da Ruta PY03 para
  escoamento (Mitigação: mapeamento de rotas rurais vicinais em direção
  a Curuguaty).
\item
  \textbf{Isolamento de Saúde:} Dependência de centros vizinhos para
  cirurgias (Mitigação: manutenção de estoque médico básico e kit trauma
  local).
\end{itemize}
\subsubsection{Dossiê de Campo}\label{dossie-Yasy_Cany-dossiuxea-de-campo}
\begin{description}[style=nextline,leftmargin=0.5cm,font=\bfseries]
\item[Coordenadas]  24°20'S 55°25'W.
\item[Topografia]  Relevo predominantemente plano a suavemente ondulado, com solos de elite (Latossolo Vermelho). Altitude \textasciitilde 240m. Área de \textasciitilde 850 km². Geologicamente muito estável, risco sísmico desprezível.
\item[Alvos Estratégicos]  Polo de produção agrícola intensiva; Gigantesco complexo de silos e armazenamento de grãos; Ponto logístico sobre a Ruta PY03. Ausência de bases militares massivas, mas centro de gravidade econômica regional e ponto de passagem para o centro-oeste do departamento.
\item[Fallout]  Ventos predominantes do Norte (NE) e Leste. Baixíssimo risco de fallout direto; posição isolada e protegida.
\item[População]  20.399 habitantes (Censo 2022 Final). Município em plena expansão demográfica urbana e rural, ligada ao vigor do agronegócio.
\item[Densidade]  \textasciitilde 24,0 hab/\allowbreak km² (Densidade baixa, oferecendo isolamento tático superior e privacidade fora do núcleo comercial asfáltico).
\item[Vias de Acesso]  Localização estratégica sobre a Ruta PY03. Conectividade asfáltica de elite para Assunção e Curuguaty. Facilidade logística tática em tempos de paz, com rotas internas que favorecem o desaparecimento rural.
\item[Serviços]  Infraestrutura institucional média. Unidade de Saúde da Família (USF) local. Dependência de Curuguaty (approx. 25 km) para serviços hospitalares de alta complexidade e logística comercial especializada.
\item[Custo de Vida]  Baixo; economia baseada na exportação de grãos e pecuária mecanizada.
\item[Preço da Terra]  US\$ 8.000-12.000/\allowbreak ha (terras mecanizadas de elite - "Terra Roxa").
\item[Hidrologia]  Baixo risco de inundações sistêmicas devido à topografia favorável e drenagem natural eficiente. Presença de diversos arroios perenes de alta qualidade hídrica (ex: Rio Jejuí Guazú ao norte).
\item[Clima]  Subtropical Úmido. Verões quentes e chuvosos; invernos amenos com geadas ocasionais.
\item[Energia]  Rede nacional (Itaipu); boa estabilidade devido à proximidade com polos industriais regionais.
\item[Água]  Abundante via poços artesianos profundos e sistemas de captação local; excelente qualidade hídrica subterrânea.
\item[Qualidade do Solo]  Latossolo Vermelho de Elite; solo profundo, estruturado e de altíssima fertilidade natural. Excelente aptidão para soja, milho e pastagens mecanizadas.
\item[Resources Locais]  Grande polo produtor de grãos e proteína animal. Elevadíssimo potencial para autossuficiência alimentar básica em nível de propriedade devido à vasta área produtiva e baixa densidade populacional total.
\item[Segurança]  Zona rural pacífica em termos de crime comum urbano. No entanto, exige atenção tática regional devido ao contexto de Canindeyú (histórico de tensões agrárias e trânsito de ilícitos). Estabilidade social baseada na cultura produtora tradicional e pequenos/\allowbreak médios fazendeiros.
\item[Leis Local]  Município consolidado. \\ \textbf{Livre de Restrição} de Fronteira. \\ \textit{Aquisição de terra rural proibida para estrangeiros limítrofes.}
\end{description}
\paragraph{Solo (SoilGrids 2.0, média ponderada 0--30
cm)}

\begin{longtable}[]{@{}ll@{}}
\toprule\noalign{}
Parâmetro & Valor \\
\midrule\noalign{}
\endhead
\bottomrule\noalign{}
\endlastfoot
pH (H$_2$O) & 5.4 \\
Carbono orgânico (SOC) & 18.77 g/\allowbreak kg \\
Argila & 26.99 \% \\
Areia & 54.23 \% \\
Silte (calc.) & 18.8 \% \\
Densidade aparente & 1.28 g/\allowbreak cm³ \\
\textbf{Aptidão agrícola} & \textbf{Média} \\
\end{longtable}

Fonte: ISRIC SoilGrids 2.0 via WCS (acesso em 2026-03-21). Coords:
-24.3333°, -55.4167°. Média ponderada camadas 0-5, 5-15, 15-30 cm.

\begin{itemize}
\tightlist
\item
  \textbf{Homicidios/\allowbreak 100k:} \textasciitilde18,0 (Regional/\allowbreak Canindeyú).
\end{itemize}

\subsection{10. Analise de Riscos}

\begin{itemize}
\tightlist
\item
  \textbf{Cenario curto prazo:} Manutenção da tranquilidade rural e
  rentabilidade da produção mecanizada.
\item
  \textbf{Cenario medio prazo:} Valorização das terras pela expansão
  urbana planejada e consolidação industrial.
\end{itemize}

\subsection{11. Redacao Final}

\subsubsection{Diagnostico Integrado}

Yasy Cañy é um dos polos de resiliência e produtividade mais
equilibrados de Canindeyú. Oferece solos de elite mundial e uma
abundância massiva de recursos naturais, funcionando como um enclave de
segurança alimentar e hídrica. Sua força reside na excepcional qualidade
da base produtiva combinada com um suporte institucional sólido para o
contexto regional. É o local ideal para quem busca autonomia total em
ambiente de alta performance, transformando a fertilidade do solo em uma
barreira definitiva de proteção econômica e tática contra crises
externas.

\subsubsection{Por que sim}

\begin{itemize}
\tightlist
\item
  Solos de altíssima produtividade mundial e abundância de águas puras.
\item
  Ambiente social seguro e resiliente sustentado pela produção.
\item
  Conectividade asfáltica de elite via Ruta PY03 para Curuguaty.
\item
  Sem restrições da Lei de Fronteira.
\end{itemize}

\subsubsection{Por que nao}

\begin{itemize}
\tightlist
\item
  Insegurança tática regional latente no departamento de Canindeyú.
\item
  Alvo estratégico logístico em cenários de conflito por suprimentos.
\item
  Valor da terra em processo de valorização acelerada.
\end{itemize}

\subsubsection{Recomendacao Tecnica}

Altamente indicado para estabelecimentos de autossuficiência tecnificada
ou investimentos produtivos de larga escala. Requer autonomia energética
e médica complementar. O GSS de 7.4 reflete a excepcional força dos
recursos e da governança local em equilíbrio com a visibilidade
estratégica. Referencia atual de score: GSS 7.4 em 2026-03-05.

\subsection{13.}

\begin{itemize}
\tightlist
\item
  \begin{enumerate}
  \def\labelenumi{\arabic{enumi})}
  \tightlist
  \item
    Dados censitários validados (INE\footnote{Instituto Nacional de Estadística (Instituto Nacional de Estatística), órgão oficial de dados demográficos e censitários.} 2022).
  \end{enumerate}
\item
  \begin{enumerate}
  \def\labelenumi{\arabic{enumi})}
  \setcounter{enumi}{1}
  \tightlist
  \item
    Relatórios de produção agrícola mecanizada de alta precisão (MAG).
  \end{enumerate}
\item
  \begin{enumerate}
  \def\labelenumi{\arabic{enumi})}
  \setcounter{enumi}{2}
  \tightlist
  \item
    Mapa de solos e bacias hidrográficas de Canindeyú (Portal
    Geoestadístico).
  \end{enumerate}
\item
  \begin{enumerate}
  \def\labelenumi{\arabic{enumi})}
  \setcounter{enumi}{3}
  \tightlist
  \item
    Registro de segurança pública departamental.
  \end{enumerate}
\end{itemize}

\subsubsection{Combustível}

\textbf{Referência:} PETROPAR /\allowbreak  postos locais (2024) \textbf{Tipo de
localidade:} Interior

\begin{longtable}[]{@{}lll@{}}
\toprule\noalign{}
Combustível & USD/\allowbreak litro & Gs/\allowbreak litro (aprox.) \\
\midrule\noalign{}
\endhead
\bottomrule\noalign{}
\endlastfoot
Gasolina 93 oct & 1.00 & 7,400 \\
Gasolina 97 oct (premium) & 1.10 & 8,140 \\
Diesel & 0.93 & 6,882 \\
\end{longtable}

\begin{quote}
Preços podem variar ±5\% conforme posto e sazonalidade. Chaco e interior
remoto apresentam maior variação.
\end{quote}

\subsubsection{Cobertura Celular}

\textbf{Fonte:} CONATEL PY /\allowbreak  operadoras (2024)

\begin{longtable}[]{@{}ll@{}}
\toprule\noalign{}
Parâmetro & Valor \\
\midrule\noalign{}
\endhead
\bottomrule\noalign{}
\endlastfoot
Cobertura 4G população (dept.) & 78\% \\
Cobertura 4G área rural & 55\% \\
Melhor operadora & Tigo \\
Qualidade rural & moderada \\
\end{longtable}

\begin{quote}
Para áreas rurais fora do núcleo urbano, recomenda-se chip Tigo como
principal e Personal como backup.
\end{quote}

\subsubsection{Internet}

\textbf{Fonte:} CONATEL /\allowbreak  Speedtest Ookla (2024)

\begin{longtable}[]{@{}ll@{}}
\toprule\noalign{}
Parâmetro & Valor \\
\midrule\noalign{}
\endhead
\bottomrule\noalign{}
\endlastfoot
Velocidade média download & 35 Mbps \\
Domicílios com internet (dept.) & 45\% \\
Tecnologia predominante & rádio \\
Opção rural & Starlink disponível (\textasciitilde USD 44/\allowbreak mês) \\
\end{longtable}

\subsubsection{Mercado Imobiliário e Terra
Rural}

\textbf{Fonte:} INDERT /\allowbreak  Clasificados.com.py (2024)

\begin{longtable}[]{@{}ll@{}}
\toprule\noalign{}
Tipo & Referência \\
\midrule\noalign{}
\endhead
\bottomrule\noalign{}
\endlastfoot
Terra agrícola alta prod. (USD/\allowbreak ha) & 7,500 \\
Imóvel urbano (USD/\allowbreak m²) & 650 \\
Aluguel 2 quartos (USD/\allowbreak mês) & 260 \\
\end{longtable}

\begin{quote}
Valores de referência departamental. Localidades menores podem ter
preços 20--40\% abaixo da capital departamental. \#\#\# Saúde
\end{quote}

\textbf{Fonte:} MSPBS /\allowbreak  IPS Paraguay (2024-2026), consolidação
departamental e proxy local

\begin{longtable}[]{@{}ll@{}}
\toprule\noalign{}
Serviço & Disponibilidade \\
\midrule\noalign{}
\endhead
\bottomrule\noalign{}
\endlastfoot
USF /\allowbreak  Posto de Saúde & sim \\
Hospital Regional & sim \\
IPS (seguro social) & não \\
Farmácia & sim \\
Distância ao hospital de referência & \textasciitilde207 km (Salto del
Guaira) \\
\end{longtable}

\textbf{Principais estabelecimentos:} USF local; referencia hospitalar
em Salto del Guaira

\textbf{Observação para imigrantes:} Atendimento primario local ou em
raio curto; casos de maior complexidade seguem para Salto del Guaira.
Cobertura privada continua recomendada para especialidades e urgências
de maior complexidade.
\subsubsection{Indicadores Sociais}

\textbf{Fontes:} PNUD 2020, INE\footnote{Instituto Nacional de Estadística (Instituto Nacional de Estatística), órgão oficial de dados demográficos e censitários.} Censo 2022, INE\footnote{Instituto Nacional de Estadística (Instituto Nacional de Estatística), órgão oficial de dados demográficos e censitários.} EPHC 2023, Ministerio
Público 2024\\
\textbf{Nota:} Valores marcados como \emph{(dept.)} referem-se ao
departamento de Canindeyú; valores \emph{(dist.)} são específicos deste
distrito. Dados marcados \emph{(est.)} são estimativas.

\begin{longtable}[]{@{}
  >{\raggedright\arraybackslash}p{(\linewidth - 4\tabcolsep) * \real{0.4231}}
  >{\raggedright\arraybackslash}p{(\linewidth - 4\tabcolsep) * \real{0.2692}}
  >{\raggedright\arraybackslash}p{(\linewidth - 4\tabcolsep) * \real{0.3077}}@{}}
\toprule\noalign{}
\begin{minipage}[b]{\linewidth}\raggedright
Indicador
\end{minipage} & \begin{minipage}[b]{\linewidth}\raggedright
Valor
\end{minipage} & \begin{minipage}[b]{\linewidth}\raggedright
Âmbito
\end{minipage} \\
\midrule\noalign{}
\endhead
\bottomrule\noalign{}
\endlastfoot
IDH (2020) & 0,685 & dept. \\
Ranking IDH nacional & 15/\allowbreak 18 & dept. \\
Esperança de vida & 71,2 anos & dept. \\
Escolaridade média & 7,3 anos & dept. \\
RNB per capita (USD PPA) & 7.100 & dept. \\
Pobreza monetária (\%) & 35,1\% & dept. \\
Pobreza extrema (\%) & 10,8\% & dept. \\
Índice de Gini & 0,467 & dept. \\
Acesso a água potável (\%) & 67,4\% & dept. \\
Acesso a saneamento (\%) & 64,8\% & dept. \\
Taxa de homicídios (est., /\allowbreak 100k hab) & \textasciitilde15--25
(estimativa, significativamente acima da média nacional) & dept. \\
Índice de segurança & baixo & dept. \\
População (Censo 2022) & 20.399 habitantes & dist. \\
Idade mediana & N/\allowbreak D & dist. \\
Taxa de fecundidade & N/\allowbreak D & dist. \\
\end{longtable}
\subsubsection{Combustível}

\textbf{Referência:} PETROPAR /\allowbreak  postos locais (2024) \textbf{Tipo de
localidade:} Interior

\begin{longtable}[]{@{}lll@{}}
\toprule\noalign{}
Combustível & USD/\allowbreak litro & Gs/\allowbreak litro (aprox.) \\
\midrule\noalign{}
\endhead
\bottomrule\noalign{}
\endlastfoot
Gasolina 93 oct & 1.00 & 7,400 \\
Gasolina 97 oct (premium) & 1.10 & 8,140 \\
Diesel & 0.93 & 6,882 \\
\end{longtable}

\begin{quote}
Preços podem variar ±5\% conforme posto e sazonalidade. Chaco e interior
remoto apresentam maior variação.
\end{quote}

\subsubsection{Cobertura Celular}

\textbf{Fonte:} CONATEL PY /\allowbreak  operadoras (2024)

\begin{longtable}[]{@{}ll@{}}
\toprule\noalign{}
Parâmetro & Valor \\
\midrule\noalign{}
\endhead
\bottomrule\noalign{}
\endlastfoot
Cobertura 4G população (dept.) & 78\% \\
Cobertura 4G área rural & 55\% \\
Melhor operadora & Tigo \\
Qualidade rural & moderada \\
\end{longtable}

\begin{quote}
Para áreas rurais fora do núcleo urbano, recomenda-se chip Tigo como
principal e Personal como backup.
\end{quote}

\subsubsection{Internet}

\textbf{Fonte:} CONATEL /\allowbreak  Speedtest Ookla (2024)

\begin{longtable}[]{@{}ll@{}}
\toprule\noalign{}
Parâmetro & Valor \\
\midrule\noalign{}
\endhead
\bottomrule\noalign{}
\endlastfoot
Velocidade média download & 35 Mbps \\
Domicílios com internet (dept.) & 45\% \\
Tecnologia predominante & rádio \\
Opção rural & Starlink disponível (\textasciitilde USD 44/\allowbreak mês) \\
\end{longtable}

\subsubsection{Mercado Imobiliário e Terra
Rural}

\textbf{Fonte:} INDERT /\allowbreak  Clasificados.com.py (2024)

\begin{longtable}[]{@{}ll@{}}
\toprule\noalign{}
Tipo & Referência \\
\midrule\noalign{}
\endhead
\bottomrule\noalign{}
\endlastfoot
Terra agrícola alta prod. (USD/\allowbreak ha) & 7,500 \\
Imóvel urbano (USD/\allowbreak m²) & 650 \\
Aluguel 2 quartos (USD/\allowbreak mês) & 260 \\
\end{longtable}

\begin{quote}
Valores de referência departamental. Localidades menores podem ter
preços 20--40\% abaixo da capital departamental. \#\#\# Saúde
\end{quote}

\textbf{Fonte:} MSPBS /\allowbreak  IPS Paraguay (2024-2026), consolidação
departamental e proxy local

\begin{longtable}[]{@{}ll@{}}
\toprule\noalign{}
Serviço & Disponibilidade \\
\midrule\noalign{}
\endhead
\bottomrule\noalign{}
\endlastfoot
USF /\allowbreak  Posto de Saúde & sim \\
Hospital Regional & sim \\
IPS (seguro social) & não \\
Farmácia & sim \\
Distância ao hospital de referência & \textasciitilde207 km (Salto del
Guaira) \\
\end{longtable}

\textbf{Principais estabelecimentos:} USF local; referencia hospitalar
em Salto del Guaira

\textbf{Observação para imigrantes:} Atendimento primario local ou em
raio curto; casos de maior complexidade seguem para Salto del Guaira.
Cobertura privada continua recomendada para especialidades e urgências
de maior complexidade.
\newpage
\secaoDiagnostico{Yby Pyta}{\mapaConteudo{-24.2500}{-55.2500}}{Yby Pytá é um dos enclaves produtivos mais discretos de Canindeyú.
Oferece solos de elite mundial e um ambiente social tranquilo, ``fora do
radar'' das grandes crises demográficas nacionais. Sua força reside na
abundância de recursos produtivos e na baixa densidade populacional, o
que garante privacidade e isolamento tático superior. O principal
desafio é a dependência de centros como Curuguaty para serviços
especializados e a necessidade de vigilância tática regional. É o local
ideal para estabelecimentos de autossuficiência que busquem solo de alta
performance em ambiente de paz social.

\subsubsection{Por que sim}

\begin{itemize}
\tightlist
\item
  Solos de altíssima produtividade e abundância de águas puras.
\item
  Baixíssima densidade populacional (privacidade absoluta).
\item
  Localização discreta e protegida de grandes eixos de conflito
  primários.
\item
  Sem restrições da Lei de Fronteira.
\end{itemize}

\subsubsection{Por que não}

\begin{itemize}
\tightlist
\item
  Infraestrutura de saúde local limitada a atendimentos primários.
\item
  Insegurança tática regional latente no departamento.
\item
  Dependência de estradas internas de terra para acesso às propriedades.
\end{itemize}

\subsubsection{Recomendação
Técnica}

Altamente indicado para estabelecimentos de autossuficiência que busquem
isolamento tático em solo de alta performance. Requer autonomia em
logística básica e saúde. O GSS de 6.5 reflete a segurança do isolamento
em equilíbrio com a limitação institucional local. Referencia atual de
score: GSS 6.5 em 2026-03-05.

\subsubsection{}

\begin{itemize}
\tightlist
\item
  \begin{enumerate}
  \def\labelenumi{\arabic{enumi})}
  \tightlist
  \item
    Dados censitários validados (INE\footnote{Instituto Nacional de Estadística (Instituto Nacional de Estatística), órgão oficial de dados demográficos e censitários.} 2022).
  \end{enumerate}
\item
  \begin{enumerate}
  \def\labelenumi{\arabic{enumi})}
  \setcounter{enumi}{1}
  \tightlist
  \item
    Relatórios de produção agrícola mecanizada (MAG).
  \end{enumerate}
\item
  \begin{enumerate}
  \def\labelenumi{\arabic{enumi})}
  \setcounter{enumi}{2}
  \tightlist
  \item
    Mapa de solos e aptidão agrícola de Canindeyú (Portal
    Geoestadístico).
  \end{enumerate}
\item
  \begin{enumerate}
  \def\labelenumi{\arabic{enumi})}
  \setcounter{enumi}{3}
  \tightlist
  \item
    Registro de segurança pública departamental.
  \end{enumerate}
\end{itemize}}
\begin{table}[ht!]
\small \begin{tabular}{p{3.0cm}p{0.8cm}p{6.0cm}} \toprule
\textbf{Critério} & \textbf{Nota} & \textbf{Justificativa} \\ \midrule
A - Ameaças Estratégicas & 6.0 & Localização remota; boa invisibilidade tática, equilibrada por riscos regionais do departamento \\
B - Risco Social & 7.5 & Densidade populacional baixa; isolamento social superior e privacidade absoluta \\
C - Riscos Naturais & 7.0 & Geologia muito estável e baixo risco de desastres naturais graves \\
D - Autossuficiência & 7.5 & Recursos agrícolas de elite e abundância hídrica subterrânea \\
E - Institucional & 4.5 & Suporte institucional básico; serviços hospitalares dependentes de cidades vizinhas \\
\midrule \textbf{GSS FINAL} & \textbf{6.5} & \textbf{Moderadamente Seguro (Ideal para quem busca isolamento rural produtivo em solo de alta performance, ciente do contexto regional).} \\ \bottomrule \end{tabular}
\end{table}
\subsubsection{Vulnerabilidades e
Mitigação}\label{vuln-Yby_Pyta-vulnerabilidades-e-mitigauxe7uxe3o}

\begin{itemize}
\tightlist
\item
  \textbf{Vulnerabilidade Tática Regional:} Insegurança ligada ao
  contexto de segurança do departamento (Mitigação: aquisição de
  propriedades com documentação impecável e investimento em segurança
  patrimonial).
\item
  \textbf{Isolamento de Saúde:} Distância significativa de hospitais
  cirúrgicos (Mitigação: manutenção de estoque médico básico e kit
  trauma local).
\item
  \textbf{Dependência de Serviços Externos:} Pequena escala comercial
  exige logística para centros maiores (Mitigação: estoque de
  suprimentos para 60 dias).
\end{itemize}
\subsubsection{Dossiê de Campo}\label{dossie-Yby_Pyta-dossiuxea-de-campo}
\begin{description}[style=nextline,leftmargin=0.5cm,font=\bfseries]
\item[Coordenadas]  24°15'S 55°15'W.
\item[Topographical]  Relevo predominantemente plano a suavemente ondulado, situado em zona de transição entre o centro do departamento e a Reserva Mbaracayú ao sul. Altitude \textasciitilde 250m. Área de \textasciitilde 807 km². Geologicamente muito estável, risco sísmico desprezível.
\item[Alvos Estratégicos]  Áreas de expansão da fronteira agrícola (soja e milho); Silos de armazenamento de grãos; Proximidade tática com a Reserva Natural del Bosque Mbaracayú (Recurso ambiental e hídrico estratégico). Ausência de alvos militares diretos, oferecendo excelente invisibilidade estratégica e proteção por anonimato.
\item[Fallout]  Ventos predominantes do Norte (NE) e Leste. Baixíssimo risco de fallout de alvos continentais primários.
\item[População]  9.154 habitantes (Censo 2022 Final). Município jovem, em fase de crescimento demográfico vinculado à produção agrícola.
\item[Densidade]  \textasciitilde 11,3 hab/\allowbreak km² (Densidade baixa, oferecendo isolamento tático e privacidade fora do núcleo urbano em formação).
\item[Vias de Acesso]  Acesso via ramais de terra e alguns trechos pavimentados que conectam com a Ruta PY03 e PY13. Localização que favorece a discrição, mantendo-se afastada dos grandes fluxos de trânsito internacional.
\item[Serviços]  Infraestrutura institucional básica. Unidade de Saúde da Família (USF) local. Dependência de Curuguaty (\textasciitilde 45 km) ou Katueté para atendimentos hospitalares especializados e logística comercial avançada.
\item[Custo de Vida]  Muito baixo; economia baseada na exportação de grãos e pecuária.
\item[Preço da Terra]  US\$ 7.000-10.000/\allowbreak ha (terras mecanizadas); US\$ 2.500-4.500/\allowbreak ha (áreas de uso misto/\allowbreak reserva).
\item[Hidrologia]  Baixo risco de inundações sistêmicas devido à topografia ondulada. Diversos arroios cruzam o distrito, garantindo irrigação natural. Risco moderado de isolamento de colônias rurais por danos em estradas de terra em períodos chuvosos (média 1.700 mm/\allowbreak ano).
\item[Clima]  Subtropical Úmido. Verões intensos e chuvosos; invernos amenos.
\item[Energia]  Rede nacional (Itaipu); instável em áreas rurais devido ao isolamento das linhas e falta de redundância.
\item[Água]  Abundante via poços artesianos profundos e proximidade com as bacias hidrográficas da reserva natural.
\item[Qualidade do Solo]  Latossolo Vermelho de Elite; solo profundo e de altíssima fertilidade natural. Excelente aptidão para soja, milho e mandioca.
\item[Resources Locais]  Grande produção de soja, milho e gado de corte. Elevadíssimo potencial para autossuficiência alimentar básica em nível de propriedade devido à vasta área produtiva e baixa população.
\item[Segurança]  Zona rural pacífica em termos de crime comum. Exige cautela tática regional por estar em zona de expansão agrícola com histórico de reivindicações de terra e proximidade com rotas sensíveis no departamento. Estabilidade social baseada na cultura agropecuária tradicional e pequenos produtores.
\item[Leis Local: Município consolidado. Livre de Restrição de Fronteira]  Localizado fora da faixa de 50km da fronteira seca internacional.
\end{description}
\paragraph{Solo (SoilGrids 2.0, média ponderada 0--30
cm)}

\begin{longtable}[]{@{}ll@{}}
\toprule\noalign{}
Parâmetro & Valor \\
\midrule\noalign{}
\endhead
\bottomrule\noalign{}
\endlastfoot
pH (H$_2$O) & 5.32 \\
Carbono orgânico (SOC) & 20.83 g/\allowbreak kg \\
Argila & 25.63 \% \\
Areia & 56.5 \% \\
Silte (calc.) & 17.9 \% \\
Densidade aparente & 1.23 g/\allowbreak cm³ \\
\textbf{Aptidão agrícola} & \textbf{Média} \\
\end{longtable}

Fonte: ISRIC SoilGrids 2.0 via WCS (acesso em 2026-03-21). Coords:
-24.25°, -55.25°. Média ponderada camadas 0-5, 5-15, 15-30 cm.

\begin{itemize}
\tightlist
\item
  \textbf{Homicidios/\allowbreak 100k:} \textasciitilde18,0 (Regional/\allowbreak Canindeyú).
\end{itemize}

\subsection{10. Analise de Riscos}

\begin{itemize}
\tightlist
\item
  \textbf{Cenario curto prazo:} Manutenção da tranquilidade rural e
  rentabilidade da agricultura familiar e mecanizada.
\item
  \textbf{Cenario medio prazo:} Impacto positivo do desenvolvimento da
  infraestrutura asfáltica vicinal.
\end{itemize}

\subsection{11. Redacao Final}

\subsubsection{Diagnostico Integrado}

Yby Pytá é um dos enclaves produtivos mais discretos de Canindeyú.
Oferece solos de elite mundial e um ambiente social tranquilo, ``fora do
radar'' das grandes crises demográficas nacionais. Sua força reside na
abundância de recursos produtivos e na baixa densidade populacional, o
que garante privacidade e isolamento tático superior. O principal
desafio é a dependência de centros como Curuguaty para serviços
especializados e a necessidade de vigilância tática regional. É o local
ideal para estabelecimentos de autossuficiência que busquem solo de alta
performance em ambiente de paz social.

\subsubsection{Por que sim}

\begin{itemize}
\tightlist
\item
  Solos de altíssima produtividade e abundância de águas puras.
\item
  Baixíssima densidade populacional (privacidade absoluta).
\item
  Localização discreta e protegida de grandes eixos de conflito
  primários.
\item
  Sem restrições da Lei de Fronteira.
\end{itemize}

\subsubsection{Por que nao}

\begin{itemize}
\tightlist
\item
  Infraestrutura de saúde local limitada a atendimentos primários.
\item
  Insegurança tática regional latente no departamento.
\item
  Dependência de estradas internas de terra para acesso às propriedades.
\end{itemize}

\subsubsection{Recomendacao Tecnica}

Altamente indicado para estabelecimentos de autossuficiência que busquem
isolamento tático em solo de alta performance. Requer autonomia em
logística básica e saúde. O GSS de 6.5 reflete a segurança do isolamento
em equilíbrio com a limitação institucional local. Referencia atual de
score: GSS 6.5 em 2026-03-05.

\subsection{13.}

\begin{itemize}
\tightlist
\item
  \begin{enumerate}
  \def\labelenumi{\arabic{enumi})}
  \tightlist
  \item
    Dados censitários validados (INE\footnote{Instituto Nacional de Estadística (Instituto Nacional de Estatística), órgão oficial de dados demográficos e censitários.} 2022).
  \end{enumerate}
\item
  \begin{enumerate}
  \def\labelenumi{\arabic{enumi})}
  \setcounter{enumi}{1}
  \tightlist
  \item
    Relatórios de produção agrícola mecanizada (MAG).
  \end{enumerate}
\item
  \begin{enumerate}
  \def\labelenumi{\arabic{enumi})}
  \setcounter{enumi}{2}
  \tightlist
  \item
    Mapa de solos e aptidão agrícola de Canindeyú (Portal
    Geoestadístico).
  \end{enumerate}
\item
  \begin{enumerate}
  \def\labelenumi{\arabic{enumi})}
  \setcounter{enumi}{3}
  \tightlist
  \item
    Registro de segurança pública departamental.
  \end{enumerate}
\end{itemize}

\subsubsection{Combustível}

\textbf{Referência:} PETROPAR /\allowbreak  postos locais (2024) \textbf{Tipo de
localidade:} Interior

\begin{longtable}[]{@{}lll@{}}
\toprule\noalign{}
Combustível & USD/\allowbreak litro & Gs/\allowbreak litro (aprox.) \\
\midrule\noalign{}
\endhead
\bottomrule\noalign{}
\endlastfoot
Gasolina 93 oct & 1.00 & 7,400 \\
Gasolina 97 oct (premium) & 1.10 & 8,140 \\
Diesel & 0.93 & 6,882 \\
\end{longtable}

\begin{quote}
Preços podem variar ±5\% conforme posto e sazonalidade. Chaco e interior
remoto apresentam maior variação.
\end{quote}

\subsubsection{Cobertura Celular}

\textbf{Fonte:} CONATEL PY /\allowbreak  operadoras (2024)

\begin{longtable}[]{@{}ll@{}}
\toprule\noalign{}
Parâmetro & Valor \\
\midrule\noalign{}
\endhead
\bottomrule\noalign{}
\endlastfoot
Cobertura 4G população (dept.) & 78\% \\
Cobertura 4G área rural & 55\% \\
Melhor operadora & Tigo \\
Qualidade rural & moderada \\
\end{longtable}

\begin{quote}
Para áreas rurais fora do núcleo urbano, recomenda-se chip Tigo como
principal e Personal como backup.
\end{quote}

\subsubsection{Internet}

\textbf{Fonte:} CONATEL /\allowbreak  Speedtest Ookla (2024)

\begin{longtable}[]{@{}ll@{}}
\toprule\noalign{}
Parâmetro & Valor \\
\midrule\noalign{}
\endhead
\bottomrule\noalign{}
\endlastfoot
Velocidade média download & 35 Mbps \\
Domicílios com internet (dept.) & 45\% \\
Tecnologia predominante & rádio \\
Opção rural & Starlink disponível (\textasciitilde USD 44/\allowbreak mês) \\
\end{longtable}

\subsubsection{Mercado Imobiliário e Terra
Rural}

\textbf{Fonte:} INDERT /\allowbreak  Clasificados.com.py (2024)

\begin{longtable}[]{@{}ll@{}}
\toprule\noalign{}
Tipo & Referência \\
\midrule\noalign{}
\endhead
\bottomrule\noalign{}
\endlastfoot
Terra agrícola alta prod. (USD/\allowbreak ha) & 7,500 \\
Imóvel urbano (USD/\allowbreak m²) & 650 \\
Aluguel 2 quartos (USD/\allowbreak mês) & 260 \\
\end{longtable}

\begin{quote}
Valores de referência departamental. Localidades menores podem ter
preços 20--40\% abaixo da capital departamental. \#\#\# Saúde
\end{quote}

\textbf{Fonte:} MSPBS /\allowbreak  IPS Paraguay (2024-2026), consolidação
departamental e proxy local

\begin{longtable}[]{@{}ll@{}}
\toprule\noalign{}
Serviço & Disponibilidade \\
\midrule\noalign{}
\endhead
\bottomrule\noalign{}
\endlastfoot
USF /\allowbreak  Posto de Saúde & sim \\
Hospital Regional & sim \\
IPS (seguro social) & não \\
Farmácia & sim \\
Distância ao hospital de referência & \textasciitilde139 km (Salto del
Guaira) \\
\end{longtable}

\textbf{Principais estabelecimentos:} USF local; referencia hospitalar
em Salto del Guaira

\textbf{Observação para imigrantes:} Atendimento primario local ou em
raio curto; casos de maior complexidade seguem para Salto del Guaira.
Cobertura privada continua recomendada para especialidades e urgências
de maior complexidade.
\subsubsection{Indicadores Sociais}

\textbf{Fontes:} PNUD 2020, INE\footnote{Instituto Nacional de Estadística (Instituto Nacional de Estatística), órgão oficial de dados demográficos e censitários.} Censo 2022, INE\footnote{Instituto Nacional de Estadística (Instituto Nacional de Estatística), órgão oficial de dados demográficos e censitários.} EPHC 2023, Ministerio
Público 2024\\
\textbf{Nota:} Valores marcados como \emph{(dept.)} referem-se ao
departamento de Canindeyú; valores \emph{(dist.)} são específicos deste
distrito. Dados marcados \emph{(est.)} são estimativas.

\begin{longtable}[]{@{}
  >{\raggedright\arraybackslash}p{(\linewidth - 4\tabcolsep) * \real{0.4231}}
  >{\raggedright\arraybackslash}p{(\linewidth - 4\tabcolsep) * \real{0.2692}}
  >{\raggedright\arraybackslash}p{(\linewidth - 4\tabcolsep) * \real{0.3077}}@{}}
\toprule\noalign{}
\begin{minipage}[b]{\linewidth}\raggedright
Indicador
\end{minipage} & \begin{minipage}[b]{\linewidth}\raggedright
Valor
\end{minipage} & \begin{minipage}[b]{\linewidth}\raggedright
Âmbito
\end{minipage} \\
\midrule\noalign{}
\endhead
\bottomrule\noalign{}
\endlastfoot
IDH (2020) & 0,685 & dept. \\
Ranking IDH nacional & 15/\allowbreak 18 & dept. \\
Esperança de vida & 71,2 anos & dept. \\
Escolaridade média & 7,3 anos & dept. \\
RNB per capita (USD PPA) & 7.100 & dept. \\
Pobreza monetária (\%) & 35,1\% & dept. \\
Pobreza extrema (\%) & 10,8\% & dept. \\
Índice de Gini & 0,467 & dept. \\
Acesso a água potável (\%) & 67,4\% & dept. \\
Acesso a saneamento (\%) & 64,8\% & dept. \\
Taxa de homicídios (est., /\allowbreak 100k hab) & \textasciitilde15--25
(estimativa, significativamente acima da média nacional) & dept. \\
Índice de segurança & baixo & dept. \\
População (Censo 2022) & 9.154 habitantes & dist. \\
Idade mediana & N/\allowbreak D & dist. \\
Taxa de fecundidade & N/\allowbreak D & dist. \\
\end{longtable}
\subsubsection{Combustível}

\textbf{Referência:} PETROPAR /\allowbreak  postos locais (2024) \textbf{Tipo de
localidade:} Interior

\begin{longtable}[]{@{}lll@{}}
\toprule\noalign{}
Combustível & USD/\allowbreak litro & Gs/\allowbreak litro (aprox.) \\
\midrule\noalign{}
\endhead
\bottomrule\noalign{}
\endlastfoot
Gasolina 93 oct & 1.00 & 7,400 \\
Gasolina 97 oct (premium) & 1.10 & 8,140 \\
Diesel & 0.93 & 6,882 \\
\end{longtable}

\begin{quote}
Preços podem variar ±5\% conforme posto e sazonalidade. Chaco e interior
remoto apresentam maior variação.
\end{quote}

\subsubsection{Cobertura Celular}

\textbf{Fonte:} CONATEL PY /\allowbreak  operadoras (2024)

\begin{longtable}[]{@{}ll@{}}
\toprule\noalign{}
Parâmetro & Valor \\
\midrule\noalign{}
\endhead
\bottomrule\noalign{}
\endlastfoot
Cobertura 4G população (dept.) & 78\% \\
Cobertura 4G área rural & 55\% \\
Melhor operadora & Tigo \\
Qualidade rural & moderada \\
\end{longtable}

\begin{quote}
Para áreas rurais fora do núcleo urbano, recomenda-se chip Tigo como
principal e Personal como backup.
\end{quote}

\subsubsection{Internet}

\textbf{Fonte:} CONATEL /\allowbreak  Speedtest Ookla (2024)

\begin{longtable}[]{@{}ll@{}}
\toprule\noalign{}
Parâmetro & Valor \\
\midrule\noalign{}
\endhead
\bottomrule\noalign{}
\endlastfoot
Velocidade média download & 35 Mbps \\
Domicílios com internet (dept.) & 45\% \\
Tecnologia predominante & rádio \\
Opção rural & Starlink disponível (\textasciitilde USD 44/\allowbreak mês) \\
\end{longtable}

\subsubsection{Mercado Imobiliário e Terra
Rural}

\textbf{Fonte:} INDERT /\allowbreak  Clasificados.com.py (2024)

\begin{longtable}[]{@{}ll@{}}
\toprule\noalign{}
Tipo & Referência \\
\midrule\noalign{}
\endhead
\bottomrule\noalign{}
\endlastfoot
Terra agrícola alta prod. (USD/\allowbreak ha) & 7,500 \\
Imóvel urbano (USD/\allowbreak m²) & 650 \\
Aluguel 2 quartos (USD/\allowbreak mês) & 260 \\
\end{longtable}

\begin{quote}
Valores de referência departamental. Localidades menores podem ter
preços 20--40\% abaixo da capital departamental. \#\#\# Saúde
\end{quote}

\textbf{Fonte:} MSPBS /\allowbreak  IPS Paraguay (2024-2026), consolidação
departamental e proxy local

\begin{longtable}[]{@{}ll@{}}
\toprule\noalign{}
Serviço & Disponibilidade \\
\midrule\noalign{}
\endhead
\bottomrule\noalign{}
\endlastfoot
USF /\allowbreak  Posto de Saúde & sim \\
Hospital Regional & sim \\
IPS (seguro social) & não \\
Farmácia & sim \\
Distância ao hospital de referência & \textasciitilde139 km (Salto del
Guaira) \\
\end{longtable}

\textbf{Principais estabelecimentos:} USF local; referencia hospitalar
em Salto del Guaira

\textbf{Observação para imigrantes:} Atendimento primario local ou em
raio curto; casos de maior complexidade seguem para Salto del Guaira.
Cobertura privada continua recomendada para especialidades e urgências
de maior complexidade.
\newpage
\secaoDiagnostico{Ybyrarovana}{\mapaConteudo{-24.5000}{-55.2500}}{Ybyrarobaná é um enclave produtivo de baixa visibilidade no coração de
Canindeyú. Oferece solos de elite mundial e um ambiente social
tranquilo, ``fora do radar'' das grandes crises demográficas
metropolitanas. Sua força reside na abundância de recursos naturais e na
baixa densidade populacional, garantindo privacidade tática superior. O
principal desafio é a fragilidade institucional local e a dependência de
centros maiores para serviços vitais, o que exige do ocupante autonomia
logística e em saúde. É o local ideal para estabelecimentos de
autossuficiência que busquem solo de alta performance em ambiente de paz
social.

\subsubsection{Por que sim}

\begin{itemize}
\tightlist
\item
  Solos de altíssima produtividade e abundância de águas puras
  subterrâneas.
\item
  Baixíssima densidade populacional (privacidade absoluta).
\item
  Localização discreta e protegida de eixos de conflito primários.
\item
  Sem restrições da Lei de Fronteira.
\end{itemize}

\subsubsection{Por que não}

\begin{itemize}
\tightlist
\item
  Infraestrutura de saúde local limitada a atendimentos primários.
\item
  Insegurança tática regional latente no departamento de Canindeyú.
\item
  Dependência de estradas internas de terra para acesso às propriedades.
\end{itemize}

\subsubsection{Recomendação
Técnica}

Altamente indicado para estabelecimentos de autossuficiência que busquem
isolamento tático em solo de alta performance. Requer autonomia em
logística básica e saúde. O GSS de 5.7 reflete a segurança do isolamento
em equilíbrio com a vulnerabilidade institucional da localidade.
Referencia atual de score: GSS 5.7 em 2026-03-05.

\subsubsection{}

\begin{itemize}
\tightlist
\item
  \begin{enumerate}
  \def\labelenumi{\arabic{enumi})}
  \tightlist
  \item
    Dados censitários validados (INE\footnote{Instituto Nacional de Estadística (Instituto Nacional de Estatística), órgão oficial de dados demográficos e censitários.} 2022).
  \end{enumerate}
\item
  \begin{enumerate}
  \def\labelenumi{\arabic{enumi})}
  \setcounter{enumi}{1}
  \tightlist
  \item
    Relatórios de produção agrícola mecanizada (MAG).
  \end{enumerate}
\item
  \begin{enumerate}
  \def\labelenumi{\arabic{enumi})}
  \setcounter{enumi}{2}
  \tightlist
  \item
    Mapa de solos e aptidão agrícola de Canindeyú (Portal
    Geoestadístico).
  \end{enumerate}
\item
  \begin{enumerate}
  \def\labelenumi{\arabic{enumi})}
  \setcounter{enumi}{3}
  \tightlist
  \item
    Registro de segurança pública departamental.
  \end{enumerate}
\end{itemize}}
\begin{table}[ht!]
\small \begin{tabular}{p{3.0cm}p{0.8cm}p{6.0cm}} \toprule
\textbf{Critério} & \textbf{Nota} & \textbf{Justificativa} \\ \midrule
A - Ameaças Estratégicas & 5.5 & Localização remota; boa invisibilidade tática, equilibrada por riscos regionais do departamento \\
B - Risco Social & 7.0 & Densidade populacional baixa; isolamento social e privacidade satisfatórios \\
C - Riscos Naturais & 7.0 & Geologia muito estável e baixo risco de desastres naturais graves \\
D - Autossuficiência & 7.0 & Bons recursos agrícolas e hídricos; limitado pela escala comercial local limitada \\
E - Institucional & 2.0 & Nota penalizada no ranking nacional para refletir a dependência extrema de insumos globais e infraestrutura local mínima - GSS 5.7 total \\
\midrule \textbf{GSS FINAL} & \textbf{5.7} & \textbf{Moderadamente Seguro (Ideal para quem busca isolamento rural em solo de alta produtividade, aceitando a precariedade institucional local).} \\ \bottomrule \end{tabular}
\end{table}
\subsubsection{Vulnerabilidades e
Mitigação}\label{vuln-Ybyrarovana-vulnerabilidades-e-mitigauxe7uxe3o}

\begin{itemize}
\tightlist
\item
  \textbf{Isolamento de Saúde:} Distância significativa de hospitais
  cirúrgicos (Mitigação: manutenção de estoque médico básico e kit
  trauma local).
\item
  \textbf{Dependência de Serviços Externos:} Pequena escala comercial
  local (Mitigação: estoque de suprimentos técnicos e básicos para 60
  dias).
\item
  \textbf{Fragilidade Elétrica:} Rede rural vulnerável a intempéries
  (Mitigação: instalação de sistemas solares fotovoltaicos e geradores).
\end{itemize}
\subsubsection{Dossiê de Campo}\label{dossie-Ybyrarovana-dossiuxea-de-campo}
\begin{description}[style=nextline,leftmargin=0.5cm,font=\bfseries]
\item[Coordenadas]  24°30'S 55°15'W.
\item[Topographical]  Relevo predominantemente plano a suavemente ondulado, situado em zona de transição agrícola. Altitude \textasciitilde 240m. Área de \textasciitilde 1.086 km². Geologicamente muito estável, risco sísmico desprezível.
\item[Alvos Estratégicos]  Áreas de produção agrícola intensiva (soja e milho); Silos de armazenamento de grãos; Proximidade tática com a Reserva Natural del Bosque Mbaracayú. Ausência de alvos militares de grande porte, oferecendo boa invisibilidade estratégica, embora situada em departamento sob vigilância regional permanente.
\item[Fallout]  Ventos predominantes do Norte (NE) e Leste. Baixíssimo risco de fallout de alvos continentais primários.
\item[População]  12.145 habitantes (Censo 2022 Final). Município jovem com perfil demográfico rural dinâmico.
\item[Densidade]  \textasciitilde 11,2 hab/\allowbreak km² (Densidade baixa, oferecendo isolamento tático e privacidade fora dos pequenos núcleos povoados).
\item[Vias de Acesso]  Acesso via ramais de terra e secundários que conectam com a Ruta PY03. Localização que favorece a discrição, mantendo-se afastada dos grandes eixos de circulação metropolitana. Malha interna de terra exige veículos robustos em períodos chuvosos.
\item[Serviços]  Infraestrutura institucional mínima. Unidade de Saúde da Família (USF) local. Dependência crítica de Curuguaty (\textasciitilde 60 km) para atendimentos hospitalares especializados e logística comercial avançada.
\item[Custo de Vida]  Muito baixo; economia baseada na exportação de grãos e pecuária.
\item[Preço da Terra]  US\$ 7.000-10.000/\allowbreak ha (terras mecanizadas); US\$ 2.500-4.500/\allowbreak ha (áreas de uso misto).
\item[Hidrologia]  Baixo risco de inundações sistêmicas devido à topografia ondulada favorável à drenagem. Presença de diversos arroios perenes. Risco moderado de isolamento de comunidades rurais por danos em estradas de terra em períodos de chuva intensa (média 1.700 mm/\allowbreak ano).
\item[Clima]  Subtropical Úmido. Verões intensos e chuvosos; invernos amenos com geadas ocasionais benéficas.
\item[Energia]  Rede nacional (Itaipu); instável em áreas rurais devido ao isolamento da rede e exposição a intempéries.
\item[Água]  Abundante via poços artesianos e proximidade com as bacias hidrográficas da reserva natural.
\item[Qualidade do Solo]  Latossolo Vermelho de Alta Performance; solo profundo, estruturado e de altíssima fertilidade natural. Excelente aptidão para soja e milho.
\item[Resources Locais]  Grande produção de grãos e gado de corte. Elevadíssimo potencial para autossuficiência alimentar básica em nível de propriedade devido à vasta área produtiva e baixa densidade populacional.
\item[Segurança]  Zona rural pacífica em termos de crime comum. No entanto, exige atenção tática regional por estar em departamento de fronteira sensível ao trânsito de ilícitos e histórico de conflitos agrários. Estabilidade social baseada na cultura agropecuária tradicional.
\item[Leis Local: Município consolidado, criado em 2011. Livre de Restrição de Fronteira]  Localizado fora da faixa de 50km da fronteira seca internacional.
\end{description}
\paragraph{Solo (SoilGrids 2.0, média ponderada 0--30
cm)}

\begin{longtable}[]{@{}ll@{}}
\toprule\noalign{}
Parâmetro & Valor \\
\midrule\noalign{}
\endhead
\bottomrule\noalign{}
\endlastfoot
pH (H$_2$O) & 5.42 \\
Carbono orgânico (SOC) & 19.88 g/\allowbreak kg \\
Argila & 27.33 \% \\
Areia & 52.67 \% \\
Silte (calc.) & 20.0 \% \\
Densidade aparente & 1.27 g/\allowbreak cm³ \\
\textbf{Aptidão agrícola} & \textbf{Média} \\
\end{longtable}

Fonte: ISRIC SoilGrids 2.0 via WCS (acesso em 2026-03-21). Coords:
-24.5°, -55.25°. Média ponderada camadas 0-5, 5-15, 15-30 cm.

\begin{itemize}
\tightlist
\item
  \textbf{Homicidios/\allowbreak 100k:} \textasciitilde18,0 (Regional/\allowbreak Canindeyú).
\end{itemize}

\subsection{10. Analise de Riscos}

\begin{itemize}
\tightlist
\item
  \textbf{Cenario curto prazo:} Manutenção da tranquilidade rural e
  rentabilidade da agricultura familiar e mecanizada.
\item
  \textbf{Cenario medio prazo:} Impacto positivo da modernização das
  redes elétricas rurais no departamento.
\end{itemize}

\subsection{11. Redacao Final}

\subsubsection{Diagnostico Integrado}

Ybyrarobaná é um enclave produtivo de baixa visibilidade no coração de
Canindeyú. Oferece solos de elite mundial e um ambiente social
tranquilo, ``fora do radar'' das grandes crises demográficas
metropolitanas. Sua força reside na abundância de recursos naturais e na
baixa densidade populacional, garantindo privacidade tática superior. O
principal desafio é a fragilidade institucional local e a dependência de
centros maiores para serviços vitais, o que exige do ocupante autonomia
logística e em saúde. É o local ideal para estabelecimentos de
autossuficiência que busquem solo de alta performance em ambiente de paz
social.

\subsubsection{Por que sim}

\begin{itemize}
\tightlist
\item
  Solos de altíssima produtividade e abundância de águas puras
  subterrâneas.
\item
  Baixíssima densidade populacional (privacidade absoluta).
\item
  Localização discreta e protegida de eixos de conflito primários.
\item
  Sem restrições da Lei de Fronteira.
\end{itemize}

\subsubsection{Por que nao}

\begin{itemize}
\tightlist
\item
  Infraestrutura de saúde local limitada a atendimentos primários.
\item
  Insegurança tática regional latente no departamento de Canindeyú.
\item
  Dependência de estradas internas de terra para acesso às propriedades.
\end{itemize}

\subsubsection{Recomendacao Tecnica}

Altamente indicado para estabelecimentos de autossuficiência que busquem
isolamento tático em solo de alta performance. Requer autonomia em
logística básica e saúde. O GSS de 5.7 reflete a segurança do isolamento
em equilíbrio com a vulnerabilidade institucional da localidade.
Referencia atual de score: GSS 5.7 em 2026-03-05.

\subsection{13.}

\begin{itemize}
\tightlist
\item
  \begin{enumerate}
  \def\labelenumi{\arabic{enumi})}
  \tightlist
  \item
    Dados censitários validados (INE\footnote{Instituto Nacional de Estadística (Instituto Nacional de Estatística), órgão oficial de dados demográficos e censitários.} 2022).
  \end{enumerate}
\item
  \begin{enumerate}
  \def\labelenumi{\arabic{enumi})}
  \setcounter{enumi}{1}
  \tightlist
  \item
    Relatórios de produção agrícola mecanizada (MAG).
  \end{enumerate}
\item
  \begin{enumerate}
  \def\labelenumi{\arabic{enumi})}
  \setcounter{enumi}{2}
  \tightlist
  \item
    Mapa de solos e aptidão agrícola de Canindeyú (Portal
    Geoestadístico).
  \end{enumerate}
\item
  \begin{enumerate}
  \def\labelenumi{\arabic{enumi})}
  \setcounter{enumi}{3}
  \tightlist
  \item
    Registro de segurança pública departamental.
  \end{enumerate}
\end{itemize}

\subsubsection{Combustível}

\textbf{Referência:} PETROPAR /\allowbreak  postos locais (2024) \textbf{Tipo de
localidade:} Interior

\begin{longtable}[]{@{}lll@{}}
\toprule\noalign{}
Combustível & USD/\allowbreak litro & Gs/\allowbreak litro (aprox.) \\
\midrule\noalign{}
\endhead
\bottomrule\noalign{}
\endlastfoot
Gasolina 93 oct & 1.00 & 7,400 \\
Gasolina 97 oct (premium) & 1.10 & 8,140 \\
Diesel & 0.93 & 6,882 \\
\end{longtable}

\begin{quote}
Preços podem variar ±5\% conforme posto e sazonalidade. Chaco e interior
remoto apresentam maior variação.
\end{quote}

\subsubsection{Cobertura Celular}

\textbf{Fonte:} CONATEL PY /\allowbreak  operadoras (2024)

\begin{longtable}[]{@{}ll@{}}
\toprule\noalign{}
Parâmetro & Valor \\
\midrule\noalign{}
\endhead
\bottomrule\noalign{}
\endlastfoot
Cobertura 4G população (dept.) & 78\% \\
Cobertura 4G área rural & 55\% \\
Melhor operadora & Tigo \\
Qualidade rural & moderada \\
\end{longtable}

\begin{quote}
Para áreas rurais fora do núcleo urbano, recomenda-se chip Tigo como
principal e Personal como backup.
\end{quote}

\subsubsection{Internet}

\textbf{Fonte:} CONATEL /\allowbreak  Speedtest Ookla (2024)

\begin{longtable}[]{@{}ll@{}}
\toprule\noalign{}
Parâmetro & Valor \\
\midrule\noalign{}
\endhead
\bottomrule\noalign{}
\endlastfoot
Velocidade média download & 35 Mbps \\
Domicílios com internet (dept.) & 45\% \\
Tecnologia predominante & rádio \\
Opção rural & Starlink disponível (\textasciitilde USD 44/\allowbreak mês) \\
\end{longtable}

\subsubsection{Mercado Imobiliário e Terra
Rural}

\textbf{Fonte:} INDERT /\allowbreak  Clasificados.com.py (2024)

\begin{longtable}[]{@{}ll@{}}
\toprule\noalign{}
Tipo & Referência \\
\midrule\noalign{}
\endhead
\bottomrule\noalign{}
\endlastfoot
Terra agrícola alta prod. (USD/\allowbreak ha) & 7,500 \\
Imóvel urbano (USD/\allowbreak m²) & 650 \\
Aluguel 2 quartos (USD/\allowbreak mês) & 260 \\
\end{longtable}

\begin{quote}
Valores de referência departamental. Localidades menores podem ter
preços 20--40\% abaixo da capital departamental. \#\#\# Saúde
\end{quote}

\textbf{Fonte:} MSPBS /\allowbreak  IPS Paraguay (2024-2026), consolidação
departamental e proxy local

\begin{longtable}[]{@{}ll@{}}
\toprule\noalign{}
Serviço & Disponibilidade \\
\midrule\noalign{}
\endhead
\bottomrule\noalign{}
\endlastfoot
USF /\allowbreak  Posto de Saúde & sim \\
Hospital Regional & sim \\
IPS (seguro social) & não \\
Farmácia & sim \\
Distância ao hospital de referência & \textasciitilde104 km (Salto del
Guaira) \\
\end{longtable}

\textbf{Principais estabelecimentos:} USF local; referencia hospitalar
em Salto del Guaira

\textbf{Observação para imigrantes:} Atendimento primario local ou em
raio curto; casos de maior complexidade seguem para Salto del Guaira.
Cobertura privada continua recomendada para especialidades e urgências
de maior complexidade.
\subsubsection{Indicadores Sociais}

\textbf{Fontes:} PNUD 2020, INE\footnote{Instituto Nacional de Estadística (Instituto Nacional de Estatística), órgão oficial de dados demográficos e censitários.} Censo 2022, INE\footnote{Instituto Nacional de Estadística (Instituto Nacional de Estatística), órgão oficial de dados demográficos e censitários.} EPHC 2023, Ministerio
Público 2024\\
\textbf{Nota:} Valores marcados como \emph{(dept.)} referem-se ao
departamento de Canindeyú; valores \emph{(dist.)} são específicos deste
distrito. Dados marcados \emph{(est.)} são estimativas.

\begin{longtable}[]{@{}
  >{\raggedright\arraybackslash}p{(\linewidth - 4\tabcolsep) * \real{0.4231}}
  >{\raggedright\arraybackslash}p{(\linewidth - 4\tabcolsep) * \real{0.2692}}
  >{\raggedright\arraybackslash}p{(\linewidth - 4\tabcolsep) * \real{0.3077}}@{}}
\toprule\noalign{}
\begin{minipage}[b]{\linewidth}\raggedright
Indicador
\end{minipage} & \begin{minipage}[b]{\linewidth}\raggedright
Valor
\end{minipage} & \begin{minipage}[b]{\linewidth}\raggedright
Âmbito
\end{minipage} \\
\midrule\noalign{}
\endhead
\bottomrule\noalign{}
\endlastfoot
IDH (2020) & 0,685 & dept. \\
Ranking IDH nacional & 15/\allowbreak 18 & dept. \\
Esperança de vida & 71,2 anos & dept. \\
Escolaridade média & 7,3 anos & dept. \\
RNB per capita (USD PPA) & 7.100 & dept. \\
Pobreza monetária (\%) & 35,1\% & dept. \\
Pobreza extrema (\%) & 10,8\% & dept. \\
Índice de Gini & 0,467 & dept. \\
Acesso a água potável (\%) & 67,4\% & dept. \\
Acesso a saneamento (\%) & 64,8\% & dept. \\
Taxa de homicídios (est., /\allowbreak 100k hab) & \textasciitilde15--25
(estimativa, significativamente acima da média nacional) & dept. \\
Índice de segurança & baixo & dept. \\
População (Censo 2022) & 12.145 habitantes & dist. \\
Idade mediana & N/\allowbreak D & dist. \\
Taxa de fecundidade & N/\allowbreak D & dist. \\
\end{longtable}
\subsubsection{Combustível}

\textbf{Referência:} PETROPAR /\allowbreak  postos locais (2024) \textbf{Tipo de
localidade:} Interior

\begin{longtable}[]{@{}lll@{}}
\toprule\noalign{}
Combustível & USD/\allowbreak litro & Gs/\allowbreak litro (aprox.) \\
\midrule\noalign{}
\endhead
\bottomrule\noalign{}
\endlastfoot
Gasolina 93 oct & 1.00 & 7,400 \\
Gasolina 97 oct (premium) & 1.10 & 8,140 \\
Diesel & 0.93 & 6,882 \\
\end{longtable}

\begin{quote}
Preços podem variar ±5\% conforme posto e sazonalidade. Chaco e interior
remoto apresentam maior variação.
\end{quote}

\subsubsection{Cobertura Celular}

\textbf{Fonte:} CONATEL PY /\allowbreak  operadoras (2024)

\begin{longtable}[]{@{}ll@{}}
\toprule\noalign{}
Parâmetro & Valor \\
\midrule\noalign{}
\endhead
\bottomrule\noalign{}
\endlastfoot
Cobertura 4G população (dept.) & 78\% \\
Cobertura 4G área rural & 55\% \\
Melhor operadora & Tigo \\
Qualidade rural & moderada \\
\end{longtable}

\begin{quote}
Para áreas rurais fora do núcleo urbano, recomenda-se chip Tigo como
principal e Personal como backup.
\end{quote}

\subsubsection{Internet}

\textbf{Fonte:} CONATEL /\allowbreak  Speedtest Ookla (2024)

\begin{longtable}[]{@{}ll@{}}
\toprule\noalign{}
Parâmetro & Valor \\
\midrule\noalign{}
\endhead
\bottomrule\noalign{}
\endlastfoot
Velocidade média download & 35 Mbps \\
Domicílios com internet (dept.) & 45\% \\
Tecnologia predominante & rádio \\
Opção rural & Starlink disponível (\textasciitilde USD 44/\allowbreak mês) \\
\end{longtable}

\subsubsection{Mercado Imobiliário e Terra
Rural}

\textbf{Fonte:} INDERT /\allowbreak  Clasificados.com.py (2024)

\begin{longtable}[]{@{}ll@{}}
\toprule\noalign{}
Tipo & Referência \\
\midrule\noalign{}
\endhead
\bottomrule\noalign{}
\endlastfoot
Terra agrícola alta prod. (USD/\allowbreak ha) & 7,500 \\
Imóvel urbano (USD/\allowbreak m²) & 650 \\
Aluguel 2 quartos (USD/\allowbreak mês) & 260 \\
\end{longtable}

\begin{quote}
Valores de referência departamental. Localidades menores podem ter
preços 20--40\% abaixo da capital departamental. \#\#\# Saúde
\end{quote}

\textbf{Fonte:} MSPBS /\allowbreak  IPS Paraguay (2024-2026), consolidação
departamental e proxy local

\begin{longtable}[]{@{}ll@{}}
\toprule\noalign{}
Serviço & Disponibilidade \\
\midrule\noalign{}
\endhead
\bottomrule\noalign{}
\endlastfoot
USF /\allowbreak  Posto de Saúde & sim \\
Hospital Regional & sim \\
IPS (seguro social) & não \\
Farmácia & sim \\
Distância ao hospital de referência & \textasciitilde104 km (Salto del
Guaira) \\
\end{longtable}

\textbf{Principais estabelecimentos:} USF local; referencia hospitalar
em Salto del Guaira

\textbf{Observação para imigrantes:} Atendimento primario local ou em
raio curto; casos de maior complexidade seguem para Salto del Guaira.
Cobertura privada continua recomendada para especialidades e urgências
de maior complexidade.
\newpage
\secaoDiagnostico{Ypejhu}{\mapaConteudo{-23.9000}{-55.4500}}{Ypejhú é um enclave de isolamento e produtividade no extremo norte de
Canindeyú. Oferece solos férteis e um isolamento demográfico que garante
privacidade absoluta, protegida pela geografia serrana. Sua força reside
na abundância de recursos naturais e na nova conectividade asfáltica.
Contudo, a proximidade com a fronteira seca e o contexto regional exigem
que o ocupante adote um perfil vigilante e autônomo. É o local ideal
para projetos de autossuficiência que valorizem a invisibilidade e a
segurança institucional de fronteira, transformando o ``final de linha''
geográfico em uma fortaleza de proteção tática.

\subsubsection{Por que sim}

\begin{itemize}
\tightlist
\item
  Máximo isolamento geográfico em território vasto (privacidade).
\item
  Abundância de águas puras e recursos proteicos (gado).
\item
  Conectividade asfáltica moderna via Ruta PY13.
\item
  Segurança institucional reforçada pela presença de forças de
  fronteira.
\end{itemize}

\subsubsection{Por que não}

\begin{itemize}
\tightlist
\item
  Insegurança tática regional latente (fronteira).
\item
  Restrições da Lei de Fronteira (2532/\allowbreak 05) para investidores
  estrangeiros.
\item
  Distância de centros urbanos com serviços de saúde de alta
  complexidade.
\end{itemize}

\subsubsection{Recomendação
Técnica}

Altamente indicado para estabelecimentos de autossuficiência que busquem
isolamento tático em solo produtivo. Requer investimento em segurança
patrimonial e autonomia em saúde. O GSS de 7.5 reflete a superioridade
dos recursos e o isolamento protegida pela vigilância institucional da
fronteira. Referencia atual de score: GSS 7.5 em 2026-03-05.

\subsubsection{}

\begin{itemize}
\tightlist
\item
  \begin{enumerate}
  \def\labelenumi{\arabic{enumi})}
  \tightlist
  \item
    Dados censitários validados (INE\footnote{Instituto Nacional de Estadística (Instituto Nacional de Estatística), órgão oficial de dados demográficos e censitários.} 2022).
  \end{enumerate}
\item
  \begin{enumerate}
  \def\labelenumi{\arabic{enumi})}
  \setcounter{enumi}{1}
  \tightlist
  \item
    Relatórios de segurança pública e inteligência de fronteira (MDI).
  \end{enumerate}
\item
  \begin{enumerate}
  \def\labelenumi{\arabic{enumi})}
  \setcounter{enumi}{2}
  \tightlist
  \item
    Mapa de solos e bacias hidrográficas de Canindeyú (MAG).
  \end{enumerate}
\item
  \begin{enumerate}
  \def\labelenumi{\arabic{enumi})}
  \setcounter{enumi}{3}
  \tightlist
  \item
    Registro de pavimentação da Ruta PY13 (MOPC).
  \end{enumerate}
\end{itemize}}
\begin{table}[ht!]
\small \begin{tabular}{p{3.0cm}p{0.8cm}p{6.0cm}} \toprule
\textbf{Critério} & \textbf{Nota} & \textbf{Justificativa} \\ \midrule
A - Ameaças Estratégicas & 7.0 & Localização remota e discreta; excelente invisibilidade tática fora da fronteira seca, equilibrada por riscos regionais \\
B - Risco Social & 8.0 & Densidade populacional extremamente baixa; isolamento social superior \\
C - Riscos Naturais & 7.0 & Geologia muito estável e relevo seguro contra inundações \\
D - Autossuficiência & 8.5 & Recursos agrícolas excelentes e abundância hídrica natural \\
E - Institucional & 7.0 & Forte segurança institucional de fronteira e presença estatal funcional \\
\midrule \textbf{GSS FINAL} & \textbf{7.5} & \textbf{Seguro (Altamente recomendado para quem busca isolamento tático e autossuficiência em solo fértil, ciente dos desafios de segurança regional).} \\ \bottomrule \end{tabular}
\end{table}
\subsubsection{Vulnerabilidades e
Mitigação}\label{vuln-Ypejhu-vulnerabilidades-e-mitigauxe7uxe3o}

\begin{itemize}
\tightlist
\item
  \textbf{Vulnerabilidade Tática de Fronteira:} Exposição a atividades
  ilícitas regionais (Mitigação: residência rural afastada 5-10 km da
  linha internacional e investimento em segurança patrimonial).
\item
  \textbf{Isolamento de Saúde:} Distância crítica de hospitais
  cirúrgicos (Mitigação: manutenção de estoque médico básico e kit
  trauma local).
\item
  \textbf{Dependência de Rodovia Única:} Risco de bloqueio logístico na
  PY13 (Mitigação: mapeamento de rotas rurais vicinais e autonomia de
  suprimentos para 60 dias).
\end{itemize}
\subsubsection{Dossiê de Campo}\label{dossie-Ypejhu-dossiuxea-de-campo}
\begin{description}[style=nextline,leftmargin=0.5cm,font=\bfseries]
\item[Coordenadas]  23°54'S 55°27'W.
\item[Topografia]  Relevo ondulado com presença de serras (Cordillera de Amambay) e vales férteis. Altitude \textasciitilde 320m. Área massiva de \textasciitilde 956 km². Geologicamente estável, risco sísmico desprezível.
\item[Alvos Estratégicos]  Ponto de fronteira seca internacional com Paranhos (Brasil); Áreas de pecuária extensiva e produção de grãos; Proximidade tática com o Parque Nacional Mbaracayú. Ausência de bases militares de grande escala, mas zona de alto controle policial e de inteligência de fronteira.
\item[Fallout]  Ventos predominantes do Norte (NE) e Leste. Baixíssimo risco de fallout direto de alvos continentais primários devido ao isolamento geográfico.
\item[População]  8.607 habitantes (Censo 2022 Final). População com forte componente demográfico rural e indígena, integrada ao agronegócio e comércio fronteiriço.
\item[Densidade]  \textasciitilde 9,0 hab/\allowbreak km² (Densidade extremamente baixa, oferecendo isolamento tático e privacidade fora do núcleo urbano fronteiriço).
\item[Vias de Acesso]  Conectividade asfáltica via Ruta PY13 (concluída recentemente). Localização em "final de linha" geográfico que favorece o anonimato estratégico, mas sujeita a controle rigoroso da fronteira seca.
\item[Serviços]  Infraestrutura básica funcional. Centro de Saúde local e USF. Dependência de Curuguaty (\textasciitilde 90 km) para atendimentos de alta complexidade e logística comercial especializada.
\item[Custo de Vida]  Muito baixo; economia baseada na pecuária, grãos e comércio de fronteira.
\item[Preço da Terra]  US\$ 5.000-8.000/\allowbreak ha (terras mecanizadas); US\$ 2.000-4.000/\allowbreak ha (áreas de pecuária/\allowbreak campo natural).
\item[Hidrologia]  Baixo risco de inundações sistêmicas devido à topografia ondulada favorável à drenagem. Presença de arroios perenes de alta qualidade hídrica tributários da bacia do Rio Paraná.
\item[Clima]  Subtropical Úmido. Verões quentes e chuvosos; invernos amenos com geadas ocasionais.
\item[Energia]  Rede nacional (Itaipu); instável em áreas rurais remotas.
\item[Água]  Abundante via nascentes naturais e poços artesianos profundos; excelente qualidade hídrica subterrânea.
\item[Qualidade do Solo]  Solo franco-argiloso fértil; aptidão para soja, milho, pastagens e erva-mate.
\item[Resources Locais]  Grande produção de gado bovino e grãos. Elevadíssimo potencial para autossuficiência alimentar básica em nível de propriedade devido à vasta área produtiva e baixa população.
\item[Segurança]  Zona de fronteira com vigilância tática permanente. Criminalidade urbana comum moderada, mas sensível a tensões do crime organizado regional típico de Canindeyú. Coesão social baseada na cultura produtora tradicional. Presença institucional de forças de segurança de fronteira.
\item[Leis Local: Município consolidado. Restrição de Fronteira (Lei 2532/\allowbreak 05)]  Aplicável a terras rurais na faixa de 50km da fronteira brasileira.
\end{description}
\paragraph{Solo (SoilGrids 2.0, média ponderada 0--30
cm)}

\begin{longtable}[]{@{}ll@{}}
\toprule\noalign{}
Parâmetro & Valor \\
\midrule\noalign{}
\endhead
\bottomrule\noalign{}
\endlastfoot
pH (H$_2$O) & 5.32 \\
Carbono orgânico (SOC) & 20.89 g/\allowbreak kg \\
Argila & 20.98 \% \\
Areia & 67.28 \% \\
Silte (calc.) & 11.7 \% \\
Densidade aparente & 1.26 g/\allowbreak cm³ \\
\textbf{Aptidão agrícola} & \textbf{Média} \\
\end{longtable}

Fonte: ISRIC SoilGrids 2.0 via WCS (acesso em 2026-03-21). Coords:
-23.9°, -55.45°. Média ponderada camadas 0-5, 5-15, 15-30 cm.

\begin{itemize}
\tightlist
\item
  \textbf{Homicidios/\allowbreak 100k:} \textasciitilde18,0 (Regional/\allowbreak Fronteira).
\end{itemize}

\subsection{10. Analise de Riscos}

\begin{itemize}
\tightlist
\item
  \textbf{Cenario curto prazo:} Manutenção da tranquilidade rural e
  controle de tráfego na fronteira.
\item
  \textbf{Cenario medio prazo:} Valorização das terras pela consolidação
  da infraestrutura asfáltica e expansão da produção.
\end{itemize}

\subsection{11. Redacao Final}

\subsubsection{Diagnostico Integrado}

Ypejhú é um enclave de isolamento e produtividade no extremo norte de
Canindeyú. Oferece solos férteis e um isolamento demográfico que garante
privacidade absoluta, protegida pela geografia serrana. Sua força reside
na abundância de recursos naturais e na nova conectividade asfáltica.
Contudo, a proximidade com a fronteira seca e o contexto regional exigem
que o ocupante adote um perfil vigilante e autônomo. É o local ideal
para projetos de autossuficiência que valorizem a invisibilidade e a
segurança institucional de fronteira, transformando o ``final de linha''
geográfico em uma fortaleza de proteção tática.

\subsubsection{Por que sim}

\begin{itemize}
\tightlist
\item
  Máximo isolamento geográfico em território vasto (privacidade).
\item
  Abundância de águas puras e recursos proteicos (gado).
\item
  Conectividade asfáltica moderna via Ruta PY13.
\item
  Segurança institucional reforçada pela presença de forças de
  fronteira.
\end{itemize}

\subsubsection{Por que nao}

\begin{itemize}
\tightlist
\item
  Insegurança tática regional latente (fronteira).
\item
  Restrições da Lei de Fronteira (2532/\allowbreak 05) para investidores
  estrangeiros.
\item
  Distância de centros urbanos com serviços de saúde de alta
  complexidade.
\end{itemize}

\subsubsection{Recomendacao Tecnica}

Altamente indicado para estabelecimentos de autossuficiência que busquem
isolamento tático em solo produtivo. Requer investimento em segurança
patrimonial e autonomia em saúde. O GSS de 7.5 reflete a superioridade
dos recursos e o isolamento protegida pela vigilância institucional da
fronteira. Referencia atual de score: GSS 7.5 em 2026-03-05.

\subsection{13.}

\begin{itemize}
\tightlist
\item
  \begin{enumerate}
  \def\labelenumi{\arabic{enumi})}
  \tightlist
  \item
    Dados censitários validados (INE\footnote{Instituto Nacional de Estadística (Instituto Nacional de Estatística), órgão oficial de dados demográficos e censitários.} 2022).
  \end{enumerate}
\item
  \begin{enumerate}
  \def\labelenumi{\arabic{enumi})}
  \setcounter{enumi}{1}
  \tightlist
  \item
    Relatórios de segurança pública e inteligência de fronteira (MDI).
  \end{enumerate}
\item
  \begin{enumerate}
  \def\labelenumi{\arabic{enumi})}
  \setcounter{enumi}{2}
  \tightlist
  \item
    Mapa de solos e bacias hidrográficas de Canindeyú (MAG).
  \end{enumerate}
\item
  \begin{enumerate}
  \def\labelenumi{\arabic{enumi})}
  \setcounter{enumi}{3}
  \tightlist
  \item
    Registro de pavimentação da Ruta PY13 (MOPC).
  \end{enumerate}
\end{itemize}

\subsubsection{Combustível}

\textbf{Referência:} PETROPAR /\allowbreak  postos locais (2024) \textbf{Tipo de
localidade:} Interior

\begin{longtable}[]{@{}lll@{}}
\toprule\noalign{}
Combustível & USD/\allowbreak litro & Gs/\allowbreak litro (aprox.) \\
\midrule\noalign{}
\endhead
\bottomrule\noalign{}
\endlastfoot
Gasolina 93 oct & 1.00 & 7,400 \\
Gasolina 97 oct (premium) & 1.10 & 8,140 \\
Diesel & 0.93 & 6,882 \\
\end{longtable}

\begin{quote}
Preços podem variar ±5\% conforme posto e sazonalidade. Chaco e interior
remoto apresentam maior variação.
\end{quote}

\subsubsection{Cobertura Celular}

\textbf{Fonte:} CONATEL PY /\allowbreak  operadoras (2024)

\begin{longtable}[]{@{}ll@{}}
\toprule\noalign{}
Parâmetro & Valor \\
\midrule\noalign{}
\endhead
\bottomrule\noalign{}
\endlastfoot
Cobertura 4G população (dept.) & 78\% \\
Cobertura 4G área rural & 55\% \\
Melhor operadora & Tigo \\
Qualidade rural & moderada \\
\end{longtable}

\begin{quote}
Para áreas rurais fora do núcleo urbano, recomenda-se chip Tigo como
principal e Personal como backup.
\end{quote}

\subsubsection{Internet}

\textbf{Fonte:} CONATEL /\allowbreak  Speedtest Ookla (2024)

\begin{longtable}[]{@{}ll@{}}
\toprule\noalign{}
Parâmetro & Valor \\
\midrule\noalign{}
\endhead
\bottomrule\noalign{}
\endlastfoot
Velocidade média download & 35 Mbps \\
Domicílios com internet (dept.) & 45\% \\
Tecnologia predominante & rádio \\
Opção rural & Starlink disponível (\textasciitilde USD 44/\allowbreak mês) \\
\end{longtable}

\subsubsection{Mercado Imobiliário e Terra
Rural}

\textbf{Fonte:} INDERT /\allowbreak  Clasificados.com.py (2024)

\begin{longtable}[]{@{}ll@{}}
\toprule\noalign{}
Tipo & Referência \\
\midrule\noalign{}
\endhead
\bottomrule\noalign{}
\endlastfoot
Terra agrícola alta prod. (USD/\allowbreak ha) & 7,500 \\
Imóvel urbano (USD/\allowbreak m²) & 650 \\
Aluguel 2 quartos (USD/\allowbreak mês) & 260 \\
\end{longtable}

\begin{quote}
Valores de referência departamental. Localidades menores podem ter
preços 20--40\% abaixo da capital departamental. \#\#\# Saúde
\end{quote}

\textbf{Fonte:} MSPBS /\allowbreak  IPS Paraguay (2024-2026), consolidação
departamental e proxy local

\begin{longtable}[]{@{}ll@{}}
\toprule\noalign{}
Serviço & Disponibilidade \\
\midrule\noalign{}
\endhead
\bottomrule\noalign{}
\endlastfoot
USF /\allowbreak  Posto de Saúde & sim \\
Hospital Regional & não \\
IPS (seguro social) & não \\
Farmácia & sim \\
Distância ao hospital de referência & \textasciitilde143 km (Salto del
Guaira) \\
\end{longtable}

\textbf{Principais estabelecimentos:} USF local; referencia hospitalar
em Salto del Guaira

\textbf{Observação para imigrantes:} Atendimento primario local ou em
raio curto; casos de maior complexidade seguem para Salto del Guaira.
Cobertura privada continua recomendada para especialidades e urgências
de maior complexidade.
\subsubsection{Indicadores Sociais}

\textbf{Fontes:} PNUD 2020, INE\footnote{Instituto Nacional de Estadística (Instituto Nacional de Estatística), órgão oficial de dados demográficos e censitários.} Censo 2022, INE\footnote{Instituto Nacional de Estadística (Instituto Nacional de Estatística), órgão oficial de dados demográficos e censitários.} EPHC 2023, Ministerio
Público 2024\\
\textbf{Nota:} Valores marcados como \emph{(dept.)} referem-se ao
departamento de Canindeyú; valores \emph{(dist.)} são específicos deste
distrito. Dados marcados \emph{(est.)} são estimativas.

\begin{longtable}[]{@{}
  >{\raggedright\arraybackslash}p{(\linewidth - 4\tabcolsep) * \real{0.4231}}
  >{\raggedright\arraybackslash}p{(\linewidth - 4\tabcolsep) * \real{0.2692}}
  >{\raggedright\arraybackslash}p{(\linewidth - 4\tabcolsep) * \real{0.3077}}@{}}
\toprule\noalign{}
\begin{minipage}[b]{\linewidth}\raggedright
Indicador
\end{minipage} & \begin{minipage}[b]{\linewidth}\raggedright
Valor
\end{minipage} & \begin{minipage}[b]{\linewidth}\raggedright
Âmbito
\end{minipage} \\
\midrule\noalign{}
\endhead
\bottomrule\noalign{}
\endlastfoot
IDH (2020) & 0,685 & dept. \\
Ranking IDH nacional & 15/\allowbreak 18 & dept. \\
Esperança de vida & 71,2 anos & dept. \\
Escolaridade média & 7,3 anos & dept. \\
RNB per capita (USD PPA) & 7.100 & dept. \\
Pobreza monetária (\%) & 35,1\% & dept. \\
Pobreza extrema (\%) & 10,8\% & dept. \\
Índice de Gini & 0,467 & dept. \\
Acesso a água potável (\%) & 67,4\% & dept. \\
Acesso a saneamento (\%) & 64,8\% & dept. \\
Taxa de homicídios (est., /\allowbreak 100k hab) & \textasciitilde15--25
(estimativa, significativamente acima da média nacional) & dept. \\
Índice de segurança & baixo & dept. \\
População (Censo 2022) & 8.607 habitantes & dist. \\
Idade mediana & N/\allowbreak D & dist. \\
Taxa de fecundidade & N/\allowbreak D & dist. \\
\end{longtable}
\subsubsection{Combustível}

\textbf{Referência:} PETROPAR /\allowbreak  postos locais (2024) \textbf{Tipo de
localidade:} Interior

\begin{longtable}[]{@{}lll@{}}
\toprule\noalign{}
Combustível & USD/\allowbreak litro & Gs/\allowbreak litro (aprox.) \\
\midrule\noalign{}
\endhead
\bottomrule\noalign{}
\endlastfoot
Gasolina 93 oct & 1.00 & 7,400 \\
Gasolina 97 oct (premium) & 1.10 & 8,140 \\
Diesel & 0.93 & 6,882 \\
\end{longtable}

\begin{quote}
Preços podem variar ±5\% conforme posto e sazonalidade. Chaco e interior
remoto apresentam maior variação.
\end{quote}

\subsubsection{Cobertura Celular}

\textbf{Fonte:} CONATEL PY /\allowbreak  operadoras (2024)

\begin{longtable}[]{@{}ll@{}}
\toprule\noalign{}
Parâmetro & Valor \\
\midrule\noalign{}
\endhead
\bottomrule\noalign{}
\endlastfoot
Cobertura 4G população (dept.) & 78\% \\
Cobertura 4G área rural & 55\% \\
Melhor operadora & Tigo \\
Qualidade rural & moderada \\
\end{longtable}

\begin{quote}
Para áreas rurais fora do núcleo urbano, recomenda-se chip Tigo como
principal e Personal como backup.
\end{quote}

\subsubsection{Internet}

\textbf{Fonte:} CONATEL /\allowbreak  Speedtest Ookla (2024)

\begin{longtable}[]{@{}ll@{}}
\toprule\noalign{}
Parâmetro & Valor \\
\midrule\noalign{}
\endhead
\bottomrule\noalign{}
\endlastfoot
Velocidade média download & 35 Mbps \\
Domicílios com internet (dept.) & 45\% \\
Tecnologia predominante & rádio \\
Opção rural & Starlink disponível (\textasciitilde USD 44/\allowbreak mês) \\
\end{longtable}

\subsubsection{Mercado Imobiliário e Terra
Rural}

\textbf{Fonte:} INDERT /\allowbreak  Clasificados.com.py (2024)

\begin{longtable}[]{@{}ll@{}}
\toprule\noalign{}
Tipo & Referência \\
\midrule\noalign{}
\endhead
\bottomrule\noalign{}
\endlastfoot
Terra agrícola alta prod. (USD/\allowbreak ha) & 7,500 \\
Imóvel urbano (USD/\allowbreak m²) & 650 \\
Aluguel 2 quartos (USD/\allowbreak mês) & 260 \\
\end{longtable}

\begin{quote}
Valores de referência departamental. Localidades menores podem ter
preços 20--40\% abaixo da capital departamental. \#\#\# Saúde
\end{quote}

\textbf{Fonte:} MSPBS /\allowbreak  IPS Paraguay (2024-2026), consolidação
departamental e proxy local

\begin{longtable}[]{@{}ll@{}}
\toprule\noalign{}
Serviço & Disponibilidade \\
\midrule\noalign{}
\endhead
\bottomrule\noalign{}
\endlastfoot
USF /\allowbreak  Posto de Saúde & sim \\
Hospital Regional & não \\
IPS (seguro social) & não \\
Farmácia & sim \\
Distância ao hospital de referência & \textasciitilde143 km (Salto del
Guaira) \\
\end{longtable}

\textbf{Principais estabelecimentos:} USF local; referencia hospitalar
em Salto del Guaira

\textbf{Observação para imigrantes:} Atendimento primario local ou em
raio curto; casos de maior complexidade seguem para Salto del Guaira.
Cobertura privada continua recomendada para especialidades e urgências
de maior complexidade.
